\chapter{Lean Export of Tensor Networks}\label{ch:lean-export-net}

% Listing style for generated Lean: the shared literate list plus the characters the export emits.
\lstdefinestyle{leannet}{
  style=lean,
  literate={→}{{$\to$}}1 {←}{{$\leftarrow$}}1 {≤}{{$\le$}}1 {≥}{{$\ge$}}1 {≠}{{$\ne$}}1
           {ω}{{$\omega$}}1 {⟩}{{$\rangle$}}1 {⟨}{{$\langle$}}1 {∀}{{$\forall$}}1
           {∃}{{$\exists$}}1 {λ}{{$\lambda$}}1 {∑}{{$\sum$}}1 {·}{{$\cdot$}}1
           {ℂ}{{$\mathbb{C}$}}1 {ℕ}{{$\mathbb{N}$}}1 {ℤ}{{$\mathbb{Z}$}}1 {ℚ}{{$\mathbb{Q}$}}1
           {↦}{{$\mapsto$}}1 {⊗}{{$\otimes$}}1 {∘}{{$\circ$}}1 {≫}{{$\gg$}}1
           {α}{{$\alpha$}}1 {β}{{$\beta$}}1 {π}{{$\pi$}}1 {∈}{{$\in$}}1 {×}{{$\times$}}1
           {₀}{{$_0$}}1 {₁}{{$_1$}}1 {₂}{{$_2$}}1 {⁻}{{$^{-}$}}1 {¹}{{$^{1}$}}1
           {√}{{$\surd$}}1 {∧}{{$\wedge$}}1 {∨}{{$\vee$}}1 {¬}{{$\neg$}}1 {↔}{{$\leftrightarrow$}}1
           {⟶}{{$\longrightarrow$}}1 {ε}{{$\varepsilon$}}1 {μ}{{$\mu$}}1 {σ}{{$\sigma$}}1
           {Δ}{{$\Delta$}}1 {Φ}{{$\Phi$}}1 {Ψ}{{$\Psi$}}1 {θ}{{$\theta$}}1 {ℝ}{{$\mathbb{R}$}}1
           {▸}{{$\triangleright$}}1 {↓}{{$\downarrow$}}1 {⊢}{{$\vdash$}}1 {∏}{{$\prod$}}1
           {₃}{{$_3$}}1 {φ}{{$\varphi$}}1 {τ}{{$\tau$}}1 {∉}{{$\notin$}}1 {!}{{!}}1,
}
\newcommand{\leanx}{\lstinline[style=leannet,frame=none,basicstyle=\ttfamily,breaklines=true,breakatwhitespace=false]}

The module \module{archytaszx.semantics.lean_net} (file \file{archytaszx/semantics/lean_net.py})
translates a theorem of the proof kernel (\cref{ch:kernel}) into a Lean~4 source file
(\tool{Lean~4}{https://lean-lang.org}~\cite{demoura2021lean4}, with
\tool{Mathlib}{https://github.com/leanprover-community/mathlib4}~\cite{mathlib2020}) that states
and proves the same equation over the \term{tensor-network semantics} of the bundled Lean library
\code!ArchytasZX! (\cref{ch:lean-core,ch:lean-boxes}). Every diagram of the proof becomes a Lean
\lean!Net!: one \lean!Node! per diagram node, every port carrying a natural-number label, a label
shared by two ports standing for one contracted index in the sense of Penrose's tensor
networks~\cite{penrose1971negative}. Every inference of the kernel theorem becomes a Lean proof
term: an isomorphism is an instance of \lean!isoB!, a rule step is an instance of the rule's local
lemma placed in context by \lean!rewrite!, and the kernel's symmetry and transitivity are
\lean!Equal.symm! and \lean!Equal.mul_trans!. A rule step that has no lemma, or whose shape the
exporter cannot match, becomes a Lean \lean!axiom! named \code!unproved_*!; the generated file
names every such axiom in its header, and Lean's \lean!#print axioms! lists it. The engine is thus
Lean-certifying: a file without \code!unproved_! axioms that Lean accepts is a Lean-certified
proof of the kernel theorem over the semantics of the ZX, ZW and qufinite ZXW
calculi~\cite{coecke2011interacting,hadzihasanovic2015zw,poor2023zxw,wang2023qufinitezxw}.

The module also restates the final theorem on the engine's own diagrams, written as Lean
\lean!Diagram! records, through the Lean-proved translation \lean!DEqual_iff!; and, for symbolic
theorems over \bboxes{}~\cite{kissinger2014pattern,carette2022bangboxes}, it emits the network of
the \bbox{} \term{family} at every multiplicity, proved by induction on the multiplicity in Lean.

This chapter specifies the module completely: the entry points, the data structures, the tables
the export accumulates, the translation of a diagram into a network, every emitted declaration
and its exact text, the proof of every kind of kernel inference, the catalogue of rule
translators, the families and inductions, the restatement on diagrams, the size controls, the
Lean invocation, and every failure. The Lean library is specified in
\cref{ch:lean-core,ch:lean-boxes,ch:lean-pathsum,ch:lean-sym}; this chapter states of it only
what the emitted text refers to. The opaque and symbolic exports are specified in
\cref{ch:lean-export}.

\section{Position in the architecture}\label{sec:lean-export-net-position}

The kernel's \py!lean! function (\cref{ch:kernel}) calls \py!lean_net.export! first for every
plain \py!Theorem! without a domain restriction, when the semantic export is wanted and no guard
hypotheses are given. When the returned source contains no line beginning
\code!axiom unproved_!, the kernel returns it. Otherwise the kernel tries the symbolic export
(\module{archytaszx.semantics.lean_sym}), and falls back on the network source with its axioms
when that fails too. The restricted export of the kernel (theorems with a domain) calls
\py!lean_net.export! with the \py!original! argument (\cref{sec:lean-export-net-original}). Other
modules import internal helpers of this module:
\module{archytaszx.semantics.lean_sym} uses \py!_Export!, \py!_join!, \py!_paren!, \py!LNet!,
\py!_tail!, \py!_split_net!, \py!_zero_net! and \py!_norm!;
\module{archytaszx.semantics.lean_numeric} uses \py!_head!, \py!_join!, \py!_paren!;
\module{archytaszx.semantics.lean_refute} uses \py!_port_dim! and other names;
\module{archytaszx.semantics.lean_famlink} uses \py!LNet!, \py!LNode!, \py!_head!, \py!_join!,
\py!_list!, \py!_pairs!, \py!_paren!, \py!_perm_to_bijection!, \py!_rename_net!,
\py!_split_net!, \py!_tail!, \py!_zero_net! and \py!_zk!;
\module{archytaszx.semantics.lean_subst} uses \py!LNet! and \py!_head!; and
\module{archytaszx.semantics.prove}, \module{archytaszx.semantics.decide},
\module{archytaszx.semantics.lean_copyphase}, \module{archytaszx.semantics.lean_chainfam},
\module{archytaszx.semantics.lean_pairs} and
\module{archytaszx.semantics.certify_decision} catch \py!NetUnsupported!. In turn, \py!export!
calls \module{archytaszx.semantics.lean_chainnode} and
\module{archytaszx.semantics.lean_famlink} when its own run fails
(\cref{alg:lean-export-net-export}). The exporters of
\cref{sec:lean-export-net-beside} prove family claims without a kernel theorem and are called by
\module{archytaszx.semantics.certify_decision} and \module{archytaszx.semantics.kernel}.

The module depends on \module{archytaszx.semantics.lean_scalar} for writing exact scalars in Lean
and for proving identities between them (the class \py!Scalars!, the sympy functions \py!Rt!,
\py!Om!, \py!Inv!, \py!NC!, the symbol \py!II!, the class \py!Field! and the exception
\py!ScalarUnsupported!), on \module{archytaszx.algebra.numfield} for number-field certificates,
on \module{archytaszx.diagram.compare} for diagram isomorphism, on the rewrite engine for rule
steps, and on \module{archytaszx.semantics.lean_qdigits} for one special case of
\py!export!.

\section{Public interface}\label{sec:lean-export-net-api}

The module exports exactly the four names of \py!__all__!: \py!NetUnsupported!, \py!check!,
\py!export! and \py!lean_project!.

\begin{signature}
\py!class NetUnsupported(Exception)!\\
A theorem or diagram outside the tensor-network export. Its message is a short English phrase
naming the obstacle (\cref{sec:lean-export-net-failures} lists every message).
\end{signature}

\begin{signature}
\py!export(theorem: object, name: str = "main", describe: Callable[[Diagram], str] | None = None,!\\
\py!       original: tuple[dict[str, int], list[tuple[Diagram, Diagram]]] | None = None) -> str!
\end{signature}

\noindent\py!export! returns the complete text of one Lean file proving \py!theorem! over the
tensor-network semantics. The parameter \py!name! is the Lean name of the final theorem.
\py!describe! maps a diagram to the text of its doc comment; by default it returns
\code!a diagram of !$n$\code! nodes! with $n$ the node count. \py!original!, when given, is a
pair of the substituted dimension values (symbol name to natural number) and the list of pairs
(diagram used by the proof, original diagram before substitution); the final theorem is then
also restated over the original diagrams (\cref{sec:lean-export-net-original}). The function is
traced under the span name \code!lean.export!.

\begin{algorithm}[\py!export!]\label{alg:lean-export-net-export}
\begin{enumerate}
\item If \py!theorem! is not a kernel \py!Theorem!, raise \py!NetUnsupported! with the message
  \code!export takes a Theorem, got !\textit{TypeName}.
\item If the inference kind is \code!contraction!, \py!original! is \py!None!, and the theorem
  has no assumptions, no domain and no variables, return
  \py!export_qudit_digits_net(theorem.lhs, theorem.rhs, name)! of
  \module{archytaszx.semantics.lean_qdigits} (the qudit path-sum export over digit variables,
  \cref{ch:lean-pathsum}) unless it raises \py!PathSumError! or \py!ScalarUnsupported!, in which
  case continue.
\item Define \textsc{Run}(\textit{everything}, \textit{linked}): create a fresh \py!_Export! with
  the description function; set its \py!alg_everything! to \textit{everything}, its
  \py!original! to \py!original!, its \py!box_syms! to \py!_box_syms(theorem)!, and its
  \py!linked! to the set \textit{linked} (empty by default). Compute
  \py!_deferred(theorem)!, the dimension equations the rule steps assumed; call
  \py!assume(a, b)! for each. If there are none, call \py!assume(a, b)! for every pair of
  \py!_wire_constraints(theorem)! and set \py!wire_alias! to \py!True! for each. Return
  \py!finish(theorem, name)! and the export object, converting a \py!ScalarUnsupported! into
  \py!NetUnsupported! with the same message.
\item Run \textsc{Run}(\py!False!). When it raises \py!NetUnsupported!:
  \begin{enumerate}
  \item if the theorem has variables, \py!original! is \py!None! and the theorem has no
    assumptions, call \py!export_chain_node(theorem.lhs, theorem.rhs, name, kernel._var)!
    (\cref{sec:lean-export-net-chainnode}); a source returned is the result;
  \item with variables, compute \py!linked_indices(theorem)!
    (\cref{sec:lean-export-net-linked}), else the empty set; when it is empty, re-raise the
    original exception;
  \item otherwise the first run is \textsc{Run}(\py!False!, linked); its exceptions propagate.
  \end{enumerate}
\item If the first run's export has a number-field certificate and at least one unproved step,
  run \textsc{Run}(\py!True!, \textit{first}\py!.linked!); if that raises \py!NetUnsupported!,
  return the first source; if it has strictly fewer unproved steps, return its source.
\item Return the first source.
\end{enumerate}
\end{algorithm}

\begin{signature}
\py!lean_project() -> str!\\
The directory of the bundled Lean project: the directory \file{lean} beside the package's
\file{semantics} directory, that is \file{archytaszx/lean}.
\end{signature}

\begin{signature}
\py!check(source: str, path: str, timeout: int = 1800) -> tuple[bool, str]!
\end{signature}

\noindent\py!check! writes \py!source! to \py!path! (UTF-8) and runs Lean on it in the bundled
project. It returns \py!(accepted, output)!.

\begin{algorithm}[\py!check!]\label{alg:lean-export-net-check}
\begin{enumerate}
\item Find \code!lake!: the first \code!lake! on the search path, else
  \file{~/.elan/bin/lake} if that file exists. If neither exists, raise \py!NetUnsupported!
  (\code!Lean (lake) is not installed!) before writing anything.
\item Write \py!source! to \py!path!.
\item If the directory \file{.lake/packages/mathlib} does not exist in the project, run
  \code!lake exe cache get! in the project; on a nonzero exit return \py!(False, stdout + stderr)!.
\item Run \code!lake build! in the project; on a nonzero exit return
  \py!(False, stdout + stderr)!.
\item Run \code!lake env lean !\textit{abspath(path)} in the project and return
  \py!(returncode == 0, stdout + stderr)!.
\end{enumerate}
Each subprocess has the time limit \py!timeout! seconds (default 1800); an expired limit raises
the standard \py!subprocess.TimeoutExpired!.
\end{algorithm}

\subsection{Module constants}\label{sec:lean-export-net-constants}

\begin{xltabular}{\linewidth}{@{}p{3.2cm}X@{}}
\toprule
Name & Value and use \\
\midrule
\texttt{ONE} & \texttt{sympy.Integer(1)}, the unit scalar model. \\
\texttt{H\_\allowbreak{}OFFSET} & $2^{16} = 65536$; the label offset of the peeled copy in a split along a
  multiplicity (\cref{sec:lean-export-net-split}). Deeper splits use $2^{16}\cdot 2^{t}$. \\
\texttt{RAW\_\allowbreak{}OFFSET} & $10^{6}$; the first dimension index of a dimension that a wire identifies with
  another (\cref{sec:lean-export-net-identified}). \\
\texttt{\_\allowbreak{}POSITIVE} & The Lean proof text
  \texttt{(by first \textbar{} omega \textbar{} norm\_\allowbreak{}num \textbar{} positivity \textbar{} (apply Nat.mul\_\allowbreak{}pos \textless{};\textgreater{} omega) \textbar{} exact Nat.pos\_\allowbreak{}pow\_\allowbreak{}of\_\allowbreak{}pos \_\allowbreak{} (by omega))},
  handed to \texttt{Scalars} as the proof that a natural expression is positive. \\
\texttt{\_\allowbreak{}KINDS} & The engine generator names \texttt{Z X F T Ti W B S} mapped to the Lean kind names of
  the same spelling; the pair-phase generator \texttt{Zp} takes the kind \texttt{ZP} through
  \texttt{zp\_\allowbreak{}kind} (\cref{sec:lean-export-net-zp}); any other generator (in particular the
  discard G) has no Lean kind. \\
\texttt{PHASE\_\allowbreak{}TABLE\_\allowbreak{}MIN} & 64; a phase vector with more rational entries and no rule is
  written as a packed table (\cref{alg:lean-export-net-phase-body}). \\
\texttt{TABLE\_\allowbreak{}CHUNK} & 256; a packed table with more entries is a list of named chunk definitions of
  this many entries (\py!table_get!). \\
\texttt{LIST\_\allowbreak{}SPLIT} & 4096; a node, summed-label, boundary-label or \texttt{labs} list with more
  entries is written as chunk definitions joined by \texttt{++} (\py!_split_list!). \\
\texttt{LIST\_\allowbreak{}CHUNK} & 2048; the entries of one list chunk definition. \\
\texttt{DOC\_\allowbreak{}MAX} & 65536; a doc comment is cut after this many characters and ends in a
  space and an ellipsis (\py!_comment!). \\
\texttt{\_\allowbreak{}PHASE\_\allowbreak{}TABLES\_\allowbreak{}MAX} & 4096; the most phase vectors the module-level cache
  \texttt{\_\allowbreak{}PHASE\_\allowbreak{}TABLES} holds; a full cache is cleared before the next entry. \\
\texttt{\_\allowbreak{}RESERVED} & The identifiers a parameter may not take: \texttt{E $\Phi$ $\Psi$ $\Delta$ labs k at by do fun in let if then else end from have show with match theorem def axiom open ex om fou Net Node Env Equal I $\pi$ pi exp sqrt F FC SpC ksOf bundL sS sG $\Delta$pos j}. \\
\texttt{\_\allowbreak{}FAMILY\_\allowbreak{}UNFOLD} & \texttt{F, Fam.env, bundL, List.lookup, ksOf, List.getD}: the definitions a
  tactic unfolds to read the environment of a family. \\
\texttt{\_\allowbreak{}THETA\_\allowbreak{}HYPS} & \texttt{(h$\theta$ : pev $\theta$p $\theta$ = 0) (h$\theta$lo : $\theta$lo \textless{} $\theta$) (h$\theta$hi : $\theta$ \textless{} $\theta$hi)}, the hypotheses
  fixing the real root $\theta$ of a number-field certificate. \\
\texttt{\_\allowbreak{}ZC} & \texttt{("Z", "ZW")}, the Lean kinds whose coefficient is read by \texttt{zc}. \\
\texttt{\_\allowbreak{}CHK} & \texttt{(by chk)}. \\
\texttt{\_\allowbreak{}FORWARD} & The converse rules whose forward rule has another name
  (\cref{sec:lean-export-net-forward}). \\
\texttt{\_\allowbreak{}RULES} & The rule translators (\cref{sec:lean-export-net-rules}). \\
\bottomrule
\end{xltabular}

\section{The semantics the export targets}\label{sec:lean-export-net-semantics}

The emitted code refers to the following definitions of the Lean library (namespace
\lean!ArchytasZX!); \cref{ch:lean-core,ch:lean-boxes} specify them. All proofs live in the
namespace \lean!ArchytasZX.Export!.

\begin{definition}[Kinds, nodes, networks]\label{def:lean-export-net-net}
The inductive type \lean!Kind! has the constructors \lean!Z!, \lean!X!, \lean!F!, \lean!T!,
\lean!Ti!, \lean!W!, \lean!B!, \lean!S!, \lean!I! (the bare wire), \lean!ZW! (the weighted Z
spider), \lean!ZB cs!, \lean!XB co ci!, \lean!WB cs! (a Z, X or W node whose last legs bundle
along the chains \lean!cs!, \lean!co!, \lean!ci! of multiplicity indices, outermost first),
\lean!Bun m b! (the copies of the generator \lean!b! along multiplicity \lean!m!), \lean!XW!
(the weighted X spider) and \lean!ZP dx dy c! (the pair phase over two legs of dimension indices
\lean!dx!, \lean!dy!, with integer coefficient \lean!c!). A \lean!Node! is the record \lean!⟨kind, outs, ins, dim, ph⟩! of a kind,
the output labels, the input labels, a dimension index and a phase index. A \lean!Net! is the
record \lean!⟨nodes, int, bnd⟩! of a node list, the summed (internal) labels and the boundary
labels.
\end{definition}

\begin{definition}[Environment and value]\label{def:lean-export-net-env}
An \lean!Env! has fields \lean!dimv : ℕ → ℕ! (the value of each dimension index),
\leanx!phase : ℕ → ℕ → ℝ! (each phase index's phase vector in turns),
\lean!labs : List (ℕ × ℕ)! (each label's dimension index), \lean!mult : ℕ → ℕ! (default
\lean!fun _ => 0!, each multiplicity's value) and \lean!coef : ℕ → ℕ → ℂ! (default
\lean!fun _ _ => 1!, each weight index's coefficient vector). With $e(\theta) = e^{2\pi i\theta}$
(Lean \lean!ex!), $\omega_d = e^{2\pi i/d}$ (Lean \lean!om!) and the Fourier entry
$\mathrm{fou}(d,j,k) = \omega_d^{jk}/\sqrt d$, the node tensors are
\begin{align*}
  t_Z(d,\varphi)(v_1,\dots,v_n) &= \textstyle\sum_{k<d} [v_1=\dots=v_n=k]\, e(\varphi_k),\\
  t_X(d,\varphi)(o;i) &= \textstyle\sum_{k<d} e(\varphi_k) \prod_o \mathrm{fou}(d,o,k)
     \prod_i \overline{\mathrm{fou}(d,i,k)},\\
  t_T(o,i) &= [o=0 \lor o=i], \qquad
  t_{T^{-1}}(o,i) = [o=i] - [o=0 \land i\ge 1 \land o \ne i],\\
  t_W(d)(o_1,\dots,o_n;i) &= \begin{cases}[\forall j,\ o_j=0] & i=0\\
     [i<d \land (o_j)_{o_j\ne0} = (i)] & i\ne0\end{cases},\\
  t_B(s,t)(k;a,b) &= [a<s \land b<t \land k = a t + b], \qquad t_I(d)(a,b) = [a=b<d],\\
  t_{ZP}(M, c)(x, y) &= e(c\,x\,y/M),
\end{align*}
the legs of \lean!ZP! read in the order outputs then inputs and $M$ the value of the node's
dimension index (\lean!tP!; a leg list of another length gives 0),
the splitter \lean!S! having the binder's entry with outputs and inputs exchanged, and the
weighted spiders \lean!ZW!, \lean!XW! replacing $e(\varphi_k)$ by \lean!coef ph k!. The value of
a network $N$ at a boundary assignment $\sigma$ is
\[
  \mathrm{eval}_E(N)(\sigma) = \sum_{\tau \text{ on } N.\mathrm{int}}\ \prod_{n\in N.\mathrm{nodes}}
  \mathrm{val}_E(\sigma\cup\tau)(n),
\]
each summed label ranging over its dimension \lean!labDim E l = E.dimv ((E.labs.lookup l).getD 0)!.
\end{definition}

\begin{definition}[Equality of scaled networks]\label{def:lean-export-net-equal}
\lean!Equal E s₁ N₁ s₂ N₂! holds when \lean!N₁.bnd! is a permutation of \lean!N₂.bnd! and
$s_1\,\mathrm{eval}_E(N_1)(\sigma) = s_2\,\mathrm{eval}_E(N_2)(\sigma)$ for every $\sigma$ in range
on \lean!N₁.bnd!. \lean!LocalEq E L iL R iR iface kL kR! states
$k_L \sum_{iL}\prod L = k_R\sum_{iR}\prod R$ at every assignment in range on \lean!iface!.
\end{definition}

The emitted proofs use these library facts, each specified in \cref{ch:lean-core,ch:lean-boxes}:
\lean!Equal.refl E s N!; \lean!Equal.symm!; \lean!Equal.trans!;
\lean!Equal.mul_trans : Equal E a X b Y → Equal E c Y e Z → Equal E (a*c) X (b*e) Z!;
\lean!Equal.rescale (h : Equal E a X b Y) (ha : a ≠ 0) (hst : s*b = t*a) : Equal E s X t Y!;
\lean!rewrite E C L R iC iL iR bnd iface s k s' hloc hL hR hface hs!, which gives
\lean!Equal E s ⟨C ++ L, iC ++ iL, bnd⟩ s' ⟨C ++ R, iC ++ iR, bnd⟩! from
\lean!LocalEq E L iL R iR iface 1 k!, two disjointness conditions, the interface condition and
\lean!s' = s * k!; \lean!plug!, the same with arbitrary scalars $k_L$, $k_R$ and no
\lean!hs!; \lean!wire_elim E s ns₁ ns₂ int bnd a b dz!, which removes the bare wire
\lean!⟨.I, [a], [b], dz, 0⟩! from between \lean!ns₁! and \lean!ns₂!, renames $b$ to $a$ and erases
$b$ from the summed labels; \lean!isoB E s A B π p h!, which gives \lean!Equal E s A s B! from the
Boolean check \lean!isoOKB E.labs A B π p = true! (relabelling by the pairs $\pi$, node
permutation $p$, node equivalence \lean!nodeEquivB!); \lean!Equal.toLocal!; \lean!Equal.rename!;
\lean!netCoef!, \lean!valCoefZ!, \lean!valCoefX!; and the family facts \lean!Equal.of_split_at!,
\lean!Equal.of_zero_at!, \lean!Spec.splitOK!, \lean!Net.split!, \lean!Net.zero!,
\lean!Net.rename!.

\section{Data structures}\label{sec:lean-export-net-data}

\subsection{Kinds on the Python side}\label{sec:lean-export-net-kinds}

A kind is a string or a tuple. The strings are the plain kinds \code!"Z" "X" "F" "T" "Ti" "W" "B" "S" "I" "ZW" "XW"!.
The tuples are \py!("ZB", cs)!, \py!("XB", co, ci)!, \py!("WB", cs)!, \py!("Bun", m, base)! and
\py!("ZP", dx, dy, c)! (\cref{sec:lean-export-net-zp}), where \py!cs!, \py!co!, \py!ci! are tuples of \term{chains} and a chain (type alias
\py!Chain = tuple[int, ...]!) is a tuple of multiplicity indices, outermost first. The helper
functions on kinds mirror Lean functions of the same meaning:

\begin{xltabular}{\linewidth}{@{}p{3.6cm}X@{}}
\toprule
Function & Behavior \\
\midrule
\texttt{\_\allowbreak{}head(kind)} & The constructor: the string itself, else the tuple's first entry. \\
\texttt{\_\allowbreak{}klean(kind)} & The Lean term: \texttt{.Z} for a string \texttt{Z};
  \texttt{.ZB [[m,...], ...]} and \texttt{.WB ...} with each chain written \texttt{[m$_1$, m$_2$]};
  \texttt{.XB co ci}; \texttt{.Bun m base} with \texttt{base} parenthesized unless it is a plain
  string; \texttt{.ZP dx dy (c)} for \texttt{("ZP", dx, dy, c)}. \\
\texttt{\_\allowbreak{}canon(kind)} & The kind with the chains of \texttt{ZB}, \texttt{WB}, \texttt{XB} sorted
  (recursively through \texttt{Bun}); a \texttt{ZP} kind unchanged. \\
\texttt{\_\allowbreak{}outcs(kind, n)}, \texttt{\_\allowbreak{}incs(kind, n)} & Mirrors of \texttt{Kind.outCs}, \texttt{Kind.inCs}: the
  chains of a copied generator's $n$ outputs (inputs); empty chains padded in front of the
  bundled chains; \texttt{Bun m b} prefixes $m$ to every chain of \texttt{b}. \\
\texttt{\_\allowbreak{}kms(kind)} & Mirror of \texttt{Kind.ms}: the set of multiplicity indices the kind reads. \\
\texttt{\_\allowbreak{}copyable(kind)} & Mirror of \texttt{copyKind}: false for \texttt{B}, \texttt{S} and \texttt{ZP}, recursive
  through \texttt{Bun}, true otherwise. \\
\texttt{\_\allowbreak{}zk(cs)}, \texttt{\_\allowbreak{}xk(co, ci)}, \texttt{\_\allowbreak{}wk(cs)} & \texttt{("ZB", cs)} (\texttt{("XB", co, ci)},
  \texttt{("WB", cs)}) when some chain list is nonempty, else \texttt{"Z"} (\texttt{"X"},
  \texttt{"W"}). \\
\texttt{\_\allowbreak{}tower(ms)} & Mirror of \texttt{Kind.tower}: \texttt{"I"} for no multiplicities, else
  \texttt{("Bun", ms[0], \_\allowbreak{}tower(ms[1:]))}. \\
\texttt{\_\allowbreak{}tower\_\allowbreak{}ms(kind)} & The multiplicities of a tower of copied bare wires, or \texttt{None} when the
  kind is not such a tower. \\
\bottomrule
\end{xltabular}

\subsection{Network records}\label{sec:lean-export-net-records}

\begin{signature}
\py!@dataclass(frozen=True) class LNode:!
\py!kind: Any; outs: tuple[int, ...]; ins: tuple[int, ...]; dim: int; ph: int!
\end{signature}
\noindent One network node. \py!lean()! returns the literal
\code!⟨!\textit{klean(kind)}\code!, [o₁, ...], [i₁, ...], !\textit{dim}\code!, !\textit{ph}\code!⟩!.
\py!labels()! returns \py!outs + ins!. \py!rename(f)! returns the node with every label passed
through \py!f!. For a binder \code!B! or splitter \code!S!, \py!dim! is the index of the first
dimension and \py!ph! the index of the second dimension.

\begin{signature}
\py!@dataclass(frozen=True) class LNet: nodes: tuple[LNode, ...]; summed: tuple[int, ...]; bnd: tuple[int, ...]!
\end{signature}
\noindent One network; \py!lean()! returns \code!⟨([n₁, ...] : List Node), [s₁, ...], [b₁, ...]⟩!.
The helpers \py!_list(xs)! (\code![x₁, x₂]!), \py!_nodes(ns)!
(\code!([n₁.lean(), ...] : List Node)!) and \py!_pairs(ps)! (\code![(a, b), ...]!) write the
components.

\begin{signature}
\py!@dataclass class Labelled: net: LNet; order: tuple[NodeId, ...]; port: dict[PortRef, int];!\\
\py!  name: str; scalar_note: str; source: Diagram | None = None!
\end{signature}
\noindent A diagram's network: the network, the diagram node id of each network node in
network order, every port's label, the Lean name (\code!N!$i$), the text of the diagram's scalar
(\py!str(d.scalar.to_sympy())!) and the diagram itself when it was translated with default
boundary labels and has no \bboxes{} (otherwise \py!None!).

\begin{signature}
\py!@dataclass(frozen=True) class _Named: text: str; model: sp.Expr; nonzero: str | None = None!
\end{signature}
\noindent A scalar written as a Lean definition applied to arguments: its Lean text, its sympy
model, and optionally a Lean proof that it is nonzero.

\begin{signature}
\py!@dataclass(frozen=True) class Proof: term: str; left: object; right: object!
\end{signature}
\noindent A Lean term proving \lean!Equal E left lhs right rhs!. \py!left! and \py!right! are
\term{product trees}: a sympy model, a \py!_Named!, or a triple \py!("*", a, b)! of two product
trees. \py!_Export.text(k)! writes a tree as \code!(!\textit{a}\code! * !\textit{b}\code!)!
recursively; \py!drop_ones(k)! writes it without the factors \code!1!; \py!nonzero(k)! proves it
nonzero by \lean!mul_ne_zero!, a \py!_Named!'s own proof or
\code!(by (try simp only [!\textit{scalar defs}\code!]); nzt)!, or \py!Scalars.nonzero_expr!;
\py!flat(k)! multiplies the models into one sympy expression.

\begin{signature}
\py!@dataclass class _Ctx: ex: _Export; before: Diagram; after: Diagram; step: RewriteStep;!\\
\py!  a: Labelled; new_nodes: dict[NodeId, LNode]; deltas: list[LNode]!
\end{signature}
\noindent What a rule translator sees of one step: the export, the diagrams before and after the
step, the step, the labeled start network, the new nodes (keyed by the new node id) labeled
with the start network's labels where the step's port mapping carries them, and the
\term{deltas}: one bare-wire node \lean!⟨.I, [x], [w], d, 0⟩! for every port-mapping entry
$\textit{old}\mapsto\textit{new}$ whose new port is not on a new node, with $x$ the old port's
label, $w$ the new port's label and $d$ the dimension index of $w$. The properties are
\py!match! (the step's match), \py!node(nid)! (the start network's node of that id),
\py!wire(w)! (the label of the wire's first end), \py!env! (the current environment text),
\py!rt(dz)! (the model $\sqrt{d}$ of dimension index \py!dz!, \py!Rt(dim_exprs[dz])!) and
\py!copy_scalar(dz, n)!.

\begin{definition}[Copy scalar]\label{def:lean-export-net-copy}
\py!copy_scalar(dz, n)! returns $(k, \textit{text}, \textit{proof})$ with
$k = \sqrt d$ for $n=0$, $k=1$ for $n=1$ and $k = (\sqrt d^{-1})^{n-1}$ otherwise, its
parenthesized Lean text, and the proof
\code!(show !$k$\code! * !$\sqrt d$\code! ^ !$n$\code! = !$\sqrt d$\code! from (!\textit{p}\code!))!
with $p$ from \py!Scalars.prove(k·√d^n, √d)!; a \py!ScalarUnsupported! becomes
\py!NetUnsupported!. It is the scalar of copying a state through an $n$-legged spider.
\end{definition}

\begin{signature}
\py!@dataclass class _Spec: l_ids: tuple[NodeId, ...]; r: Sequence[tuple[NodeId | None, LNode]];!\\
\py!  lemma: str; k: Any = 1; i_l: tuple[int, ...] | None = None; i_r: ... = None; l: tuple[LNode, ...] | None = None!
\end{signature}
\noindent A rule step's \term{local equation}: the ids of the matched nodes $L$ (in the lemma's
order), the replacement nodes $R$ with their diagram ids (\py!None! for a bare wire), the Lean
proof text of \lean!LocalEq E L iL R iR iface 1 k!, the step scalar $k$ (a model or product
tree), the labels summed inside $L$ and inside $R$ (computed when \py!None!), and the left nodes
when they differ from the start network's nodes of \py!l_ids!. A lemma text may contain the
placeholder \code!IFACE!, replaced by the interface list
(\cref{sec:lean-export-net-region}).

\begin{signature}
\py!@dataclass(frozen=True) class _MixLeg: ids: tuple[NodeId, NodeId]; lit: str; labs: tuple[int, int]; v: int!
\end{signature}
\noindent One leg of a mixed-dimension Z box (\cref{sec:lean-export-net-mixed}).

\begin{signature}
\py!@dataclass(frozen=True) class _Shim: fields: dict[str, Any]!
\end{signature}
\noindent A pattern match read through attribute names: \py!getattr(shim, name)! returns
\py!fields[name]!, else raises \py!AttributeError!.

\begin{signature}
\py!@dataclass class _BoxedMirror: nodes; wires; outs; ins; boxes!
\end{signature}
\noindent The Python mirror of a Lean boxed \lean!Diagram!
(\cref{sec:lean-export-net-boxed-literal}).

\section{The export state}\label{sec:lean-export-net-state}

One \py!_Export! object accumulates everything one export writes. Its constructor takes the
description function. \Cref{tab:lean-export-net-state} lists its fields with their initial
values; the following subsections specify how each table grows.

\begin{xltabular}{\linewidth}{@{}p{3.3cm}p{3.2cm}X@{}}
\caption{Fields of \texttt{\_\allowbreak{}Export}.}\label{tab:lean-export-net-state}\\
\toprule
Field & Initial & Meaning \\
\midrule
\endfirsthead
\toprule
Field & Initial & Meaning \\
\midrule
\endhead
\texttt{sym} & \texttt{\_\allowbreak{}Sym()} & Parameter names (\cref{sec:lean-export-net-sym}). \\
\texttt{dims}, \texttt{dim\_\allowbreak{}exprs} & \texttt{[]} & Lean text and sympy expression of each dimension index. \\
\texttt{dim\_\allowbreak{}key} & \texttt{\{\}} & Canonical expression text to dimension index. \\
\texttt{phases} & \texttt{["fun \_\allowbreak{} => 0"]} & Lean text of each phase index; index 0 is the zero vector. \\
\texttt{phase\_\allowbreak{}specs} & \texttt{[(\{\}, None)]} & Entries (index to turns) and rule of each phase index. \\
\texttt{phase\_\allowbreak{}key} & \texttt{\{"": 0\}} & Phase body text (and weight text) to phase index. \\
\texttt{phase\_\allowbreak{}tables} & \texttt{\{\}} & Phase index to the packed table (numerators, denominator) of an unweighted vector. \\
\texttt{poly\_\allowbreak{}phases} & \texttt{False} & Whether phase tables are fitted by digit polynomials (set by the exact refuter). \\
\texttt{phase\_\allowbreak{}polys}, \texttt{poly\_\allowbreak{}of} & \texttt{\{\}} & Phase index to its digit polynomial, and phase vector to its digit polynomial or \texttt{None}. \\
\texttt{table\_\allowbreak{}names}, \texttt{table\_\allowbreak{}defs} & \texttt{\{\}}, \texttt{[]} & Numerator tuple to its chunked table name \texttt{PT}$j$, and the chunk declarations. \\
\texttt{weight\_\allowbreak{}tab} & \texttt{\{\}} & Phase index to the Lean body of its coefficient vector. \\
\texttt{coef\_\allowbreak{}models} & \texttt{\{\}} & Weight body to its models (index to sympy model). \\
\texttt{labs} & \texttt{\{\}} & Label to dimension index. \\
\texttt{nb} & \texttt{None} & The boundary size shared by the unboxed top-level diagrams. \\
\texttt{next\_\allowbreak{}label} & 0 & The next fresh label. \\
\texttt{nets} & \texttt{\{\}} & (structure key, slots) to \texttt{Labelled}. \\
\texttt{boxed} & \texttt{\{\}} & Network name to (boxed diagram, node order). \\
\texttt{boxed\_\allowbreak{}mirror} & \texttt{\{\}} & Lean diagram name to \texttt{\_\allowbreak{}BoxedMirror}. \\
\texttt{literal\_\allowbreak{}info} & \texttt{\{\}} & Lean diagram name to (diagram, per-node data). \\
\texttt{original} & \texttt{None} & The \texttt{original} argument of \texttt{export}. \\
\texttt{wire\_\allowbreak{}alias} & \texttt{False} & Whether wires identify dimensions read through \texttt{dimIds}. \\
\texttt{raw\_\allowbreak{}ix}, \texttt{raw\_\allowbreak{}pairs} & \texttt{\{\}} & Identified dimensions' own indices and the indices they read as. \\
\texttt{decls} & \texttt{[]} & The emitted declarations, in order. \\
\texttt{checked} & \texttt{set()} & Lean diagram names with a \texttt{\_\allowbreak{}chk} theorem. \\
\texttt{count} & 0 & The counter of \texttt{fresh} names. \\
\texttt{unproved} & \texttt{[]} & Names of the emitted axioms. \\
\texttt{scalar\_\allowbreak{}gap} & \texttt{""} & Why the diagrams' own scalars are not compared, if they are not. \\
\texttt{bund}, \texttt{bund\_\allowbreak{}of} & \texttt{\{\}} & (base index, multiplicity index) to bundle index, and back. \\
\texttt{mult\_\allowbreak{}ix} & \texttt{\{\}} & Multiplicity expression text to multiplicity index. \\
\texttt{mult\_\allowbreak{}names}, \texttt{mult\_\allowbreak{}exprs}, \texttt{mult\_\allowbreak{}bare} & \texttt{[]} & Lean text, sympy expression and
  "is a bare symbol" of each multiplicity index. \\
\texttt{ks\_\allowbreak{}vals}, \texttt{ks\_\allowbreak{}syms} & \texttt{[]} & The current multiplicity values (Lean text, sympy). \\
\texttt{slots} & \texttt{None} & Boundary labels for the current translation, or default. \\
\texttt{alias} & \texttt{\{\}} & Dimension expression text to the expression assumed equal to it. \\
\texttt{shadow}, \texttt{shadow\_\allowbreak{}pairs} & \texttt{\{\}}, \texttt{[]} & A dimension's own index apart from its alias,
  and (own, canonical) index pairs. \\
\texttt{mults} & \texttt{set()} & Raw names of multiplicity symbols. \\
\texttt{sp\_\allowbreak{}binders}, \texttt{sp\_\allowbreak{}args}, \texttt{sp\_\allowbreak{}hyps}, \texttt{sp\_\allowbreak{}hyp\_\allowbreak{}args} & \texttt{""} & The parameter binders,
  arguments, positivity hypotheses and their arguments without multiplicities. \\
\texttt{scalar\_\allowbreak{}defs} & \texttt{[]} & Names of emitted scalar definitions (\texttt{\_\allowbreak{}sS}, \texttt{\_\allowbreak{}sG}). \\
\texttt{unfold} & \texttt{"E"} & The definitions a tactic unfolds to read the environment. \\
\texttt{family} & \texttt{False} & Whether a multiplicity exists (a family export). \\
\texttt{alg\_\allowbreak{}cert}, \texttt{alg\_\allowbreak{}everything} & \texttt{None}, \texttt{False} & The number-field certificate and
  whether every closed algebraic atom goes into it. \\
\texttt{box\_\allowbreak{}syms} & \texttt{set()} & Symbols of \bbox{} multiplicities in the theorem. \\
\texttt{linked} & \texttt{set()} & Induction indices proved as linked inductions (\cref{sec:lean-export-net-linked}). \\
\texttt{linked\_\allowbreak{}pending} & \texttt{[]} & (tag, induction theorem, start network, goal network) of each linked induction not yet written. \\
\texttt{fpos\_\allowbreak{}lemmas} & \texttt{""} & The \texttt{Fpos} application of the linked lemmas, kept before the binders change. \\
\texttt{in\_\allowbreak{}table} & \texttt{False} & Whether a node table is being written. \\
\texttt{sc} & \texttt{Scalars(...)} & The scalar writer (\cref{sec:lean-export-net-scalars}). \\
\texttt{env}, \texttt{binders}, \texttt{args}, \texttt{hyps}, \texttt{hyp\_\allowbreak{}args} & set by \texttt{finish} & The current
  environment text and the binders of the current theorem. \\
\bottomrule
\end{xltabular}

\subsection{Parameter names}\label{sec:lean-export-net-sym}

\py!_Sym! maps sympy symbols to Lean binder names. \py!params! maps a symbol's raw name to the
pair (Lean name, Lean type) with type one of \code!ℕ!, \leanx!ℝ!, \code!ℂ!, \leanx!ℤ!.

\begin{algorithm}[\py!_Sym.name(sym, kind)!]\label{alg:lean-export-net-name}
\begin{enumerate}
\item If the raw name is already registered, return its Lean name.
\item Replace every character outside \code![A-Za-z0-9_]! by \code!_!. If the result does not
  begin with a letter or lies in \py!_RESERVED!, prefix \code!p_!.
\item While the result equals a Lean name already taken, append \code!'!.
\item Register (result, \py!kind!) and return the result.
\end{enumerate}
\end{algorithm}

\py!kind_of(sym, default)! returns a registered symbol's type, else \py!default!. Three raw names
are registered with fixed Lean names: \code!k_lean_index! as \code!k : ℕ! (the index variable of
a phase or weight rule), \code!alg_theta! as \leanx!θ : ℝ!, and the induction index and step
symbols of a top-level induction as \code!n! and \code!k! (\cref{sec:lean-export-net-induction}).
A \bbox{} multiplicity symbol $m$ inside a node table is replaced by the symbol
\code!table_!$m$ (natural, nonnegative), so that a table body never mentions a multiplicity the
family quantifies over (\py!table_sym!, \py!table_expr!).

\begin{definition}[Rendering a sympy expression]\label{def:lean-export-net-render}
\py!_Sym.render(expr, ty, sym_kind)! writes \py!expr! as a Lean term of type
$\textit{ty}\in\{\mathbb N,\mathbb R,\mathbb C\}$, a free symbol taking the type \py!sym_kind!
unless already registered:
\begin{itemize}
\item $1 \mapsto$ \code!1!; an integer $v \mapsto$ \code!(v : ty)! or \code!(-|v| : ty)!;
  a rational $p/q \mapsto$ \code!((p : ty) / q)! (rejected for \code!ℕ!).
\item $i \mapsto$ \code!Complex.I! (only for \code!ℂ!); $\pi \mapsto$ \code!Real.pi! for
  \leanx!ℝ!, \code!(Real.pi : ℂ)! for \code!ℂ! (rejected for \code!ℕ!).
\item A symbol: its name when its type is \textit{ty}; \code!(name : ty)! for the casts
  $\mathbb N\to\mathbb R$, $\mathbb N\to\mathbb C$, $\mathbb R\to\mathbb C$,
  $\mathbb Z\to\mathbb R$, $\mathbb Z\to\mathbb C$; otherwise rejected.
\item A sum or product: the parenthesized terms joined by \code! + ! or \code! * !.
\item A power $b^e$: \code!(b ^ e)! for a nonnegative integer $e$; \code!(b ^ e')! for a
  symbolic $e$ that renders as a natural number; else rejected for \code!ℕ!;
  \leanx!(b ^ |e|)⁻¹! for a negative integer; for $e = p/2$ the root
  \code!(Real.sqrt b : ty)! (\code!Real.sqrt b! for \leanx!ℝ!) raised to $|p|$, inverted when
  $p<0$; any other power rejected.
\item $\exp(x) \mapsto$ \code!Complex.exp x! (only for \code!ℂ!).
\end{itemize}
Every rejection raises \py!NetUnsupported! (for example
\code!a rational 1/2 where a natural number is required!,
\leanx!symbol n of type ℝ used as ℕ!, \code!cannot write ... in Lean!).
\end{definition}

\subsection{Scalars}\label{sec:lean-export-net-scalars}

The field \py!sc! is a \py!Scalars! object of \module{archytaszx.semantics.lean_scalar}
(\cref{ch:lean-export}) with three callbacks: a natural symbol's Lean name is
\py!sym.name(table_sym(sym), "ℕ")!; another symbol's name and type are
\py!sym.name(table_sym(sym), kind_of(sym, "ℂ"))! and \py!kind_of(sym, "ℂ")!; and the positivity
proof of any natural expression is \py!_POSITIVE!. \py!_Export.scalar(s)! returns the model
\py!sc.model(s.to_sympy())! of an engine \py!Scalar!, converting \py!ScalarUnsupported! into
\py!NetUnsupported!. A model is a sympy expression over the atoms \py!Rt(x)! ($\sqrt x$, Lean
\lean!rt x!), \py!Om(m)! ($\omega_m$, Lean \lean!om m!), \py!Inv(e)! ($e^{-1}$), \py!NC(x)! (the
natural $x$ cast to $\mathbb C$), \py!II! ($i$), symbols and, with a number field, the root
$\theta$; a subterm without a model becomes an opaque complex symbol \code!opaque_!$j$.
\py!sc.text(model)! writes it as a \code!ℂ! term, \py!sc.nonzero_expr(model)! proves it nonzero,
and \py!sc.prove(lhs, rhs)! returns a tactic text \code!by ...! proving
\code!lhs = rhs! by \lean!ring! or \lean!linear_combination! over the defining relations
$\sqrt x^2 = x$, $e\cdot e^{-1}=1$, $\omega_m^m = 1$, $i^2 = -1$ and $p(\theta)=0$. The export
writes such a tactic as \code!(by ...)! where a term is wanted and strips the leading \code!by !
(the first three characters) where it splices the tactic into a \lean!simp only ...; all_goals (...)!
block.

\subsection{Dimension table}\label{sec:lean-export-net-dims}

Every dimension the export meets gets a \term{dimension index}: its position in \py!dims!.
\py!dim(d)! is \py!dim_expr(d.to_sympy())!. \py!dim_expr(raw)! takes the canonical expression
\py!canon(raw)!, returns its index from
\py!dim_key! when present, and otherwise appends \py!sym.render(e, "ℕ", "ℕ")! to \py!dims! and
\py!e! to \py!dim_exprs!, adds the free symbols of \py!e! to \py!sc.nats! and returns the new
index. Indices are assigned in order of first use, so the first dimension met is index 0.

\py!assume(a, b)! identifies two dimension expressions: with $r_a = $ \py!canon(a)! and
$r_b = $ \py!canon(b)! distinct, it sorts them by their text and records
\py!alias[str(hi)] = lo!. \py!canon(e)! follows \py!alias! from \py!str(e)! until it leaves the
table. All identified dimensions therefore share one index.

\py!raw_dim(d)! is the index of $d$ itself: \py!dim(d)! when $d$ is its own canonical form,
otherwise a \term{shadow index} appended once per expression (recorded in \py!shadow!) with the
pair (shadow index, canonical index) appended to \py!shadow_pairs!.

\py!raw_index(dim)! is the diagram's own index of a dimension: \py!dim(dim)! unless
\py!wire_alias! is set and $\textit{dim}$ is not canonical; then an index
\py!RAW_OFFSET!${}+j$ ($j$ the count of such dimensions so far), with
\py!raw_pairs[raw] = dim(dim)!.

\subsection{Phase and weight tables}\label{sec:lean-export-net-phases}

Every phase vector gets a \term{phase index}; index 0 is the zero vector \code!fun _ => 0!.
\py!phase(pv)! returns 0 for \py!None! or a zero vector; otherwise it sets \py!in_table! for the
duration of the call and runs \py!_phase!.

\begin{algorithm}[Phase body]\label{alg:lean-export-net-phase-body}
\py!phase_body(pv)! writes the phases of \py!pv! in turns as the body of a Lean
\code!fun k!:
\begin{enumerate}
\item If \py!phase_poly(pv)! gives a digit polynomial \textit{dp}, return
  \code!(((!\textit{dp}\py!.lean()!\leanx! : ℤ) : ℝ) / !\textit{dp}\py!.den!\code!)!.
\item If \py!_phase_table(pv)! gives a packed table (\textit{nums}, \textit{den}): with more than
  \py!TABLE_CHUNK! numerators return
  \code!(((!\py!table_get(nums)!\leanx! : ℤ) : ℝ) / !\textit{den}\code!)!, else
  \code!(((([!$n_0$\code!, !$n_1$\leanx!, ...] : List ℤ).getD k 0 : ℤ) : ℝ) / !\textit{den}\code!)!.
\item Otherwise, for each entry (index $i$, phase $\varphi$) of \py!pv.entries()! in ascending index order,
  render \leanx!table_expr(φ.to_sympy_turns())! as \leanx!ℝ!.
\item Fold from the last entry: start with \code!0! and wrap
  \code!if k = !$i$\code! then !\textit{text}\code! else !\textit{body}.
\item If \py!pv.rule! is not \py!None! (a phase given by a rule in the index variable), register
  \code!k_lean_index! as \code!k : ℕ!, render the rule with the index variable replaced by it, and
  return \code!(!\textit{body}\code!) + (if k = 0 then 0 else !\textit{rule}\code!)!.
\end{enumerate}
\end{algorithm}

\begin{definition}[Packed phase tables]\label{def:lean-export-net-packed}
\py!_phase_table(pv)! returns \py!None! when \py!pv! has a rule. Otherwise it returns the cached
value of \py!_PHASE_TABLES[pv]! (keyed by the phase vector's equality), or computes it: when
\py!pv! has more than \py!PHASE_TABLE_MIN! entries and every entry in turns is rational $p_i/q_i$,
the pair (\textit{nums}, \textit{den}) with $\textit{den} = \operatorname{lcm}_i q_i$ and
\textit{nums} the list of length $1 + \max i$ with $p_i\cdot\textit{den}/q_i$ at index $i$ and 0
at the indices without an entry; else \py!None!. The result (\py!None! included) is stored in
the cache, which is cleared first when it holds \py!_PHASE_TABLES_MAX! entries. The table is the
phase vector's entries as numerators over one denominator; the Lean body reads entry $k$ with
default 0.

\py!table_get(nums)! returns the Lean term of entry \code!k! of a long table. A new numerator
tuple gets the name \code!PT!$j$ ($j$ the number of tables named so far) and appends to
\py!table_defs! one declaration per chunk $c$ of \py!TABLE_CHUNK! entries,
\code!/-- Chunk !$c$\code! of a phase table. -/ def PT!$j$\code!_!$c$\leanx! : List ℤ := [...]!,
then \code!/-- A phase table in chunks of 256. -/ def PT!$j$\leanx! : List (List ℤ) := [PT!$j$\code!_0, ...]!.
The term is \code!(PT!$j$\code!.getD (k / 256) []).getD (k % 256) 0!.

\py!phase_poly(pv)! returns \py!None! unless \py!poly_phases! is set and \py!pv! has a concrete
dimension; then it returns, cached in \py!poly_of!, \py!digit_poly(nums, den, dim)! of
\module{archytaszx.semantics.digit_poly} for the packed table (\py!None! without one): a
polynomial $F$ with integer coefficients in the mixed-radix digits of $k$ (most significant
first, at most 4 digits, degree at most 3 per digit) with
$\textit{nums}_k \equiv F(k) \pmod{\textit{den}}$ for every $k < \dim$. Its Lean text
\py!dp.lean()! is \code!((!$\sum c\prod x_i^{e_i}$\code!) % !\textit{den}\code!)!, each digit
written \code!((k / !$b$\code! % !$b_i$\leanx! : ℕ) : ℤ)! with $b$ the product of the lower bases
(\code!k / b! for the first digit, \code!k % b_i! for the last).
\end{definition}

\py!_phase(pv)! computes the body and, for a weighted vector, its weight text
\py!weights(pv)!; the key is the body, or \textit{body}\code!|!\textit{weights}. A new key gets
the next index; the text \code!fun k => !\textit{body} goes to \py!phases!, the pair (entries in
turns, rule in turns) to \py!phase_specs!, and the weight text to \py!weight_tab! under the new
index. For an unweighted vector the new index also records its digit polynomial in
\py!phase_polys! when it has one, else its packed table in \py!phase_tables! when it has one; the
refutation exporters read both (\cref{ch:lean-export}).

\begin{algorithm}[Weight vector]\label{alg:lean-export-net-weights}
\py!weights(pv)! writes the coefficient vector of a weighted phase vector as the body of a Lean
\code!fun k!:
\begin{enumerate}
\item At a symbolic dimension, return \py!weights_rule(pv)!.
\item At the concrete dimension $d$: the models $c_k = $ \py!scalar(pv.coefficient(k))! for
  $1\le k<d$; fold from the largest $k$: start with \code!1! and wrap
  \code!if k = !$k$\code! then !\textit{text}($c_k$)\code! else !\textit{body}. Record the
  models in \py!coef_models[body]! and return the body. The coefficient at $k=0$ and beyond the
  dimension is therefore \code!1!.
\end{enumerate}
\py!weights_rule(pv)!: if the weight rule depends on the index variable, register
\code!k_lean_index! and write the rule's model with the index replaced; otherwise the body is
\code!1! (no rule) or the rule's model text. Each explicit weight entry wraps the body as above
in descending order. When the vector also has phase entries or a phase rule, the result is
\code!(if k = 0 then 1 else !\textit{body}\code!) * ex (!\textit{phase body}\code!)!; otherwise
\code!if k = 0 then 1 else !\textit{body}.
\end{algorithm}

\py!kind_ph(kind, pv, aux)! returns the Lean kind and phase index of a node: kinds other than
\code!Z!, \code!X! keep \py!aux! as their phase index; a \code!Z! or \code!X! spider gets
\py!phase(pv)! and becomes \code!ZW! or \code!XW! when \py!pv! is weighted.

\subsection{Label table}\label{sec:lean-export-net-labels}

Every port of a network carries a label; \py!labs! maps each label to the dimension index of the
wire it stands for. \py!label(dim)! returns \py!next_label!, increments it, records
\py!labs[lab] = dim! and calls \py!peel_reg(lab, H_OFFSET)!. \py!peel_reg(lab, off)! declares
the label \py!lab + off! with the base dimension of \py!lab!'s bundle when \py!lab!'s dimension
index is a bundle index (\cref{sec:lean-export-net-bundles}); it never overwrites an existing
entry. Boundary slot $i$ of a top-level diagram carries label $i$ (outputs first, then inputs),
and wire labels start at the boundary size.

\subsection{Multiplicities and bundles}\label{sec:lean-export-net-bundles}

A \bbox{} multiplicity gets a \term{multiplicity index}. \py!mult_of(expr)! returns the index of
\py!str(expr)!; a new bare symbol is added to \py!mults! and \py!sc.nats! and named
\py!sym.name(_mult_symbol(key), "ℕ")!; a new compound expression first registers each free
symbol (sorted by text) and is rendered as \code!ℕ!. The new index's text, expression and
bareness go to \py!mult_names!, \py!mult_exprs!, \py!mult_bare!; \py!family! becomes true and
\py!unfold! becomes \py!_FAMILY_UNFOLD!. \py!mult_index(name)! is \py!mult_of(_mult_symbol(name))!.

A \term{bundle} is the leg of a node inside $m$ copies seen from outside: one leg of dimension
$d^{\mu}$ standing for $\mu$ legs of dimension $d$. \py!bundle_dim(base, m)! returns the
dimension index of the bundle along multiplicity index $m$ of the dimension index \py!base!:
a new index whose \py!dims! text and \py!dim_exprs! entry copy those of \py!base! (the Lean
family computes the power, \cref{sec:lean-export-net-family-tables}), recorded in \py!bund! and
\py!bund_of!. \py!dim_for(base, chain)! applies \py!bundle_dim! along the chain from its
innermost entry outward. \py!chain(i)! returns the multiplicities an index bundles along,
outermost first; \py!root(i)! the plain index under it; \py!label_chain(x)! the chain of label
$x$'s dimension; \py!isb(m)! the predicate "the label's chain begins with $m$";
\py!leg_tails(n, skip)! maps each leg (side \code!"o"! or \code!"i"!, position) of $n$ whose label
chain, after dropping the first \py!skip! entries, is nonempty, to the rest of that chain.

\subsection{Number fields}\label{sec:lean-export-net-field}

Scalars that the scalar model has no atom for, such as algebraic numbers that are not roots of
unity or square roots, are written over one real algebraic number $\theta$.

\begin{algorithm}[\py!setup_field(theorem, everything)!]\label{alg:lean-export-net-field}
\begin{enumerate}
\item Collect \py!_field_values(theorem)!: every diagram scalar of the theorem, its premises and
  the diagrams of their inference data, every weighted spider's coefficients (indices
  $1,\dots,d-1$ at a concrete dimension, the explicit weights at a symbolic one) and every
  step's introduced scalar.
\item Model each value with a scratch \py!Scalars!; the wanted atoms are the opaque
  subterms. With \py!everything!, add every \py!numfield.atoms! of each closed value.
\item Keep the closed algebraic atoms (\py!_algebraic_atom!: no free symbols and every atom
  algebraic). With none, or when \py!numfield.certificate! raises, do nothing.
\item Register \code!alg_theta! as \leanx!θ : ℝ!. Each atom gets the form
  $\sum_k r_k\theta^k + i\sum_k s_k\theta^k$ from the certificate; the relation
  $p(\theta) = 0$ is proved by
  \leanx!(by have h := pevC_root θp θ hθ; simp only [pevC, θp] at h; push_cast at h; linear_combination h)!;
  inverses come from \py!numfield.inverse!. Install the \py!Field! in \py!sc.alg! and keep the
  certificate in \py!alg_cert!.
\end{enumerate}
\end{algorithm}

With a certificate, \py!set_params! appends \py!_THETA_HYPS! to the hypotheses and
\leanx!hθ hθlo hθhi! to their arguments, and the file begins with the tables of
\py!field_tables()!: \leanx!def θp : List ℤ! (the polynomial's coefficients, constant term first),
\leanx!θlo! and \leanx!θhi! (the isolating interval, as reals), and the theorem
\leanx|θ_isolated : ∃! x : ℝ, θlo < x ∧ x < θhi ∧ pev θp x = 0|, proved by \lean!root_unique!
with the derivative \leanx!pdv θp x! rewritten as the Bernstein form
$\sum_{j=0}^{n} w_j (x-\ell)^j (h-x)^{n-j}$ with $w_j = b_j\binom{n}{j}/(h-\ell)^n$ for the
certificate's Bernstein coefficients $b_j$ of the derivative (all positive), and positivity by
nested \lean!add_pos! of \lean!mul_pos! terms; the Bernstein basis certifies the sign of the
derivative on the interval~\cite{rouillier2004isolation}.

\section{From a diagram to a network}\label{sec:lean-export-net-translate}

\subsection{Node boxes and port boxes}\label{sec:lean-export-net-boxes}

\py!node_chain(d)! maps each node inside node-scoped \bboxes{} to the tuple of its boxes,
outermost first, and the chain of their multiplicity indices.

\begin{algorithm}[\py!node_chain!]\label{alg:lean-export-net-node-chain}
\begin{enumerate}
\item If any \bbox{} is extended (a chain, indexed, pair or copy-scalar box), raise
  \code!a chain, indexed, pair or copy-scalar bang box!.
\item For each node, sort its node-scoped boxes by depth (number of parents). Require that the
  outermost has no parent and each next box's parent is the previous box; else raise
  \code!node boxes that do not nest!.
\item Record (box ids, multiplicity indices of their multiplicities).
\end{enumerate}
\end{algorithm}

\subsection{Translation}\label{sec:lean-export-net-translate-alg}

\begin{algorithm}[\py!translate(d, slots, alloc, register=True, dims_out=None)!]\label{alg:lean-export-net-translate}
\begin{enumerate}
\item Compute the node chains. For each port in a port scope, collect its boxes (by ascending box
  id). The single top box (whose parent is not among them) must have as parent the innermost node
  box of the port's node (or none); the others must nest one inside the next. Otherwise raise
  \code!a port box outside its node's innermost box! or \code!port boxes that do not nest!. Record
  the port's \term{port chain}: the multiplicity indices of its boxes, outer first.
\item The boundary references are the boundary outputs followed by the boundary inputs. If
  \py!slots! is given and its length differs, raise \code!the boundary does not match!.
\item For boundary slot $i$ with port $r$: \textit{base} $=$ \py!dim(port dimension)!;
  \textit{want} $=$ \py!dim_for(base, own + extra)! with \textit{own} the node chain of $r$'s node
  and \textit{extra} the port chain of $r$. The label is \py!slots[i]! or $i$. With
  \py!register!, record \py!labs[lab] = want! (raise \code!a boundary slot changes dimension! when
  a different value is present) and call \py!peel_reg(lab, H_OFFSET)!. Record a nonempty
  \textit{extra} as the port's tail.
\item For each wire in the order of \py!Wire.sort_key()!: both ends must have the same dimension
  index (else \code!a wire joins dimensions !$d_a$\code! and !$d_b$); neither end may be in a port
  scope (else \code!a port box on a wire!). Of the two ends' node box tuples, one must be a prefix
  of the other (else \code!a wire between two node boxes!); the longer one's chain is the wire's
  chain and the shorter end, if its chain is shorter, gets as tail the rest of the chain. The
  label is \py!alloc(dim, chain)!; both ends carry it, and it joins the summed labels.
\item For each node in ascending id order: the Lean kind \py!zp_kind(node)! for a \code!Zp!
  generator, else from \py!_KINDS! (else \code!generator !$g$\code! has no Lean kind!); its output
  labels then its input labels (a port without a label raises
  \code!port !$r$\code! is neither wired nor on the boundary!); for a \code!Zp! node the dimension
  index is the modulus index of \py!zp_kind! and the phase index 0, for any other node the
  dimension and auxiliary index from \py!node_dims! and the kind and phase from \py!kind_ph!; the
  node with its tailed legs moved last by \py!_bundle_node!; for a node in node boxes, the kind
  wrapped in \code!("Bun", m, ...)! for each multiplicity of its chain from innermost to
  outermost, after checking \py!_copyable! (else \code!a boxed !$k$\code! node!).
\item Return the network (nodes, summed labels in wire order, boundary labels \py!slots! or
  $0,\dots,n-1$), the node order and the port labels.
\end{enumerate}
\end{algorithm}

\py!node_dims(d, nid)! returns (dimension index, auxiliary index): for a binder \code!B! the
indices of its two input dimensions; for a splitter \code!S! those of its two output dimensions;
otherwise (\py!dim(resolve_dim(node))!, 0). When \py!resolve_dim! raises
\py!DenoteDomainError! (legs of different dimensions), the set of the legs' (and the phase's)
dimension indices must have one element; otherwise, unless the export sets \py!mixed_ok!, raise
\code!node !$n$\code! has legs of different dimensions!; with \py!mixed_ok! the first leg's
index is used and every other index is recorded in \py!mixed! paired with it.

\subsubsection{Pair-phase nodes}\label{sec:lean-export-net-zp}

A \code!Zp! node of the engine (a node with a pair phase, \cref{ch:algebra}) becomes a node of
the Lean kind \lean!ZP!. \py!zp_kind(node)! returns the pair (\py!("ZP", dx, dy, c)!, $M$):
\begin{enumerate}
\item The legs are the outputs followed by the inputs. Fewer or more than two legs, or a phase
  that is not a \py!PairPhase!, raise \code!a Zp node without two legs and a pair phase!.
\item The coefficient is $c = \mathrm{cancel}(\varphi / (x\,y))$ with $\varphi$ the pair phase in
  turns (\py!phase.pair.to_sympy_turns()!) and $x$, $y$ the two index symbols \py!INDEX!,
  \py!INDEX2!. When $c$ still contains either symbol, or its numerator (after
  \py!sympy.together!) is not a nonzero integer, raise
  \code!a pair phase !$\varphi$\code! that is not c x y / M!.
\item $dx$ and $dy$ are \py!dim! of the two legs' dimensions; the numerator is the integer $c$;
  $M$ is \py!dim_expr! of the denominator.
\end{enumerate}
The node is written \code!⟨.ZP dx dy (c), outs, ins, M, 0⟩!; its tensor is $e(c\,x\,y/M)$ at the
two leg values (\cref{def:lean-export-net-env}). A \lean!ZP! node is not copyable, so a
\code!Zp! node inside a node box raises \code!a boxed !$k$\code! node! with $k$ the kind tuple.

\begin{algorithm}[\py!_bundle_node(n, tails)!]\label{alg:lean-export-net-bundle-node}
Given the tails (side and position to chain), move the bundled legs last in a fixed order:
\begin{enumerate}
\item No tails: return $n$. A weighted spider with tails raises
  \code!a weighted spider with bundled legs!.
\item \code!Z!: sort all tailed legs by (chain, side with outputs 0 and inputs 1, position);
  the unbundled outputs stay outputs; the inputs become the unbundled inputs followed by the
  tailed legs in that order; the kind is \code!("ZB", chains)!.
\item \code!X!: the unbundled outputs followed by the tailed outputs sorted by (chain, position);
  likewise for the inputs; the kind is \code!("XB", output chains, input chains)!.
\item \code!W! with no tailed input: the unbundled outputs followed by the tailed outputs sorted
  by (chain, position); the kind is \code!("WB", chains)!.
\item Any other case raises \code!a bundled leg on a !$k$\code! node!.
\end{enumerate}
\end{algorithm}

\subsection{Labeled networks and their names}\label{sec:lean-export-net-labelled}

\py!labelled(d, slots=None)! returns the cached \py!Labelled! for the key
(\py!structure_key(d)!, slots) or creates one:

\begin{enumerate}
\item When \py!slots! is \py!None! and $d$ has no \bboxes{}, its boundary size must equal the
  first such diagram's (\py!nb!; else \code!the diagrams have different boundaries!); the first
  one sets \py!nb! and raises \py!next_label! to at least \py!nb!. When \py!slots! is \py!None!,
  \py!next_label! is raised to at least the boundary size.
\item Translate with \py!alloc(dim, chain) = label(dim_for(dim, chain))!.
\item Name the network \code!N!$j$ with $j$ the number of cached networks; its source is $d$ when
  translated with default slots and without \bboxes{}. A boxed diagram translated with default
  slots is recorded in \py!boxed!.
\item Emit the declarations of \py!_net_def(doc, N!$j$\py!, net)! with \textit{doc} the
  \py!_comment! of \textit{describe(d)}.
\end{enumerate}

\py!_net_def(doc, name, net)! returns \code!/-- !\textit{doc}\code! -/! followed by
\code!def !\textit{name}\code! : Net := !\textit{net.lean()} when none of the node, summed and
boundary lists has more than \py!LIST_SPLIT! entries. Otherwise each list goes through
\py!_split_list! under the names \textit{name}\code!n!, \textit{name}\code!s!,
\textit{name}\code!b! (types \code!Node!, \code!ℕ!, \code!ℕ!), the chunk declarations come first
(nodes, then summed, then boundary labels), and the network is
\code!/-- !\textit{doc}\code! -/ def !\textit{name}\code! : Net := ⟨!\textit{nodes}\code!, !\textit{summed}\code!, !\textit{bnd}\code!⟩!,
the node list written \code!([...] : List Node)! when it was not split.
\py!_split_list(name, ty, items)! returns no declaration and the literal \code![...]! for at most
\py!LIST_SPLIT! items; otherwise one declaration
\code!/-- Chunk !$c$\code! of `!\textit{name}\code!`. -/ def !\textit{name}\code!_!$c$\code! : List !\textit{ty}\code! := [...]!
per run of \py!LIST_CHUNK! items and the term
\textit{name}\code!_0 ++ !\textit{name}\code!_1 ++ ...!.

\py!structure_key(d)! is the tuple (nodes, wires, boxes, boundary inputs, boundary outputs) with
nodes as (id, generator name, phase, input dimension texts, output dimension texts) in id
order, the phase compared as the engine's phase object (not its text), wires as their sorted
sort keys, and boxes as (sorted node scope, sorted port-scope sort keys) in box id order.
Diagrams of one structure share one network and name. Doc comments pass through
\py!_comment!, which cuts a text longer than \py!DOC_MAX! characters to its first
\py!DOC_MAX! characters followed by a space and the ellipsis character U+2026, then replaces \code!-/! by \code!- /! and
\code!/-! by \code!/ -!.

\section{Structure of the emitted file}\label{sec:lean-export-net-file}

\py!render(theorem)! returns \py!source(theorem.assumptions)!. \py!source(assumptions, lead="")!
assembles the file in this exact order:

\begin{enumerate}
\item When \py!original! is set, append \py!original_block()! to the declarations and clear
  \py!original! (\cref{sec:lean-export-net-original}).
\item The header comment: the lines \code!/-!, then
\begin{lstlisting}[style=grammar]
  ArchytasZX proof export over the tensor-network semantics of the ArchytasZX Lean
  library: every diagram is a Net, every rule step an instance of a lemma proved
  there. `#print axioms` lists any step stated without proof.
\end{lstlisting}
  then the line \py!lead! when nonempty, then the note lines of
  \cref{tab:lean-export-net-notes}, then \code!-/!.
\item An empty line, \code!import ArchytasZX!, an empty line,
  \code!namespace ArchytasZX.Export!, an empty line, \code!open ArchytasZX!, an empty line.
\item The options, one per line: \code!set_option linter.unusedVariables false!,
  \code!set_option maxRecDepth 100000!, \code!set_option maxHeartbeats 2000000!,
  \code!set_option linter.unusedTactic false!, \code!set_option linter.unreachableTactic false!,
  then an empty line.
\item The six tactic macros of \cref{sec:lean-export-net-macros}, each preceded by an empty line.
\item An empty line, then the tables (\cref{sec:lean-export-net-tables}) and the declarations in
  emission order, joined by one empty line each.
\item An empty line, \code!end ArchytasZX.Export!, and a final newline.
\end{enumerate}

\begin{xltabular}{\linewidth}{@{}p{4.3cm}X@{}}
\caption{Header note lines, in order (each indented two spaces).}\label{tab:lean-export-net-notes}\\
\toprule
Condition & Line \\
\midrule
each assumption $a$ & \texttt{assumes }$a$ \\
\texttt{alias} nonempty & \texttt{dimensions assumed equal are one dimension here: }$k_1$\texttt{ = }$v_1$\texttt{, ...}
  (each aliased text with its canonical expression, sorted by text) \\
\texttt{constraints} set & \texttt{assumes the dimension constraints }$c_1$\texttt{, ... (hypotheses of the theorem)} \\
number field & the line
  \texttt{$\theta$ is the real root of $\theta$p between $\theta$lo and $\theta$hi (theorem $\theta$\_\allowbreak{}isolated); the hypotheses h$\theta$, h$\theta$lo, h$\theta$hi state it},
  then one line \textit{form}\texttt{ stands for }\textit{atom} per atom \\
each opaque scalar & \textit{name}\texttt{ stands for the scalar }\textit{expr} \\
unproved steps & \texttt{unproved steps: }\textit{name}$_1$\texttt{, ...} \\
scalar gap & \texttt{the diagrams\textquotesingle{} own scalars are not compared: }\textit{reason} \\
\bottomrule
\end{xltabular}

\subsection{Tactic macros}\label{sec:lean-export-net-macros}

\py!_macros(unfold, pos_first)! emits each macro as the three lines
\code!set_option hygiene false in!, \code!/-- !\textit{doc}\code! -/!,
\code!macro "!\textit{name}\code!" : tactic => `(tactic| !\textit{body}\code!)!. The argument
\textit{unfold} is \py!unfold! (\code!E! for a plain export, \py!_FAMILY_UNFOLD! for a family);
\textit{pos\_first} is, for a family,
\code!exact Fam.env_dimv_pos _ _ _ (!\textit{FPOS}\code! _)!, else empty, where \textit{FPOS} is
\py!fpos_lemmas! when set and \py!fpos()! otherwise. \py!fpos()! is \code!Fpos! followed by
\textit{sp\_args} and the words of \textit{sp\_hyp\_args} that begin with \code!h_! (the
positivity hypotheses only; the number-field hypotheses are left out).
With $U$ the unfold list:

\begin{xltabular}{\linewidth}{@{}p{1.1cm}p{3.2cm}X@{}}
\toprule
Name & Doc comment & Body \\
\midrule
\texttt{chk} & Decides a structural side condition. &
  \texttt{first \textbar{} decide +kernel \textbar{} decide \textbar{} (simp only [}$U$\texttt{]; decide)} \\
\texttt{pos} & A dimension is positive. &
  \texttt{first \textbar{} }\textit{pos\_first}\texttt{ \textbar{} } (only when nonempty), then
  \texttt{(simp [U, $\Delta$] \textless{};\textgreater{} first \textbar{} decide \textbar{} omega \textbar{} positivity \textbar{} (apply Nat.mul\_\allowbreak{}pos \textless{};\textgreater{} omega) \textbar{} (apply pow\_\allowbreak{}pos; omega))} \\
\texttt{phz} & A phase vector is zero. & \texttt{(intro k \_\allowbreak{}; simp [U, $\Phi$, ex\_\allowbreak{}zero])} \\
\texttt{nzt} & A product of nonzero scalars is nonzero. &
  \texttt{repeat (first \textbar{} exact one\_\allowbreak{}ne\_\allowbreak{}zero \textbar{} apply mul\_\allowbreak{}ne\_\allowbreak{}zero \textbar{} apply pow\_\allowbreak{}ne\_\allowbreak{}zero \textbar{} apply inv\_\allowbreak{}ne\_\allowbreak{}zero \textbar{} exact Complex.I\_\allowbreak{}ne\_\allowbreak{}zero \textbar{} exact Complex.exp\_\allowbreak{}ne\_\allowbreak{}zero \_\allowbreak{} \textbar{} (apply rt\_\allowbreak{}ne; first \textbar{} omega \textbar{} norm\_\allowbreak{}num \textbar{} positivity \textbar{} (apply pow\_\allowbreak{}pos; omega)) \textbar{} (apply om\_\allowbreak{}ne\_\allowbreak{}zero; first \textbar{} omega \textbar{} norm\_\allowbreak{}num) \textbar{} (apply Nat.cast\_\allowbreak{}ne\_\allowbreak{}zero.2; first \textbar{} omega \textbar{} positivity \textbar{} norm\_\allowbreak{}num) \textbar{} norm\_\allowbreak{}num)} \\
\texttt{phs} & A phase vector is the sum of two. &
  \texttt{(intro k \_\allowbreak{}; simp only [U, $\Phi$]; rw [$\leftarrow$ ex\_\allowbreak{}add]; apply ex\_\allowbreak{}congr\_\allowbreak{}int; (try split\_\allowbreak{}ifs) \textless{};\textgreater{} first }
  followed by \texttt{\textbar{} (refine $\langle$}$z$\texttt{, ?\_\allowbreak{}$\rangle$; push\_\allowbreak{}cast; linarith)} for
  $z = 0, -1, 1, -2, 2$ in that order, then \texttt{)} \\
\texttt{zcz} & A coefficient vector is `[k = 0]`. &
  \texttt{(intro k hk; simp only [U, $\Delta$] at hk; simp only [zc, U, $\Phi$]; all\_\allowbreak{}goals (try simp only [$\Psi$]); all\_\allowbreak{}goals ((try split\_\allowbreak{}ifs) \textless{};\textgreater{} first \textbar{} (exfalso; omega) \textbar{} simp [ex\_\allowbreak{}zero]))} \\
\bottomrule
\end{xltabular}

\noindent Here \textit{U} stands for the unfold list, written in place.

\subsection{Tables}\label{sec:lean-export-net-tables}\label{sec:lean-export-net-family-tables}

Let \textit{binders} be \code!(!$n$\code! : !$T$\code!)! for every registered parameter except
\code!k_lean_index! and the multiplicity symbols, in registration order, and \textit{args} their
names. The tables are:

The list is built in this order and then extended:
\begin{enumerate}
\item The chunk declarations \py!table_defs! of the long phase tables
  (\cref{def:lean-export-net-packed}), in the order they were named.
\item \leanx!/-- Dimension values. -/ def Δ binders : ℕ → ℕ! with one line
  \code!  | !$i$\code! => !\textit{dims[i]} per index and the default line \code!  | _ => 1!.
\item \leanx!/-- Phase vectors in turns. -/ noncomputable def Φ binders : ℕ → ℕ → ℝ! with one
  line per phase index and \code!  | _ => fun _ => 0!.
\item The chunk declarations of \code!labs! (\py!_split_list!) when it
  has more than \py!LIST_SPLIT! entries (name \code!labs!, type \code!(ℕ × ℕ)!).
\item \code!/-- Each label's dimension index. -/ def labs : List (ℕ × ℕ) := [(l, i), ...]! sorted
  by label (or the \code!++! of its chunks).
\item \leanx!/-- The environment. -/ noncomputable def E binders : Env := { dimv := Δ args, phase := Φ args, labs := labs }!,
  with \leanx!, coef := Ψ args! before the closing brace when weights exist (double spaces
  collapsed to one).
\item When weights exist, \leanx!/-- Weight vectors. -/ noncomputable def Ψ binders : ℕ → ℕ → ℂ!
  with \code!  | !$i$\code! => fun k => !\textit{body} per weighted phase index (ascending) and
  \code!  | _ => fun _ => 1! is inserted at position 2 of the list (after \leanx!Φ! when there
  are no phase-table chunks).
\item Let \textit{first} be the entry now at position 0 (\leanx!Δ! when there are no phase-table
  chunks). When a number field exists, the four declarations of \py!field_tables()!
  (\cref{sec:lean-export-net-field}) are put in front.
\item For a family only, inserted directly after \textit{first}:
  \leanx!/-- Every dimension value is positive. -/ theorem Δpos binders hyps : ∀ i, 0 < Δ args i!
  proved by \leanx!intro i; unfold Δ; split <;> first | omega | positivity | (apply pow_pos; omega) | norm_num!,
  where \textit{hyps} are \code!(h_!$n$\code! : 0 < !$n$\code!)! for the natural parameters not
  named \code!table_*!.
\item For a family: \code!/-- Each bundle dimension index with its base index and multiplicity. -/ def bundL : List (ℕ × ℕ × ℕ) := [(i, b, m), ...]!;
  \leanx!/-- The family, its multiplicities in order: ... -/ noncomputable def F binders : Fam := { dims := Δ args, bund := bundL, phase := Φ args, labs := labs }!
  (with the \lean!coef! field as above; the doc lists \textit{raw}\code! = !\textit{name} per
  multiplicity in index order); \leanx!theorem Fpos binders hyps : ∀ i, 0 < (F args).dims i := fun i => Δpos args h_... i!
  (\code!(F )! written \code!F! without parameters);
  \code!noncomputable def FC : Fam := { dims := fun _ => 1, bund := bundL, phase := fun _ _ => 0, labs := labs }!
  (the family's structure alone); and
  \code!noncomputable def SpC (m : ℕ) : Spec := ⟨FC, m, ksOf []⟩!.
\end{enumerate}

In a family the environment of the current multiplicity values $v_0,\dots$ is the text
\py!env_text(vals)! $=$ \code!((F !\textit{args}\code!).env (ksOf [!$v_0$\code!, ...]))!, the
family's dimension of a bundle index $i$ with base $b$ along chain $(m_1,\dots)$ being
$\Delta(b)^{\prod_j \mathit{ks}(m_j)}$ (\lean!Fam.dimv!).

\subsection{A representative file}\label{sec:lean-export-net-example}

\begin{example}\label{ex:lean-export-net-fusion}
The diagram has two Z spiders on dimension 2: node 0 with input 0 and outputs 0, 1; node 1 with
input 0 and output 0; one wire from output 1 of node 0 to input 0 of node 1; boundary outputs
(node 0, output 0), (node 1, output 0) and boundary input (node 0, input 0). One
\code!spider_fusion! step merges them into node 2. The call
\py!export(kernel.rule(d, step))! returns the following file (macro section omitted; it is the
text of \cref{sec:lean-export-net-macros} with $U$ $=$ \code!E!).

\begin{center}
\begin{tikzpicture}[zxd]
  \node[zs] (a) at (0,0) {};
  \node[zs] (b) at (0.8,0.9) {};
  \draw[wire] (0,-0.8) -- (a); \draw[wire] (a) -- (-0.5,1.8);
  \draw[wire] (a) -- (b); \draw[wire] (b) -- (0.8,1.8);
  \node at (2,0.5) {$\Longrightarrow$};
  \node[zs] (c) at (3.3,0.4) {};
  \draw[wire] (3.3,-0.8) -- (c); \draw[wire] (c) -- (2.9,1.8); \draw[wire] (c) -- (3.7,1.8);
  \node[dimlabel] at (0.2,-0.6) {$2$};
\end{tikzpicture}
\end{center}

\begin{lstlisting}[style=leannet]
/-- Dimension values. -/
def Δ  : ℕ → ℕ
  | 0 => (2 : ℕ)
  | _ => 1

/-- Phase vectors in turns. -/
noncomputable def Φ  : ℕ → ℕ → ℝ
  | 0 => fun _ => 0
  | _ => fun _ => 0

/-- Each label's dimension index. -/
def labs : List (ℕ × ℕ) := [(0, 0), (1, 0), (2, 0), (3, 0)]

/-- The environment. -/
noncomputable def E : Env := { dimv := Δ , phase := Φ , labs := labs }

/-- a diagram of 2 nodes -/
def N0 : Net := ⟨([⟨.Z, [0, 3], [2], 0, 0⟩, ⟨.Z, [1], [3], 0, 0⟩] : List Node), [3], [0, 1, 2]⟩

/-- a diagram of 1 nodes -/
def N1 : Net := ⟨([⟨.Z, [0, 1], [2], 0, 0⟩] : List Node), [], [0, 1, 2]⟩

theorem step1 :
    Equal E 1 N0 1 N1 := by
  refine Equal.trans (isoB E 1 _ ⟨([⟨.Z, [0, 3], [2], 0, 0⟩, ⟨.Z, [1], [3], 0, 0⟩] : List Node), [3], [0, 1, 2]⟩ [] [0, 1] (by chk)) ?_
  have hR := rewrite E ([] : List Node) ([⟨.Z, [0, 3], [2], 0, 0⟩, ⟨.Z, [1], [3], 0, 0⟩] : List Node) ([⟨.Z, [0, 1], [2], 0, 0⟩] : List Node) [] [3] [] [0, 1, 2] [0, 2, 1] 1 1 1 (fusionZ E 0 0 0 0 3 [0, 3] [2] [1] [3] [0, 1] [2] [0, 2, 1] (by chk) (by chk) (by chk) (by chk) (by phs)) (by chk) (by chk) (by chk) (one_mul _).symm
  refine Equal.trans hR ?_
  refine isoB E 1 _ N1 [] [0] ?_
  chk

/-- The start and goal diagrams with their own scalars. -/
theorem main :
    Equal E 1 N0 1 N1 :=
  Equal.rescale (step1)
    one_ne_zero
    (by ring)

/-- a diagram of 2 nodes, as the engine's diagram. -/
def D0 : Diagram :=
  { nodes := [⟨0, .Z, 0, 0, [0, 0], [0]⟩, ⟨1, .Z, 0, 0, [0], [0]⟩],
    wires := [(⟨0, true, 1⟩, ⟨1, false, 0⟩)],
    outs := [⟨0, true, 0⟩, ⟨1, true, 0⟩],
    ins := [⟨0, false, 0⟩] }

/-- D0's wires, packed. -/
def WB0 : ℕ := 0x31

/-- D0's dimension index at each label, 1 bits each. -/
def LT0 : ℕ := 0x0

theorem D0_bd : DBridge.bridgeD 2 3 1 3 WB0 LT0 D0 N0 labs = true := by decide +kernel

theorem D0_net : D0.toNet 3 = N0 := (DBridge.bridgeD_ok 2 3 1 3 WB0 LT0 D0 N0 labs D0_bd).2

theorem D0_chk : D0.check labs 3 = true := (DBridge.bridgeD_ok 2 3 1 3 WB0 LT0 D0 N0 labs D0_bd).1

/-- a diagram of 1 nodes, as the engine's diagram. -/
def D1 : Diagram :=
  { nodes := [⟨2, .Z, 0, 0, [0, 0], [0]⟩],
    wires := [],
    outs := [⟨2, true, 0⟩, ⟨2, true, 1⟩],
    ins := [⟨2, false, 0⟩] }

/-- D1's wires, packed. -/
def WB1 : ℕ := 0x0

/-- D1's dimension index at each label, 1 bits each. -/
def LT1 : ℕ := 0x0

theorem D1_bd : DBridge.bridgeD 2 4 1 3 WB1 LT1 D1 N1 labs = true := by decide +kernel

theorem D1_net : D1.toNet 3 = N1 := (DBridge.bridgeD_ok 2 4 1 3 WB1 LT1 D1 N1 labs D1_bd).2

theorem D1_chk : D1.check labs 3 = true := (DBridge.bridgeD_ok 2 4 1 3 WB1 LT1 D1 N1 labs D1_bd).1

/-- The theorem on the engine's diagrams, through the proved translation `DEqual_iff`. -/
theorem main_diagrams :
    DEqual E 1 D0 1 D1 := by
  rw [DEqual_iff (o₁ := 3) (o₂ := 3) (L := labs) rfl D0_chk D1_chk, D0_net, D1_net]
  exact main
\end{lstlisting}

The boundary labels are 0, 1 (outputs) and 2 (input); the wire gets label 3. The step's region is
both nodes, the context is empty, the consumed wire 3 is summed in $L$, the interface is
$[0, 2, 1]$ (the region's unsummed labels in order of appearance), and $R$ is the merged spider.
The final \lean!isoB! renames nothing (the pairs list is empty) and permutes no node. In
\code!WB0!, the port key of $(0,\text{out},1)$ is $(2\cdot0+0)\cdot 2+1 = 1$ and of
$(1,\text{in},0)$ is $(2\cdot1+1)\cdot2+0 = 6$, with 3 key bits: $1 + 6\cdot 2^3 = 49 =$
\code!0x31!.
\end{example}

\subsection{Naming conventions}\label{sec:lean-export-net-names}

\begin{xltabular}{\linewidth}{@{}p{4.1cm}X@{}}
\toprule
Name & Declaration \\
\midrule
\texttt{N}$j$ & The network of the $j$-th distinct (structure, slots) pair translated. \\
\texttt{D}$j$ & The engine diagram of \texttt{N}$j$ (\texttt{Diagram} record). \\
\texttt{WB}$j$, \texttt{LT}$j$ & Packed wires and label table of \texttt{D}$j$. \\
\texttt{LB}$j$, \texttt{MN}$j$, \texttt{MS}$j$, \texttt{NP}$j$ & Bundled label table, node-id mask, box
  masks and plain network of a boxed \texttt{D}$j$. \\
\texttt{D}$j$\texttt{\_\allowbreak{}bd}, \texttt{\_\allowbreak{}bb}, \texttt{\_\allowbreak{}net}, \texttt{\_\allowbreak{}chk}, \texttt{\_\allowbreak{}ok} & Bridge theorems of
  \texttt{D}$j$. \\
\texttt{N}$j$\texttt{F}, \texttt{D}$j$\texttt{F} & Start and goal at their own dimensions
  (\cref{sec:lean-export-net-identified}). \\
\texttt{step}$c$ & A proved step; $c$ is the shared counter \texttt{count}. \\
\texttt{unproved\_\allowbreak{}}\textit{stem}\texttt{\_\allowbreak{}}$c$ & An axiom for a step without proof; \textit{stem} is
  the rule name (or \texttt{iso}, \texttt{algebraic}) with every character outside
  \texttt{[A-Za-z0-9\_\allowbreak{}]} replaced by \texttt{\_\allowbreak{}}. \\
\texttt{bun}$c$, \texttt{bun}$c$\texttt{L}, \texttt{bun}$c$\texttt{R} & The boxed rewrite at every
  multiplicity and its two region networks (\cref{sec:lean-export-net-lift}). \\
\texttt{ind}$c$ & A nested induction lemma; its helpers are \texttt{ind}$c$\texttt{\_\allowbreak{}sS} and so on. \\
\textit{name}, \textit{name}\texttt{M} & The final theorem; \texttt{M} marks the merged theorem when
  dimensions were identified. \\
\textit{name}\texttt{\_\allowbreak{}diagrams}, \textit{name}\texttt{\_\allowbreak{}original} & The restatements on the engine's
  diagrams and on the original diagrams. \\
\textit{pre}\texttt{sS}, \textit{pre}\texttt{sG}, \textit{pre}\texttt{zcase}, \textit{pre}\texttt{scase} &
  Start scalar, goal scalar, base case and step case of an induction (\textit{pre} empty at top
  level, \textit{name}\texttt{\_\allowbreak{}} for shifts and lifts). \\
\texttt{$\Delta$}, \texttt{$\Phi$}, \texttt{$\Psi$}, \texttt{labs}, \texttt{E}, \texttt{F}, \texttt{FC}, \texttt{SpC},
  \texttt{bundL}, \texttt{$\Delta$pos}, \texttt{Fpos} & The tables. \\
\texttt{EF}, \texttt{labsF}, \texttt{dimIds} & Tables of identified dimensions. \\
\texttt{$\Delta$o}, \texttt{$\Phi$o}, \texttt{$\Psi$o}, \texttt{Eo}, \texttt{Fo} & Tables of the original diagrams. \\
\texttt{$\theta$p}, \texttt{$\theta$lo}, \texttt{$\theta$hi}, \texttt{$\theta$\_\allowbreak{}isolated} & The number field. \\
\bottomrule
\end{xltabular}

\py!fresh(stem)! increments \py!count! and returns \textit{stem}\,\textit{count}; steps, axioms,
boxed rewrites and nested inductions share the counter.

\section{Proof terms}\label{sec:lean-export-net-proofs}

\subsection{Steps and axioms}\label{sec:lean-export-net-steps}

Let \textit{env} be the current environment text (\code!E!, \code!(E !\textit{args}\code!)! or a
family environment). \py!statement(sl, a, sr, b)! is
\code!Equal !\textit{env}\code! !$(sl)$\code! !$a$\code! !$(sr)$\code! !$b$, each scalar passed
through \py!_paren! (kept when it is a plain identifier or number of
\code![A-Za-z0-9_.']+! or one balanced parenthesized group, else wrapped in parentheses).
\py!_join(*parts)! joins the nonempty parts by single spaces.

\begin{itemize}
\item \py!theorem(a, b, k, tactic)! emits
  \code!theorem step!$c$\code! !\textit{binders hyps}\code! :!, then on the next line indented
  four spaces \textit{statement}(\code!1!, $a$, text($k$), $b$)\code! := by!, then each tactic
  line indented two spaces. It returns the proof
  \code!(step!$c$\code! !\textit{args hyp\_args}\code!)! with left scalar 1 and right scalar $k$.
\item \py!axiom(stem, note, a, b, k)! emits \code!/-- !\textit{note}\code! -/! and
  \code!axiom unproved_!\textit{stem}\code!_!$c$\code! !\textit{binders}\code! : !\textit{statement}(\code!1!, $a$, text($k$), $b$),
  records the name in \py!unproved!, and returns \code!(unproved_..._!$c$\code! !\textit{args}\code!)!
  with scalars 1 and $k$. An axiom takes no positivity hypotheses.
\end{itemize}

\subsection{Dispatch on the inference}\label{sec:lean-export-net-term}

\py!term(t)! (traced as \code!lean.export.step!) returns a \py!Proof! for the kernel theorem $t$
by its inference kind:

\begin{xltabular}{\linewidth}{@{}p{2.3cm}X@{}}
\toprule
Kind & Proof \\
\midrule
\texttt{refl} & \texttt{(Equal.refl }\textit{env}\texttt{ 1 N}$j$\texttt{)}, scalars 1, 1, with
  \texttt{N}$j$ the network of the left side at the current slots. \\
\texttt{symm} & \texttt{(}$p$\texttt{).symm} of the premise's proof $p$, scalars exchanged. \\
\texttt{trans} & \texttt{(Equal.mul\_\allowbreak{}trans }$p_1$\texttt{ }$p_2$\texttt{)}, left tree
  $(\ast, l_1, l_2)$, right tree $(\ast, r_1, r_2)$. \\
\texttt{iso} & \texttt{iso(lhs, rhs)} (\cref{sec:lean-export-net-iso}). \\
\texttt{algebraic} & \texttt{algebraic(lhs, rhs)} (\cref{sec:lean-export-net-algebraic}). \\
\texttt{rule} & \texttt{rule(t)} (\cref{sec:lean-export-net-rule}). \\
\texttt{induction} & in a family export, \texttt{linked\_\allowbreak{}induction(t)} when the index is in
  \texttt{linked} (\cref{sec:lean-export-net-linked}), else \texttt{nested\_\allowbreak{}induction(t)}
  (\cref{sec:lean-export-net-nested}). \\
\texttt{w\_\allowbreak{}fusion} & \texttt{w\_\allowbreak{}fuse(t)} (\cref{sec:lean-export-net-wfuse}). \\
other & raise \texttt{a }\textit{kind}\texttt{ inference has no network proof}. \\
\bottomrule
\end{xltabular}

\noindent The product trees of \code!trans! make the final scalar comparison exact: the proof of
a chain of steps $k_1, k_2, \dots$ proves \lean!Equal E (1*1*...) N (k₁*k₂*...) M!.

\subsection{Isomorphisms}\label{sec:lean-export-net-iso}

\py!iso(x, y)! labels both diagrams; equal names give \lean!Equal.refl!. Otherwise it computes
the label renaming \py!_iso_pairs(x, y, a, b)! and the term \py!iso_term("1", a.net, b.net, pairs, b.name)!
and emits a step whose tactic is \code!exact !\textit{term}. When either raises
\py!NetUnsupported!, it emits the axiom \code!unproved_iso_!$c$ with the note
\code!isomorphic comparison views (!\textit{reason}\code!)!.

\begin{algorithm}[\py!_iso_pairs(x, y, a, b)!]\label{alg:lean-export-net-iso-pairs}
\begin{enumerate}
\item Copy both diagrams, set their scalars to one and their parameters to none, and compute a
  node isomorphism with symmetric legs (\py!isomorphism(gx, gy, symmetric_legs=True)!); none raises
  \code!no node isomorphism!.
\item A port $r$ fits a port $r'$ when the isomorphism maps $r$'s node to $r'$'s node, the
  directions agree unless the node is a Z spider, and the positions agree unless the node is a Z
  or X spider.
\item For each wire of $x$ in sort-key order, take the first unused wire of $y$ in sort-key order
  whose ends fit the wire's ends in either orientation; map the wire's label to the label of that
  wire's first end. No such wire raises \code!no wire correspondence!.
\item The mapped labels must be exactly the summed labels of $x$'s network (else
  \code!the summed labels do not correspond!). Return \py!_perm_to_bijection(mapping)!.
\end{enumerate}
\end{algorithm}

\py!_perm_to_bijection(mapping)! extends a renaming to a bijection of its keys and values: a
non-injective mapping raises \code!the label renaming is not injective!; the values that are not
keys (ascending) are paired with the keys that are not values (in key order); the identity pairs
are dropped.

\begin{algorithm}[\py!iso_term(s, a, b, pairs, name=None)!]\label{alg:lean-export-net-iso-term}
Rename every node of $a$ by the pairs (\py!_ren!: the first pair with a matching key wins, other
labels fixed). For each node of $b$ in order, pick the first unused renamed node of $a$
equivalent to it by \py!_equiv_b!; record its position. A node without a partner raises
\code!the networks do not match node for node!. Return
\code!isoB !\textit{env}\code! !$(s)$\code! _ !\textit{name or b.lean()}\code! !\textit{pairs}\code! [!$p_1$\code!, ...] (by chk)!.
\end{algorithm}

\py!_equiv_b(n, m)! mirrors the Lean \lean!nodeEquivB!: the canonical normalized kinds, the
dimensions and the phases agree, and the leg groups of \py!_segments! have the same lengths and
agree group by group, a symmetric group up to order and another exactly.
\py!_norm(n)! mirrors \lean!Node.norm!: in a \code!ZB!, \code!XB! or \code!WB! node it moves the
tail legs with an empty chain to the plain legs (\py!_norm_tail!) and drops their chains, turning
the kind plain when no chain remains. \py!_segments(n)! of the normalized node gives the groups:

\begin{xltabular}{\linewidth}{@{}p{3cm}X@{}}
\toprule
Kind & Groups (legs, symmetric) \\
\midrule
\texttt{Z}, \texttt{ZW} & (outputs + inputs, yes) \\
\texttt{X}, \texttt{XW} & (outputs, yes), (inputs, yes) \\
\texttt{W} & (outputs, yes), (inputs, no) \\
\texttt{ZB cs} & (outputs + plain inputs, yes), then per distinct chain $c$ in sorted order
  (the tail inputs of chain $c$, yes) \\
\texttt{XB co ci} & (plain outputs, yes), per chain of \texttt{co} as above, (plain inputs, yes),
  per chain of \texttt{ci} \\
\texttt{WB cs} & (plain outputs, yes), per chain, (inputs, no) \\
\texttt{Bun m b} & \texttt{b = Z}: (outputs + inputs, yes); \texttt{b = X}: (outputs, yes),
  (inputs, yes); otherwise (outputs, no), (inputs, no) \\
other & (outputs, no), (inputs, no) \\
\bottomrule
\end{xltabular}

\py!_equiv(n, m)! mirrors the older \lean!nodeEquiv! (exact kinds; Z legs as one multiset, X
outputs and inputs as multisets, other nodes exact); the translators do not use it.

\begin{algorithm}[\py!_match(a, b, dims_a, dims_b, fix_bnd)!]\label{alg:lean-export-net-match}
Find a label bijection carrying network $a$ onto $b$ node for node, boundary to boundary (each
boundary label to itself when \py!fix_bnd!), or \py!None!:
\begin{enumerate}
\item Different node, boundary or summed counts give \py!None!.
\item The signature of a node is (canonical normalized kind, dimension, phase, the lengths and
  symmetry flags of its segments). Order $a$'s nodes by the number of $b$'s nodes with the same
  signature, fewest first.
\item With \py!fix_bnd!, fix every boundary label of $a$ to itself. A label $x$ may map to $y$
  when $x$ already maps to $y$, or $x$ is unbound, $y$ is unused, both are boundary labels or
  both are not, and \py!dims_a(x) == dims_b(y)!.
\item Search depth first: for the next node of $a$, try each unused node of $b$ with the same
  signature; align the segments in order, an asymmetric segment position by position, a
  symmetric segment by trying every distinct target for its first leg and recursing on the rest;
  undo the bindings on backtracking.
\item Return the forward map on success.
\end{enumerate}
\end{algorithm}

\subsection{Algebraic meets}\label{sec:lean-export-net-algebraic}

\py!algebraic(x, y)! proves $x = y$ for two diagrams of one structure whose spider coefficients
agree. The two networks must have the same node order (else \code!the diagrams have different nodes!).
For each node pair, unchanged nodes link by \code!(fun _ => rfl)!; a changed node keeps its
labels and takes the other's kind and phase, linked by
\code!(valCoefZ !\textit{env n r}\code! (by decide) (by decide) rfl rfl rfl !\textit{hc}\code!)!
for Z or weighted Z nodes, or \lean!valCoefX! with the function \lean!xc! for X nodes, where
\textit{hc} is \py!_coef_hyp! (\cref{sec:lean-export-net-coef}); any other changed node raises
\code!a changed node that is not a spider!. The links form
\code!(List.Forall₂.cons !$\ell_1$\code! (... List.Forall₂.nil))!. The wire renaming is the map
from each wire's label in $x$ to its label in $y$ (it must cover the summed labels). The step's
tactic is
\begin{lstlisting}[style=lean]
refine Equal.trans (netCoef ENV Na MID rfl rfl CHAIN) ?_
exact ISO
\end{lstlisting}
with \code!MID! the middle network literal and \code!ISO! the \lean!isoB! from it onto $y$'s
network. Any \py!NetUnsupported! yields the axiom \code!unproved_algebraic_!$c$
(\code!equal algebraic values (!\textit{reason}\code!)!).

\subsection{Rule steps}\label{sec:lean-export-net-rule}

\py!rule(t)! (traced as \code!lean.export.rule!) proves one rewrite step. Let \textit{step} be
the inference's \py!RewriteStep! and $k_e$ its introduced scalar's model (1 without a step).

\begin{algorithm}[\py!rule!]\label{alg:lean-export-net-rule}
\begin{enumerate}
\item Label the step's two diagrams at the current slots.
\item With a step but no translator in \py!_RULES!, try \py!_through_forward(t)!
  (\cref{sec:lean-export-net-forward}); a rebuilt theorem is proved by \py!term! instead.
\item Without a step or a translator, emit the axiom \code!unproved_!\textit{rule}\code!_!$c$
  with note \code!rule !\textit{rule} and scalar $k_e$.
\item Otherwise \py!rule_tactic!; a \py!NetUnsupported! gives the axiom with note
  \code!rule !\textit{rule}\code! (!\textit{reason}\code!)!. On success emit a step theorem with
  the translator's scalar.
\end{enumerate}
\end{algorithm}

\begin{algorithm}[\py!rule_tactic(before, after, step, a, b, builder)!]\label{alg:lean-export-net-rule-tactic}
\begin{enumerate}
\item If a new node or a removed node lies in node boxes, hand over to \py!bundled_tactic!
  (\cref{sec:lean-export-net-lift}).
\item Label each new node: a port that the step's port mapping hits takes the start label of the
  old port; any other port takes a fresh label per wire of \textit{after} (a port without a wire
  raises \code!new port !$r$\code! has no wire!). Kind and phase as in the translation; bundled
  legs moved by \py!_bundle_node! with \py!leg_tails!.
\item Form the deltas (\cref{sec:lean-export-net-data}).
\item Call the translator with the context; it returns a \py!_Spec!.
\item The left summed labels default to the labels of the consumed wires with both ends in the
  region; the right summed labels default to the sorted fresh labels; the left nodes default to
  the start network's nodes of \py!l_ids!. Replace \code!IFACE! in the lemma by the interface
  list.
\item Return the scalar and \py!region_lines!.
\end{enumerate}
\end{algorithm}

\subsection{The region tactic}\label{sec:lean-export-net-region}

\py!_iface(L, iL)! is the list of labels of the nodes $L$ not in $iL$, in order of first
appearance: the region's interface.

\begin{algorithm}[\py!region_lines(a, b, l_ids, L, R, r_ids, iL, iR, k, lemma)!]\label{alg:lean-export-net-region}
Let $C$ be the start network's nodes outside the region in network order, $iC$ the summed labels
not in $iL$, \textit{bnd} the boundary. If $|iC|+|iL|$ differs from the number of summed labels
or $iL$ repeats a label, raise \code!a consumed wire label is not summed!. With $K$ the text of
$k$:
\begin{enumerate}
\item \code!refine Equal.trans (isoB !\textit{env}\code! 1 _ ⟨C ++ L, iC ++ iL, bnd⟩ [] p (by chk)) ?_!:
  the start reordered so that the region comes last.
\item \code!have hR := rewrite !\textit{env}\code! C L R iC iL iR bnd iface 1 (K) (K) (LEMMA) (by chk) (by chk) (by chk) (one_mul _).symm!
  and \code!refine Equal.trans hR ?_! (\py!_stated!: the rewrite is stated by \lean!have!, its type
  elaborated without the goal, then chained in front of the goal).
\item While a node of the current network is a tower of copied bare wires (\py!_tower_ms! not
  \py!None!) with multiplicities $ms$, output $x$ and input $w$: reorder it last by
  \lean!isoB! at scalar $K$; then
  \code!have hR := rewrite !\textit{env}\code! REST [n] [⟨.I, [x], [w], bz, 0⟩] int [] [] bnd [x, w] (K) 1 (K) (bunI !\textit{env}\code! ms dim bz x w [x, w] (Fam.env_dimv_of _ _ _ _ _ rfl rfl ⟨rfl, rfl⟩) (by pos) (by chk) ×4) (by chk) ×3 (mul_one _).symm!
  and \code!refine Equal.trans hR ?_!
  with \code!bz! the dimension index of $x$; the tower becomes one plain bare wire.
\item While a node is a bare wire \code!I! with output $x$ and input $w$: if $w$ is not summed or
  $x = w$, raise \code!a bare wire joins two boundary slots!. Reorder by \lean!isoB! at scalar $K$
  into $\textit{pre} \mathbin{+\!\!+} [\text{wire}] \mathbin{+\!\!+} \textit{post}$ (the identity
  order), then
  \code!have hR := wire_elim !\textit{env}\code! (K) pre post int bnd x w dz (by chk) ×5! and
  \code!refine Equal.trans hR ?_!;
  rename $w$ to $x$ in the remaining nodes and erase $w$ from the summed labels.
\item The final \lean!isoB! onto $b$: the renaming from \py!_match_labels(current, ids, b)!,
  which pairs the current nodes with $b$'s nodes by diagram id (a missing id, a changed arity, an
  inconsistent label, a moving boundary label or a mismatch of summed labels raises); on failure,
  \py!_match(current, b.net, ..., fix_bnd=True)!, whose failure re-raises the first error. The
  term is \code!isoB !\textit{env}\code! (K) _ N!$j$\code! pairs p (by chk)!; the step ends with
  \code!refine isoB ... p ?_! and the line \code!chk!.
\end{enumerate}
\end{algorithm}

\subsection{Rules without a translator}\label{sec:lean-export-net-forward}

A step of a rule without a translator may still be proved as the reverse of a step of its
reverse rule.

\begin{algorithm}[\py!_through_forward(t)!]\label{alg:lean-export-net-forward}
\begin{enumerate}
\item Try \py!_tidy(t, t.rhs, [])!: the result alone, tidied back to the start.
\item The reverse rule of a name ending in \code!_converse! is \py!_FORWARD[name]! when present,
  else the name without the suffix; the reverse of any other rule is \textit{name}\code!_converse!.
  When the reverse equals the name or lacks a rule or a translator, return \py!None!.
\item For each match of the reverse rule in \code!t.rhs!, apply it; on success return
  \py!_tidy(t, result, [(t.rhs, step)])! when not \py!None!.
\end{enumerate}
\end{algorithm}

\begin{xltabular}{\linewidth}{@{}p{7.2cm}X@{}}
\caption{The table \texttt{\_\allowbreak{}FORWARD}.}\label{tab:lean-export-net-forward}\\
\toprule
Converse rule & Forward rule \\
\midrule
\texttt{x\_\allowbreak{}fusion\_\allowbreak{}converse} & \texttt{spider\_\allowbreak{}fusion} \\
\texttt{triangle\_\allowbreak{}inverse\_\allowbreak{}cancellation\_\allowbreak{}ti\_\allowbreak{}t\_\allowbreak{}converse} & \texttt{triangle\_\allowbreak{}inverse\_\allowbreak{}cancellation} \\
\texttt{fourier\_\allowbreak{}effect\_\allowbreak{}color\_\allowbreak{}change\_\allowbreak{}converse} & \texttt{fourier\_\allowbreak{}state\_\allowbreak{}color\_\allowbreak{}change\_\allowbreak{}swapped} \\
\texttt{fourier\_\allowbreak{}effect\_\allowbreak{}color\_\allowbreak{}change\_\allowbreak{}swapped\_\allowbreak{}converse} & \texttt{fourier\_\allowbreak{}state\_\allowbreak{}color\_\allowbreak{}change} \\
\texttt{triangle\_\allowbreak{}zero\_\allowbreak{}state\_\allowbreak{}t\_\allowbreak{}converse}, \texttt{triangle\_\allowbreak{}zero\_\allowbreak{}state\_\allowbreak{}ti\_\allowbreak{}converse} & \texttt{triangle\_\allowbreak{}zero\_\allowbreak{}state} \\
\texttt{connective\_\allowbreak{}states\_\allowbreak{}states\_\allowbreak{}converse}, \texttt{connective\_\allowbreak{}states\_\allowbreak{}effects\_\allowbreak{}converse} & \texttt{connective\_\allowbreak{}states} \\
\texttt{connective\_\allowbreak{}states\_\allowbreak{}swapped\_\allowbreak{}states\_\allowbreak{}converse}, \texttt{connective\_\allowbreak{}states\_\allowbreak{}swapped\_\allowbreak{}effects\_\allowbreak{}converse} & \texttt{connective\_\allowbreak{}states\_\allowbreak{}swapped} \\
\bottomrule
\end{xltabular}

\begin{algorithm}[\py!_tidy(t, current, steps)!]\label{alg:lean-export-net-tidy}
\begin{enumerate}
\item A diagram is \textit{reached} when it is isomorphic to \code!t.lhs!, its comparison view is
  isomorphic to that of \code!t.lhs!, or (when computable) its algebraic view is isomorphic to
  that of \code!t.lhs!.
\item Depth-first search from \textit{current}: return when reached; stop when the diagram has at
  most as many nodes as \code!t.lhs! or its view key was seen. Otherwise try each match of
  \code!identity_removal! then of \code!spider_fusion!, skipping steps that consume a node of
  \code!t.lhs!, and recurse.
\item With no path, or no steps at all, return \py!None!. Otherwise chain the kernel theorems of
  the steps by \py!kernel.trans!, close with \py!kernel.iso(current, t.lhs)! or, when that raises,
  \py!kernel.algebraic(current, t.lhs)!, and return \py!kernel.symm! of the chain; a
  \py!KernelError! returns \py!None!.
\end{enumerate}
\end{algorithm}

\subsection{The final theorem}\label{sec:lean-export-net-finish}

\py!finish(theorem, name)! (traced as \code!lean.export.finish!) writes the whole proof.

\begin{algorithm}[\py!finish!]\label{alg:lean-export-net-finish}
\begin{enumerate}
\item Set the environment to \code!E! and empty binders. An \code!induction! inference goes to
  \py!finish_induction!, a \code!shift! inference to \py!finish_shift!
  (\cref{sec:lean-export-net-induction}).
\item Register every variable of the theorem as a multiplicity; set up the number field.
\item Label the start and goal. \textit{merged} is true when dimensions were identified by
  deferred constraints (\py!alias! nonempty, \py!wire_alias! false) and the theorem has no
  variables; then compute the faithful networks of both sides (\cref{sec:lean-export-net-identified}).
\item Collect every dimension, phase and scalar of the theorem (\py!collect!: both sides, the
  premises recursively, the diagrams \code!peeled!, \code!other_peeled!, \code!own_k!,
  \code!other_k!, \code!replaced!, \code!exposed! of the inference data, and each step's scalar;
  the premises of an induction whose index is in \py!linked! are not visited).
  Model the start's and goal's scalars ($s_S$, $s_G$), or record that one has no model.
\item \py!set_params!: binders \code!(!$n$\code! : !$T$\code!)! of every parameter except
  \code!k_lean_index! and multiplicities; hypotheses \code!(h_!$n$\code! : 0 < !$n$\code!)! for
  the natural ones not named \code!table_*!; the number-field hypotheses; environment
  \code!(E !\textit{args}\code!)! or \code!E!.
\item In a family, keep these as the \code!sp_*! fields, take the multiplicity values to be their
  names, add binders \code!(!$m$\code! : ℕ)! for the bare multiplicities, and set the environment
  to \py!env_text(ks_vals)!. When aliases exist in a family, set \py!wire_alias! and clear
  \textit{merged}. A merged export names the theorem \textit{name}\code!M!.
\item Prove the theorem by \py!term!, obtaining the trees $l$, $r$. When linked inductions are
  pending, pass the proof through \py!link_tables! (\cref{sec:lean-export-net-linked}).
\item With both scalars modeled, prove $s_S\cdot r = s_G\cdot l$ by \py!sc.prove! and emit
\begin{lstlisting}[style=grammar]
/-- The start and goal diagrams with their own scalars. -/
theorem NAME BINDERS HYPS :
    Equal ENV (sS) Nstart (sG) Ngoal :=
  Equal.rescale PROOF
    NONZERO(l)
    (HST)
\end{lstlisting}
  where, when scalar definitions exist, the two last arguments are
  \code!(by (try simp only [!\textit{defs}\code!]); nzt)! and
  \code!(by simp only [!\textit{defs}\code!]; all_goals (!\textit{HST without "by "}\code!))!.
\item When a scalar has no model or the identity has no proof, record the reason in
  \py!scalar_gap! and emit instead the theorem with the trees $l$, $r$ written without unit
  factors and the doc comment \code!start = !$s_S$\code!, goal = !$s_G$, proved by
  \code!simpa only [one_mul, mul_one] using !\textit{proof}.
\item Restate: merged exports go to \py!unmerge!, family exports to
  \py!family_diagram_theorem!, plain exports to \py!diagram_theorem!
  (\cref{sec:lean-export-net-restate}). Return \py!render(theorem)!.
\end{enumerate}
\end{algorithm}

\section{Restatement on the engine's diagrams}\label{sec:lean-export-net-restate}

The final theorem speaks of networks. The restatement speaks of the engine's diagrams as Lean
\lean!Diagram! records and links them to the networks by the Lean-proved translation
\lean!DEqual_iff : E.labs = L → D₁.check L o₁ → D₂.check L o₂ → (DEqual E s₁ D₁ s₂ D₂ ↔ Equal E s₁ (D₁.toNet o₁) s₂ (D₂.toNet o₂))!.

\subsection{Diagram literals}\label{sec:lean-export-net-literal}

A \lean!Diagram! has fields \lean!nodes : List DNode! (each \lean!⟨id, kind, dim, ph, outs, ins⟩!
with \lean!outs!, \lean!ins! the dimension index of each leg), \lean!wires : List (Port × Port)!,
\lean!outs!, \lean!ins! (the boundary ports) and \lean!boxes! (default empty). A \lean!Port! is
\lean!⟨node, out, pos⟩!.

\begin{algorithm}[\py!diagram_literal(d, net, name, own=False, labs=None)!]\label{alg:lean-export-net-literal}
\begin{enumerate}
\item The wire offset \textit{off} is the first summed label, or the boundary size when there is
  none. When the summed labels are not \textit{off}, \textit{off}$+1$, \dots, return \py!None!.
\item For each node in id order (paired with the network's nodes): the leg dimension indices by
  \py!raw_dim! (with \py!own!) or \py!raw_index!; the node dimension is the network's, replaced by
  \py!raw_index! of the first leg's dimension when not \py!own!, the node has legs, it is not a
  binder or splitter, and the network's index equals \py!dim! of that leg. Write
  \code!⟨id, kind, dim, ph, [outs], [ins]⟩!.
\item Write the wires in sort-key order, the boundary outputs and inputs; a port is
  \code!⟨node, true|false, pos⟩!. Name the record \code!D!$j$ for the network \code!N!$j$ and emit
  \code!/-- !\textit{describe}\code!, as the engine's diagram. -/ def D!$j$\code! : Diagram := { nodes := ..., wires := ..., outs := ..., ins := ... }!.
  Without \py!own!, record the per-node data in \py!literal_info!.
\item The diagram expression is \py!dim_ids(dname)! (\cref{sec:lean-export-net-identified}).
\item Compute \py!_bridge_plan! with the leg dimensions read through \py!raw_pairs! and the
  label table \py!labs! (default \code!labs!/\py!self.labs!). Without a plan, emit
  \code!theorem D!$j$\code!_net : !\textit{dexpr}\code!.toNet !\textit{off}\code! = N!$j$\code! := by decide +kernel!.
  With the plan (\textit{width}, \textit{bits}, \textit{wb}, $b$, \textit{table}) emit
  \code!def WB!$j$\code! : ℕ := 0x...!, \code!def LT!$j$\code! : ℕ := 0x...! (doc
  \code!D!$j$\code!'s dimension index at each label, !$b$\code! bits each.!),
  \code!theorem D!$j$\code!_bd : DBridge.bridgeD ARGS = true := by decide +kernel!, and the two
  projections \code!D!$j$\code!_net : dexpr.toNet off = N!$j$ \code!:= (DBridge.bridgeD_ok ARGS D!$j$\code!_bd).2!
  and \code!D!$j$\code!_chk : dexpr.check L off = true := (... ).1!, with
  \code!ARGS = width bits b off WB!$j$\code! LT!$j$\code! dexpr N!$j$\code! L!; record the name in
  \py!checked!.
\item Return (\textit{dexpr}, \textit{off}).
\end{enumerate}
\end{algorithm}

\begin{algorithm}[\py!_bridge_plan(d, net, off, legs, labs, order=None)!]\label{alg:lean-export-net-bridge}
The Python mirror of \lean!DBridge.bridgeD!; it returns
(\textit{width}, \textit{bits}, \textit{packed}, \textit{eb}, \textit{tab}) or \py!None!.
\begin{enumerate}
\item A port key is (node id, is output, position). Let \textit{bnd} be the boundary ports,
  \textit{wires} the wire port pairs in sort-key order, $n_b$, $n_w$ their counts, and the node
  ids those of \py!order! or ascending. Return \py!None! when $n_b > \textit{off}$, the ids are not
  distinct (any order is accepted), the summed labels are not $\textit{off},\dots,\textit{off}+n_w-1$, the
  boundary labels are not $0,\dots,n_b-1$, or the node counts differ.
\item The ends are the boundary ports followed by both ends of every wire; they must be
  distinct.
\item For every node and every leg (outputs then inputs) with its dimension index and label:
  a label below \textit{off} must equal the leg's boundary slot, a label at or above it must
  equal \textit{off} plus the leg's wire index; one label must not carry two dimension indices.
  The number of legs must equal the number of ends.
\item Every label of \py!labs! that is a boundary slot or a wire label must carry the dimension
  index recorded for it, and there must be exactly $n_b + n_w$ such labels.
\item $\textit{width} = \max(1, \max_n \max(|\text{outs}|, |\text{ins}|))$;
  $\mathrm{pkey}(\textit{node},\textit{out},\textit{pos}) =
  (2\,\textit{node} + [\lnot\textit{out}])\cdot\textit{width} + \textit{pos}$;
  \textit{bits} the maximum bit length of the keys of the ends (at least 1).
\item Pack the wires from the last to the first:
  $\textit{packed} \leftarrow \mathrm{pkey}(a) + 2^{\textit{bits}}\mathrm{pkey}(b) +
  2^{2\textit{bits}}\textit{packed}$.
\item \textit{eb} is the maximum bit length of the recorded dimension indices (at least 1) and
  $\textit{tab} = \sum_\ell \textit{index}(\ell)\cdot 2^{\textit{eb}\cdot \ell}$.
\end{enumerate}
\end{algorithm}

The packed form makes the Lean check linear in the size of the diagram: \lean!bridgeD! decides by
\lean!decide +kernel! that the literal's \lean!toNet! is the network and that its \lean!check!
holds, reading wires and label dimensions from two natural numbers instead of searching lists.

\subsection{The diagram theorem}\label{sec:lean-export-net-diagram-theorem}

\py!diagram_theorem(name, binders, args, sl, a, sr, b, negated=False, env="", label="")! does
nothing when either network lacks a source diagram, or when aliases exist without
\py!wire_alias!, or when a literal returns \py!None!. Otherwise it emits
\begin{lstlisting}[style=grammar]
/-- The theorem on the engine's diagrams, through the proved translation `DEqual_iff`. -/
theorem LABEL :
    ∀ BINDERS, DEqual ENV (sl) Da (sr) Db := by
  intro ARGS
  rw [DEqual_iff (o₁ := OA) (o₂ := OB) (L := labs) rfl CHK_A CHK_B, Da_net, Db_net]
  exact NAME ARGS
\end{lstlisting}
with \code!LABEL! the given label or \textit{name}\code!_diagrams!, the \code!∀! and
\code!intro! lines omitted without binders, \code!CHK! the \code!_chk! theorem when the
literal has one and \code!(by decide +kernel)! otherwise, and \code!Db_net! omitted when both
sides are one diagram. With \py!negated! the statement is \code!¬ ∀ BINDERS, ...! and the proof is
\code!intro h; apply NAME; intro ARGS; have h := h ARGS; rw [...] at h; exact h!.

\subsection{Identified dimensions}\label{sec:lean-export-net-identified}

A rewrite step may assume that two dimension expressions are equal (a
\py!DimensionConstraint!), and a wire may join ports of two dimension expressions that the
theorem asserts equal. In both cases the dimensions share one index (\py!assume!), so the
networks are well typed.

\paragraph{Deferred constraints, merged export.} When the steps assumed equations
(\py!_deferred! returns the pairs (assumed, equal\_to) of \py!rewrite.rule.assumed!) and the
theorem has no variables, the main theorem is proved over the merged networks under the name
\textit{name}\code!M!, and \py!unmerge(theorem, outer, inner, start, goal)! states \textit{name}
over the networks at their own dimensions:
\begin{enumerate}
\item \py!faithful(d, lab)! rebuilds a network with every node and label at its own dimension's
  index (\py!raw_dim!; a binder's or splitter's two dimensions likewise); a label whose two ends
  have different own indices raises \code!a wire joins dimensions assumed equal!. No shadow pair
  raises \code!no dimension of the start or goal is assumed equal to another!; an existing
  \py!scalar_gap! raises \code!the diagrams' scalars are not compared!.
\item The constraints are the texts \leanx!(Δ_own) = (Δ_canonical)! of every shadow pair; the
  hypotheses are \code!(hc!$i$\code! : !$c_i$\code!)!.
\item Emit \code!def labsF!, the labels at their own indices;
  \leanx!noncomputable def EF binders : Env := { dimv := Δ args, phase := Φ args, labs := labsF, coef := ... }!;
  \code!def N!$j$\code!F : Net! for the start and (when different) the goal at their own
  dimensions; and
\begin{lstlisting}[style=leannet]
/-- The diagrams are equal wherever the assumed dimension equalities hold. -/
theorem NAME binders hyps (hc0 : c₀) ... :
    Equal EF (sS) NsF (sG) NgF :=
  Equal.of_mapDim' EF E PAIRS (by
    intro p hp
    simp only [List.mem_cons, List.not_mem_nil, or_false] at hp
    rcases hp with rfl | ... <;> simp [EF, Δ, hc0, ...])
    NsF NgF Ns Ng (by decide) (by decide)
    (by simp only [EF, E]; congr 1)
    (NAMEM args hyp_args)
\end{lstlisting}
  where \code!PAIRS! is the list of shadow pairs.
\item Unless in a family: when either own-index diagram fails \py!own_check! (its boundary slots
  and wire labels do not agree with \code!labsF!), set \py!wire_alias! for the call and restate
  the merged theorem through \lean!dimIds! under the label \textit{name}\code!_diagrams!;
  otherwise emit the diagram literals at own indices (names \code!D!$j$\code!F!, table
  \code!labsF!) and the theorem \textit{name}\code!_diagrams! stating
  \code!DEqual EF (sS) D!$j$\code!F (sG) D!$k$\code!F! with the constraint hypotheses, proved by
  \lean!DEqual_iff! with \code!(L := labsF)! and \code!exact NAME args hc0 ...!.
\end{enumerate}

\paragraph{Wire constraints, read-through.} When no step deferred a constraint but wires join
different dimension expressions (\py!_wire_constraints!), each pair is assumed and
\py!wire_alias! is set. The diagram literals then write each identified dimension at its own
index \py!RAW_OFFSET!${}+j$ and are read through \code!dimIds!: \py!dim_ids(dname)! emits once
\begin{lstlisting}[style=grammar]
/-- Dimensions identified with another (by a wire of the diagrams or by an equality the proof
assumes), each with the index it is read as: R is EXPR, read as index C (TEXT); ... -/
def dimIds : List (ℕ × ℕ) := [(R, C), ...]
\end{lstlisting}
and returns \code!(D!$j$\code!.mapDim (dimMap dimIds))!, \lean!dimMap ps x = (ps.lookup x).getD x!.

\subsection{The original diagrams}\label{sec:lean-export-net-original}

When \py!export! receives \py!original! $=(\sigma, \textit{sources})$, \py!original_block()!
restates the last emitted theorem whose name ends in \code!_diagrams! over the original diagrams,
whose dimension symbols $\sigma$ substituted.
\begin{enumerate}
\item With identified dimensions or a number field, raise
  \code!an original-diagram statement with identified dimensions!; without a
  \code!_diagrams! theorem, raise \code!no diagram statement to restate over the original diagrams!.
\item For every diagram name \code!D!$j$ in its statement, find the original paired with the
  restricted diagram (by identity or structure key; else \code!no original diagram for D!$j$).
  Node by node, the generator and leg counts must agree (else
  \code!the original diagram differs from the restricted one!). Each leg's index gets the
  original dimension: it must substitute to the export's dimension (else
  \code!an original dimension does not substitute to the export's!) and one index must not get two
  dimensions (\code!two original dimensions share one index!). A node dimension outside its legs
  takes the phase's dimension (\code!a node dimension without legs or phase!). A phase index gets
  the original phase vector (\code!a phase index without an original phase!,
  \code!two original phases share one index!). An original weight vector at an unweighted index
  raises \code!an original weight vector at an unweighted index!.
\item Emit \leanx!Δo!, \leanx!Φo! (and \leanx!Ψo!) with the substituted symbols as binders
  \code!(!$x$\code! : ℕ)! before the other parameters; the theorems \leanx!Δo_at!, \leanx!Φo_at!
  (and \leanx!Ψo_at!) stating that the original tables at the substituted values are the
  export's, proved by \code!funext! and case analysis on the index
  (\leanx!rcases i with _ | ... | i <;> first | rfl | (simp [Δo, Δ]; done) | (norm_num [Δo, Δ]; done)!
  and the analogues); the environment \code!Eo! or family \code!Fo!; the theorem
  \code!Eo_at! (\code!Fo_at!) by \code!simp only!; and
\begin{lstlisting}[style=grammar]
/-- The theorem on the original diagrams, their substituted dimension symbols free and equal to
the restriction's values. -/
theorem NAME_original (x : ℕ) (h_x_eq : x = v) ... :
    STATEMENT-over-(Eo x ...) := by
  subst h_x_eq
  ...
  simp only [Eo_at]
  exact NAME
\end{lstlisting}
\end{enumerate}

\section{Families and !-boxes}\label{sec:lean-export-net-families}

A theorem with \bboxes{} (\cref{ch:diagrams}) is a statement about every multiplicity. The
export writes one \term{family network}: a node inside node-scoped boxes along multiplicities
$m_1, \dots$ (outermost first) is one node of kind \code!Bun m₁ (Bun m₂ ... base)!, which stands
for all its copies; a leg crossing out of $\mu$ copies is one bundled leg of dimension
$d^{\mu}$; a port in port-scoped boxes is a bundled leg of a \code!ZB!, \code!XB! or \code!WB!
node. The Lean value of such a network at multiplicity values $\mathit{ks}$ is the value of the
expanded diagram (\lean!Diagram.expandK!, \cref{ch:lean-boxes}). Every multiplicity is a binder
\code!(!$m$\code! : ℕ)! of the theorem, and the environment is
\code!((F ...).env (ksOf [m₁, ...]))!.

\subsection{Split and zero}\label{sec:lean-export-net-split}

The induction on a multiplicity $m$ uses two network operations, mirrored in Python with the
Lean definitions' exact results.

\begin{definition}[Split]\label{def:lean-export-net-split}
\py!_split_net(net, m, isb, off)! mirrors \lean!Net.split (SpC m) off!: each node is replaced by
\py!_split_node(n, m, off)!, and every summed or boundary label $x$ that is a bundle along $m$
outermost is replaced by the pair $x, x+\textit{off}$. \py!_split_node!: a \code!ZB! node appends,
for each tail input whose chain begins with $m$, the label $b+\textit{off}$ to its inputs and the
chain without $m$ to its chains (\py!_peel!); \code!XB! does so on both sides, \code!WB! on its
outputs; a \code!Bun m b! node becomes itself followed by a copy of kind $b$ with every label
shifted by \textit{off} (the peeled last copy); other nodes stay.
\end{definition}

\begin{definition}[Zero]\label{def:lean-export-net-zero}
\py!_zero_net(net, m, isb)! mirrors \lean!Net.zero (SpC m)!: \code!ZB!, \code!XB!, \code!WB!
nodes drop their tail legs whose chain begins with $m$ (\py!_drop!, the kind becoming plain when
no chain remains); \code!Bun m ...! nodes vanish; the bundle labels along $m$ leave the summed and
boundary labels.
\end{definition}

\py!_rename_net(net, pairs)! mirrors \lean!Net.rename!. The offsets are \py!H_OFFSET!${}\cdot 2^{t}$
for the $t$-th split of one proof, so that peeled labels of nested splits never collide; before
a split at offset \textit{off}, \py!peel_reg(x, off)! declares every peeled label's dimension.

\py!slots_for(d, target)! finds boundary labels for a diagram so that it matches a given split
or zeroed network node for node: it translates $d$ without registering, with boundary labels
$2\cdot10^{9}+i$ and wire labels from $10^{9}+1$ upward, runs \py!_match(net, target, ..., fix_bnd=False)!
and returns the image of each slot (\code!a diagram does not match its family form! when none).
\py!iso_b(env, a_text, a, b_text, b)! is the \lean!isoB! term between two networks found by
\py!_match(..., fix_bnd=True)! (\code!two networks of the proof do not match!).

\subsection{Induction}\label{sec:lean-export-net-induction}

A kernel induction (\cref{ch:proofs}) has premises (base, expose, step) and data \code!index!
(the multiplicity $n$), \code!step! (the symbol $k$ with $n = k+1$), \code!side! (the side whose
last copy is peeled), \code!peeled!, \code!other_peeled!, \code!exposed!, \code!replaced! and
\code!region!.

\py!finish_induction(t, name)! sets \py!wire_alias! when aliases exist, registers the index as
\code!n : ℕ! and the step as \code!k : ℕ! (\py!prepare_induction(t, top=True)!), collects,
sets the parameters, keeps them as \code!sp_*!, and calls \py!induction_lemma(t, name, "")!.
\py!prepare_induction! raises \code!an induction over a fixed multiplicity! when the index is not
a variable of the theorem.

\begin{algorithm}[\py!induction_lemma(t, name, pre)!]\label{alg:lean-export-net-induction}
Let $n$ be the Lean name of the index, $k$ of the step, $mi$ the index's multiplicity index,
\textit{spc} $=$ \code!(SpC !$mi$\code!)!, $H={}$\py!H_OFFSET!. A non-bare multiplicity depending
on the index raises \code!a box whose multiplicity depends on the induction index!. The other
bare multiplicities become binders; \textit{ret}$(x, v)$ is \code!(!$x$\code! !\textit{args}\code! !$v$\code!)!.
\begin{enumerate}
\item Emit \code!noncomputable def !\textit{pre}\code!sS !\textit{binders}\code! (!$n$\code! : ℕ) : ℂ := !$s_S$
  and likewise \textit{pre}\code!sG!; add both to \py!scalar_defs!.
\item Base case: zero both family networks along $mi$; choose slots so that the base premise's
  start matches the zeroed start; set the multiplicity to \code!0!; prove the premise; bridge by
  \lean!isoB! from \code!(Net.zero spc Ns)! to the premise's start and from its goal to
  \code!(Net.zero spc Ng)!; emit
\begin{lstlisting}[style=grammar]
theorem PREzcase BINDERS HYPS :
    Equal ENV0 (PREsS ARGS 0) (Net.zero SPC Ns) (PREsG ARGS 0) (Net.zero SPC Ng) :=
  Equal.rescale (Equal.mul_trans I1 (Equal.mul_trans P0 I2))
    (by (try simp only [DEFS]); nzt)
    (by simp only [DEFS]; all_goals (HST))
\end{lstlisting}
  with \code!HST! proving $s_S(0)\cdot r = s_G(0)\cdot l$.
\item Step case: split the peeled side's network (and the other's) at $H$; slots for
  \code!peeled!; multiplicity value $k$; then the chain
  \lean!isoB! (split onto peeled) $\to$ expose premise $\to$ \lean!isoB! (exposed onto the plug
  form) $\to$ \code!plug! of the hypothesis $\to$ \lean!isoB! (onto \code!replaced!) $\to$ step
  premise $\to$ \lean!isoB! (\code!other_peeled! onto the other split). The plug: the region nodes
  of \code!exposed! are matched to the own side's network (\code!the hypothesis region does not match!),
  the renaming is extended with fresh labels for the other side's summed labels, and the term is
  \code!(plug ENV CTX (Net.rename π Nown).nodes (Net.rename π Noth).nodes iC (Net.rename π Nown).int (Net.rename π Noth).int bnd (Net.rename π Nown).bnd (own k) (oth k) (Equal.toLocal _ _ _ (Equal.rename ENV π Nown Noth (by chk) (by chk) ih)) (by chk) ×3)!
  with \code!ih.symm! when the goal side is peeled. Emit
\begin{lstlisting}[style=grammar]
theorem PREscase BINDERS HYPS (k : ℕ)
    (ih : Equal ENVk (PREsS ARGS k) Ns (PREsG ARGS k) Ng) :
    Equal ENVk (PREsS ARGS (k + 1)) (Net.split SPC H Ns) (PREsG ARGS (k + 1)) (Net.split SPC H Ng) :=
  Equal.rescale CHAIN (by ...; nzt) (by simp only [DEFS]; all_goals (HST))
\end{lstlisting}
\item The theorem:
\begin{lstlisting}[style=grammar]
/-- The start and goal families with their own scalars, for every `n`. -/
theorem NAME BINDERS HYPS :
    ∀ n : ℕ, Equal ENVn (sS ARGS n) Ns (sG ARGS n) Ng := by
  intro n
  induction n with
  | zero => exact Equal.of_zero_at FAM MI H [.., 0, ..] (by simp) Ns Ng
      (by decide : Spec.splitOK SPC H Ns = true) (by decide : Spec.splitOK SPC H Ng = true)
      (by decide) (PREzcase ARGS HYP_ARGS)
  | succ k ih => exact Equal.of_split_at FAM MI H k [.., k, ..] (by simp) Ns Ng
      (by decide : ...) (by decide : ...) BUNDPOS (PREscase ARGS HYP_ARGS k ih)
\end{lstlisting}
  with \code!BUNDPOS! $=$ \code!(fun _ _ _ => Fam.dimv_pos _ _ _ (!\textit{FPOS}\code! _))!, \textit{FPOS}
  $=$ \py!fpos()! (\cref{sec:lean-export-net-macros}).
  At top level (empty \textit{pre}), restate by \py!family_diagram_theorem! with the extra binder
  \code!(n : ℕ)!.
\end{enumerate}
\end{algorithm}

\subsection{Nested induction}\label{sec:lean-export-net-nested}

\py!nested_induction(t)! proves an induction inside a family proof as its own lemma: it saves the
environment, multiplicity values, slots and binders; registers the step symbol with the Lean name
\textit{n}\code!k! (primes appended until fresh); collects; emits \py!induction_lemma(t, ind!$c$\code!, ind!$c$\code!_)!;
restores the state; and returns the proof
\code!(ind!$c$\code! !\textit{sp\_args}\code! !\textit{other values}\code! !\textit{sp\_hyp\_args}\code! !\textit{current value}\code!)!
whose scalars are the \py!_Named! leaves \code!(ind!$c$\code!_sS ...)! and \code!(ind!$c$\code!_sG ...)!
with their models at the current values.

\subsection{Linked inductions}\label{sec:lean-export-net-linked}

The module \module{archytaszx.semantics.lean_famlink} proves an induction of a family export
whose goal's node phases read the induction index $n$ (a phase vector with the free symbol $n$).
The export writes such a phase through the phase parameter \code!table_!$n$
(\cref{sec:lean-export-net-sym}); a linked induction proves the lemma
\lean!theorem P : ∀ n : ℕ, Equal ((F n).env (ksOf [n])) s₀ X s₁ Y!
with that parameter set to $n$ itself, by induction on $n$, inside a namespace of its own, over
the export's tables and the normalizer's family lemmas (the section on symbolic-dimension
exports of \cref{ch:lean-export}; \cref{ch:lean-sym}).

\begin{signature}
\py!linked_indices(theorem: object) -> set[str]!\\
The index of every \code!induction! inference in \py!theorem! and its premises (depth first)
whose left or right side has a node phase with a free symbol of that name.
\end{signature}

\paragraph{Use in the export.} \py!export! reruns with \py!linked! $=$
\py!linked_indices(theorem)! when its first run fails (\cref{alg:lean-export-net-export}). In
that run \py!collect! does not visit the premises of a linked induction, and \py!term! sends it
to \py!linked_induction(t)!: with $mi$ the index's multiplicity index, it labels both sides at
the current slots, takes the tag \code!Lnk!$c$ (\py!fresh!), appends (tag, $t$, start network,
goal network) to \py!linked_pending!, and returns the proof
\code!(Lnk!$c$\code!.P !$v$\code!)!, $v$ the current value \py!ks_vals[mi]!, whose scalars are
the side scalars with the multiplicity expression replaced by its current symbol.
\py!finish! then calls \py!link_tables(proof)!:
\begin{enumerate}
\item \py!linked_lemmas()! builds the block of each pending induction by
  \py!Linked(self, start, goal, t.lhs, t.rhs, index, tag).build()!, appends it to the
  declarations, and empties \py!linked_pending!.
\item For each linked index $n$ (sorted) with a registered parameter \code!table_!$n$ and a
  multiplicity index, map that parameter's Lean name to the multiplicity's Lean name.
\item Keep \py!fpos()! as \py!fpos_lemmas!; drop the mapped parameters from \py!binders! and
  \py!args! (the binder \code!(!$x$\code! : ℕ)! and every whole-word occurrence, then collapse
  spaces); replace each whole-word occurrence by the multiplicity name in \py!sp_args!, \py!env!
  and the proof term.
\end{enumerate}
The final theorem, \textit{name}\code!_diagrams! and every point claim of the family are then
statements at $F\,n$ with the phase parameter equal to the index.

\paragraph{The lemma.} \py!Linked! extends \py!lean_sym._SymExport! and reuses its tables,
networks and family lemmas (\py!tables!, \py!texts!, \py!fam!, \py!scal_at!, \py!ratio!,
\py!call!, \py!env_at!, \py!prelude!). Its constructor raises \py!NetUnsupported! with:
\code!a linked induction beside other multiplicities! (the bare multiplicities of the export are
not exactly the index); \code!the induction index reads no phase table! (no parameter
\code!table_!$n$); \code!a linked induction whose start's phases read the index!;
\code!the goal's phases do not read the index!; and
\code!a linked induction with a common scalar or constraints! (the normalizer tables of the two
side scalars give a common scalar or dimension constraints). For each goal node whose phase
vector $\varphi$ reads $n$, with $p$ its phase index, it records the row maps
$g_{\mathrm{step}}(p) = \mathrm{phase}(\varphi[n \mapsto n-1])$ and
$g_{0}(p) = \mathrm{phase}(\varphi[n \mapsto 0])$.

\begin{algorithm}[\py!Linked.build()!]\label{alg:lean-export-net-linked}
Let $X$, $Y$ be the start and goal networks, $m$ the index's multiplicity index,
$\textit{spc} = $ \code!(SpC !$m$\code!)! and $H = 2^{\max(16,\ \mathrm{bitlen}(\max \text{label}))}$.
\begin{enumerate}
\item A phase row read by $X$ whose entries mention $n$ raises
  \code!the start reads a phase row of the index!. Every label of $X$ and $Y$ whose chain begins
  with $m$ is declared at offset $H$ (\py!peel_reg!).
\item $Y_0$ is $Y$ with the phase rows mapped by $g_0$ (\py!_map_ph!, mirror of
  \lean!Net.mapPh (phMap g)!, binders and splitters untouched); the inverse map must be one to
  one and carry $Y_0$ back to $Y$ (else
  \code!the goal's phases at zero share a row with its other phases!), and likewise the zeroed
  networks (\code!the zeroed goal does not read its rows back!). $Y_s$ is $Y$ mapped by
  $g_{\mathrm{step}}$. Emit \code!def L!$i$\code!z! (the zero network of $Y_0$, its phases at the
  index 0) and \code!def L!$i$\code!y! ($Y_s$, its phases at the index minus 1), $i$ the number
  of family lemmas so far.
\item The zero lemma: the normalizer's family lemma from \code!(Net.zero spc X)! to
  \code!L!$i$\code!z! at the scalars at 0.
\item The spider: exactly one \code!ZB! node $z$ of $X$ with an input label bundled along $m$
  and summed (else \code!a linked node box meeting other than one Z spider!); no \code!Bun!
  node along another multiplicity (\code!towers along two multiplicities!); $z$ has a boundary
  output $o$ (\code!the spider meeting the node box has no boundary output!).
\item Split $X$ and $Y$ at $H$. The split of $X$ must contain the nodes of $X$ other than $z$
  (\code!the split network does not contain the network!) and exactly one \code!ZB! node $z_s$
  with $z$'s outputs (\code!the split spider is not found!), whose inputs are $z$'s inputs
  followed by extra legs and whose chains are $z$'s followed by empty chains (else
  \code!the split spider's legs are not the unfused form!). With a fresh label $w$, $z_s$ unfuses
  into $z_a$ ($z$ with output $o$ replaced by $w$) and $z_b$ $=$ \code!⟨.Z.., [o], [w, extra...], dim, 0⟩!.
\item Emit \code!def L!$i$\code!s! (the split, its spider last), \code!def L!$i$\code!u! (the
  spider unfused) and \code!def L!$i$\code!r! (the rest after the hypothesis); the renaming
  $\pi$ maps $o \mapsto w$ and every summed label of $Y_s$ to a fresh label
  (\py!_perm_to_bijection!).
\item The step theorem
\begin{lstlisting}[style=grammar]
/-- The step at `n = k_ + 1`. -/
theorem Q BINDERS (k_ : ℕ)
    (ih : Equal ENVk SK0 X SK1 Y) :
    Equal ENVk1 SK10 (Net.split SPC H X) SK11 (Net.split SPC H Y) :=
  Equal.congr_s (Equal.mul_trans (Equal.mul_trans (Equal.mul_trans (Equal.mul_trans
    (Equal.mul_trans ISO1 UNFUSE) ISO2) PLUG) ISO3) LT) (by ring) (by ring)
\end{lstlisting}
  where \code!ISO1! is \lean!isoB! from the split onto \code!L!$i$\code!s!, \code!UNFUSE! the
  \code!plug! of the reversed \lean!fusionZG! instance, \code!ISO2! the reorder onto the plug form,
  \code!PLUG! the \code!plug! of the hypothesis \code!ih! carried from $F\,k$ to $F\,(k+1)$ by
  \lean!Equal.of_phRows! (the goal's rows read through $g_{\mathrm{step}}$, the dimension tables
  equal by \lean!Fam.env_dimv_congr!) and renamed by $\pi$, \code!ISO3! the reorder onto
  \code!L!$i$\code!r!, and \code!LT! the normalizer's family lemma from \code!L!$i$\code!r! to the
  split goal at the phase parameter $k+1$, with the per-copy ratios of both scalars.
\item The induction
\begin{lstlisting}[style=grammar]
/-- The induction, the family's phase parameter at the index. -/
theorem P BINDERS : ∀ n : ℕ, Equal ENVn S0 X S1 Y := by
  intro n
  induction n with
  | zero => exact Equal.of_zero_at FAM0 M H [..] (by simp) X Y OK0 OK1 (by decide +kernel) ZERO
  | succ k_ ih => exact Equal.of_split_at FAM1 M H k_ [..] (by simp) X Y OK0 OK1
      (fun _ _ _ => Fam.dimv_pos _ _ _ FPOS1) (Q k_ ih)
\end{lstlisting}
  with \code!ZERO! the zero lemma carried by \lean!Equal.of_phRows! through the inverse of $g_0$,
  \code!OK0!, \code!OK1! the split checks of $X$ and $Y$, the binders without the phase
  parameter, and \code!FPOS1! $=$ \code!(Fpos ARGS (by omega : 0 < d) ... _)! at $k+1$.
\item Return \code!namespace Lnk!$c$, the normalizer prelude, its scalar tables, the network
  declarations, the family lemmas (last made first), the two theorems and \code!end Lnk!$c$.
\end{enumerate}
\end{algorithm}

\lean!Equal.of_phRows! (\file{ArchytasZX/PhaseMap.lean}, \cref{ch:lean-sym}) is written by
\py!transfer(e0, e1, g, n, N, m, M, m1, hd)!:
\code!Equal.of_phRows e0 e1 (phMap []) (phMap PAIRS) ROWS_N ROWS_M hd rfl rfl KEEP (fun _ _ => rfl) MOVED (fun _ _ => rfl) n m n m1 (by decide +kernel)×4!,
\code!PAIRS! the sorted row map, \code!ROWS! the sorted phase rows of the networks' nodes other
than binders and splitters (\py!_rows!), \code!KEEP! and \code!MOVED! the proofs that each row
agrees (\py!_mem_cases!: \code!(by intro q hq; simp only [...] at hq; rcases hq with rfl | ... <;> T)!,
\code!(fun q hq => by simp at hq)! for no row) with $T$ $=$ \code!rfl! for the kept rows and
\py!_ROW_EQ! $=$ \leanx!first | rfl | (funext j; simp [Fam.env, F, Φ, phMap] <;> split_ifs <;> push_cast <;> ring)!
for the moved rows.

\subsection{Shifts and lifts}\label{sec:lean-export-net-shift}

A \code!shift! inference proves a family at $n = a k + c$ from its peeled successors at
$a k$; with \code!scale! $a = 0$ it is a \term{lift} under several multiplicity clauses.

\py!shift_lemma(t, name)! (scale $a\ge1$, count $c$, default 1): the scalar definitions as for an
induction; the networks split $c$ times at offsets $H, 2H, 4H, \dots$ ($H\cdot2^j$); slots for
\code!peeled! on the last split; the multiplicity value \textit{base} $=k$ or \code!(a * k)!; the
chain \lean!isoB!, body premise, \lean!isoB!; the theorem \textit{name}\code!_scase! at
\textit{base}$\,+c$; then \textit{name} itself by $c$ nested \lean!Equal.of_split_at! from the
innermost split outward, with doc comment
\code!The start and goal families with their own scalars, at `!$n$\code! = !\textit{base}\code! + c`.!
A non-bare multiplicity depending on the index raises
\code!a box whose multiplicity depends on the shifted index!. The closing scalar goal tries
\code!(simp [pow_mul']; done)!, then \code!(ring_nf; simp [pow_mul']; done)!, then the
\py!sc.prove! tactic.

\py!lift_lemma(t, name)! (scale 0): each clause (index, scale, $c$, step symbol or none) is
handled in order of increasing nesting depth of its multiplicity among the labels; for each, the
networks are split $c$ times and, when the clause has no step symbol, zeroed once, each at a fresh
offset $H\cdot2^{t}$ (with $t$ starting past the bit length of the largest label when boxes nest).
The proof applies \lean!Equal.of_zero_at! and \lean!Equal.of_split_at! from the last operation to
the first around the \code!_scase! theorem. Errors: \code!an instance at a fixed multiplicity!,
\code!a box whose multiplicity depends on a lifted multiplicity!,
\code!a lifted box inside a box at n >= c!.

\subsection{Rewriting inside node boxes}\label{sec:lean-export-net-lift}

A rule step whose new or removed nodes lie in node boxes is proved on the peeled copies and
lifted to every multiplicity by induction.

\begin{algorithm}[\py!bundled_tactic!]\label{alg:lean-export-net-bundled}
\begin{enumerate}
\item Require a family export with multiplicity values and exactly one (boxes, chain) among the
  touched nodes (else \code!a boxed rewrite across boxes!); every new node must have that box
  tuple (else \code!a boxed rewrite reaching outside its box!).
\item Label new nodes as in \cref{alg:lean-export-net-rule-tactic}, fresh labels bundled along the
  chain, legs bundled beyond the chain by \py!leg_tails(..., len(chain))!, kinds wrapped in
  \code!Bun! along the chain.
\item Each delta's new label must bundle exactly along the chain (else
  \code!a boxed rewrite rewiring a bundled leg!); the delta node is the tower \py!_tower(chain)!
  on the root dimension.
\item \py!lift(chain, 0, ...)! returns the lifted local equation; \py!region_lines! does the rest.
\end{enumerate}
\end{algorithm}

\begin{algorithm}[\py!lift(chain, depth, view, ...)!]\label{alg:lean-export-net-lift}
\begin{enumerate}
\item Empty chain: call the translator on the view; fill the default summed labels; replace
  \code!IFACE!; return the spec.
\item Otherwise let $m$ be the chain's head and $\textit{off} = H\cdot2^{\textit{depth}}$.
  \textit{up} shifts a bundle label along $m$ by \textit{off} (declaring it); \textit{base\_of}
  unwraps \code!Bun m b! to $b$ with shifted labels. Recurse on the rest of the chain in the
  unwrapped view, with the environment at the multiplicity value $j$ (a fresh name \code!j!, primes
  appended; non-bare multiplicities named \code!c!$k$).
\item From the inner spec ($L_b$, $R_b$, $iL_b$, $iR_b$, $k_1$), require every label to be shifted
  (\code!a boxed rewrite on unbundled legs!) and every node of $R_b$ copyable
  (\code!a boxed rule producing an uncopyable node!). Wrap both sides back in \code!Bun m! with the
  shift undone; emit \code!def bun!$c$\code!L : Net! and \code!def bun!$c$\code!R : Net! (the
  boxed regions, interface their unsummed labels).
\item Emit the theorem
\begin{lstlisting}[style=grammar]
/-- The boxed rewrite at every multiplicity. -/
theorem bunC SPBINDERS (others : ℕ) (j : ℕ) SPHYPS :
    Equal ENVj 1 bunCL ((K1) ^ j) bunCR := by
  induction j with
  | zero => exact Equal.of_zero_at FAM M OFF [.., 0, ..] (by simp) bunCL bunCR OKL OKR (by decide)
      (Equal.rescale (isoB ENV0 1 (Net.zero SPC bunCL) (Net.zero SPC bunCR) [] [] (by chk)) one_ne_zero (by simp))
  | succ j ih => exact Equal.of_split_at FAM M OFF j [.., j, ..] (by simp) bunCL bunCR OKL OKR BUNDPOS
      (Equal.rescale CHAIN (by norm_num) (by ring))
\end{lstlisting}
  where \code!CHAIN! is \lean!isoB! (split $L$ onto $L_b \mathbin{+\!\!+} L$) $\to$
  \code!plug ... 1 ((K1) ^ j) (Equal.toLocal _ _ _ ih)! (the $j$ remaining copies) $\to$ \lean!isoB!
  (reorder) $\to$ \code!plug ... 1 (K1) (LEMMA)! (the peeled copy) $\to$ \lean!isoB! (onto the split
  $R$), with interface the doubled list $x, x+\textit{off}$ of each interface label.
\item Return the spec with the unwrapped left ids, the wrapped right nodes, the lemma
  \code!Equal.toLocal _ _ _ (bun!$c$\code! ... (!\textit{current value}\code!) ...)! and the scalar
  \py!_Named("(K1) ^ (v)", k₁^v, "(pow_ne_zero _ ...)")!.
\end{enumerate}
\end{algorithm}

\subsection{W fusion into copies}\label{sec:lean-export-net-wfuse}

\py!w_fuse(t)! proves a kernel \code!w_fusion! inference: a W node whose single bundled output
(chain one multiplicity $m$) is summed and feeds exactly one node of kind \code!Bun m Z! or
\code!Bun m ZW! with that one input, one output and the same dimension. The goal must contain a
\code!WB [[m]]! node on the copies' output and a weighted Z spider before it. Scalars of the two
sides must agree (\code!the fused copies carry a scalar!). The lemma is
\code!wFuseZ ENV base (by simp) dz m bz ... (by pos) (Fam.env_dimv_of _ _ _ _ _ rfl rfl ⟨rfl, rfl⟩) (by chk)×3 (by decide)×4 HC (by simp only [zc]; W0)!
with \code!HC! the coefficient hypothesis of the base spider against the new weighted spider; the
step is \py!region_lines! with the W node and the copies as region. Failures:
\code!no W leg fans into copies of one Z spider!, \code!the copies' legs bundle elsewhere!,
\code!the W input is bundled!, \code!the goal has no W node on the copies' leg!,
\code!the goal has no weighted Z!.

\subsection{The family restatement on boxed diagrams}\label{sec:lean-export-net-boxed-literal}

\py!family_diagram_theorem(name, sl, a, sr, b, extra)! calls \py!family_diagrams! with the
current binders, the multiplicity values \code!(ksOf [...])!, the environment, the family text and
\code!Fpos ...!.

\begin{algorithm}[\py!family_diagrams!]\label{alg:lean-export-net-family-diagrams}
\begin{enumerate}
\item A side without boxes whose partner is boxed with an unboxed boundary
  (\py!_unboxed_boundary!: no boundary port on a boxed node or in a port scope) is treated as a
  boxed diagram too.
\item Emit \py!boxed_literal! of both; on \py!None! stop. If the mirrors' key signatures or
  boundary data (dimensions read through \py!raw_pairs!) differ, delete the literals just emitted
  and stop.
\item Emit
\begin{lstlisting}[style=leannet]
/-- The family theorem on the expansions of the engine's boxed diagrams, boundaries in key
order, through `DEqual_expandK_of_checks`. -/
theorem NAME_diagrams :
    ∀ BINDERS,
      DEqual ENV (sl) (Da.expandK KS) (sr) (Db.expandK KS) := by
  intro ARGS
  exact Diagram.DEqual_expandK_of_checks KS _ _
    (Diagram.FamOK_congr (F := FC) (F' := FAM) rfl rfl Da_ok)
    (Diagram.FamOK_congr (F := FC) (F' := FAM) rfl rfl Db_ok)
    (by decide +kernel) (by decide +kernel) (FPOS)
    (Da_net ▸ Db_net ▸ NAME ARGS)
\end{lstlisting}
  (a refutation negates it and applies \lean!Diagram.Equal_of_DEqual_expandK!).
\end{enumerate}
\end{algorithm}

\py!boxed_literal(a)! writes a boxed diagram as a \lean!Diagram! with \lean!boxes!, each box
\code!⟨id, mult index, none | some parent, [node ids], [ports]⟩!, node kinds unwrapped to their
base (\code!ZB!$\to$\code!Z!, \code!XB!$\to$\code!X!, \code!WB!$\to$\code!W!, \code!Bun! removed);
it records a \py!_BoxedMirror! and tries \py!boxed_plan!, the Python mirror of
\lean!BBridge.bridgeB!: the plain network (labels by slot and \textit{off}$+j$), the
\lean!bridgeD! plan, the bundled table (entry bits $bb = \max(1,\text{bit lengths})$,
$tb = \sum \textit{index}\cdot2^{bb\cdot\ell}$ over the $n_b+n_w$ labels), the box node masks
$\sum_x 2^{x}$ and the node-id mask. With a plan it emits \code!NP!$j$, \code!WB!$j$,
\code!LT!$j$, \code!LB!$j$, \code!MN!$j$, \code!MS!$j$ and
\code!theorem D!$j$\code!_bb : BBridge.bridgeB ... FC = true := by decide +kernel! with projections
\code!D!$j$\code!_net : dexpr.toNetB off = N!$j$ and \code!D!$j$\code!_ok : dexpr.FamOK FC off!;
without, both theorems \code!by decide +kernel!.

\py!_BoxedMirror! mirrors \lean!Diagram.walk! (box positions from a start parent, outermost first,
fuel the box count), \lean!nodeBoxes!, \lean!portBoxes! (walk from the node's innermost box),
\lean!nodeChain!, \lean!tail! (for a wired port, the part of the far end's node chain beyond the
near end's; for a boundary port, its port boxes' multiplicities), \lean!legChain!
(node chain then tail), \lean!portDim!, \lean!bndData! (dimension and leg chain per boundary port)
and \lean!keySig!. The key signature has one entry per boundary slot $i$ with port $q$: let $bs$ be the node boxes of $q$'s node, outermost first, $n_b = |bs|$ and $n_p$ the number of port boxes of $q$; for each box $b\in bs$ let $\mathit{mp}(b)$ be the least boundary slot on the same side as $i$ (both outputs or both inputs) whose port's node lies in $b$'s node scope, or 0 when none; let $\mathit{mps}$ be that list and $r$ the position of $i$ in it ($|\mathit{mps}|$ when absent). The entry is the pair $(\mathit{mps}[0{:}r],\ [r, r+1, \dots, n_b-1] \mathbin{+\!\!+} [n_b+n_p-1, \dots, n_b])$.

\section{Family exports beside the network proof}\label{sec:lean-export-net-beside}

Five modules write a family theorem for diagram pairs of fixed shapes without replaying a kernel
proof step by step: \module{archytaszx.semantics.lean_chainnode},
\module{archytaszx.semantics.lean_chainfam}, \module{archytaszx.semantics.lean_pairs} and
\module{archytaszx.semantics.lean_copyphase} prove a pair of boxed diagrams equal at every copy
count, and \module{archytaszx.semantics.lean_finite} joins the exports of finitely many instances.
All but the first build on the normalizer's export \py!export_symbolic! (\cref{ch:lean-export},
\cref{ch:lean-sym}) and read its emitted text back by regular expressions. Every one returns the
complete Lean file as a string.

\subsection{A chain box against a node box}\label{sec:lean-export-net-chainnode}

\begin{signature}
\py!export_chain_node(lhs: Diagram, rhs: Diagram, name: str = "main", var: Callable[[str], str] = lambda s: s) -> str | None!
\end{signature}
\noindent The pair must be exactly the following, else the result is \py!None!.
\begin{itemize}
\item \py!lhs! (\py!_is_start!): nodes $\{0, 1\}$, no wires, one box holding both nodes, no port
  scope, no parent, the single link (output 0 of node 1, input 0 of node 1), boundary inputs
  (node 1 in 0, node 0 in 0), boundary outputs (node 0 out 0, node 1 out 0), multiplicity a bare
  symbol.
\item \py!rhs! (\py!_is_goal!): nodes $\{0, 1\}$, no wires, one box holding node 1 only, no port
  scope, parent or links, boundary inputs (node 0 in 0, node 1 in 0), boundary outputs (node 1
  out 0, node 0 out 0), the same multiplicity symbol.
\item Every node of both a phaseless Z($1\to1$) of one common leg dimension (\py!_plain_id!), a
  symbol or a positive integer; both scalars 1.
\end{itemize}
The start is a chain box whose node 1 carries a wire from copy to copy; the goal is the node box
over node 1 beside one identity on the carried wire. The file holds the fixed literals
\code!D0! (the start with the carried wire's boundary ends behind the identities 2 and 3, kind
\lean!.I!), \code!C0 : DChain! and \code!D1!, the theorems \code!D0_eq : D0 = ChainNode.DL!,
\code!C0_eq : C0 = ChainNode.cL!, \code!D1_eq : D1 = ChainNode.DR! (each \code!by decide!), the
environment \code!E! (every leg of dimension $d$, every phase zero; a parameter
\code!(d : ℕ)! for a symbolic dimension) and
\begin{lstlisting}[style=leannet]
theorem main_diagrams :
    ∀ (d : ℕ) (n : ℕ),
      Equal (E d) 1 ((D0.expandChain C0 n).toNet (2 * n + 2)) 1 ((D1.expand (ksOf [n])).toNet (2 * n + 2)) := by
  intro d n
  rw [D0_eq, C0_eq, D1_eq]
  exact ChainNode.cyclic_every (E d) (fun _ => rfl) (fun _ => rfl) n _ le_rfl
\end{lstlisting}
with the names from \py!var! (\py!kernel._var! in \py!export!). The theorem is named
\textit{name}\code!_diagrams!; the header is the network header with the line
\code!a chain box carrying one wire against a node box holding it once (ArchytasZX/ChainNode.lean, cyclic_every), for every !$n$.

\subsection{Two closed chain boxes}\label{sec:lean-export-net-chainfam}

\begin{signature}
\py!export_chain_family(left: Diagram, right: Diagram, name: str = "main") -> str!
\end{signature}
\noindent proves two closed chain boxes equal at every copy count $k$ through \lean!chain_pair!
(\file{ArchytasZX/ChainPair.lean}, \file{ChainBridge.lean}; \cref{ch:lean-boxes}). Every
refusal raises \py!PathSumUnsupported!.

\paragraph{Sides.} \py!_chain_side(d)! requires one box with links that is a plain linked node
box (not cyclic, no port scope, index, parent, copy scalar, pair, tree, feed or alignment;
\code!the chain box is not a plain linked node box!; \code!not one chain box!), every leg of one
dimension (\code!the legs differ in dimension!) and only phaseless \code!Z!, \code!X!, \code!F!
nodes (\code!generator !$g$, \code!a phased node!). It lists the nodes (id, kind, outputs,
inputs) by id and the wires by sort key, puts a fresh phaseless Z($1\to1$) (ids from one above
the largest) behind every boundary port that is a link end, and records the box's nodes, the
links, the dimension, the multiplicity and the scalar.

\paragraph{Labels.} \py!_ChainLabels(s, W)! mirrors Lean's \lean!chLab0!, \lean!chLabT! at the
width $W$: with \textit{bnd} the outputs then inputs, the base wires (off every link) split into
inside wires $w_{\mathrm{in}}$ (both ends in the box; a crossing wire raises
\code!a wire crosses the chain box off its links!) and outside wires $w_{\mathrm{out}}$; each link
port's partner is the other end of its wire (\code!a link port without a wire!); the head nodes
are the partners of the links' first ports. A port's outside label is its boundary position, else
$|\textit{bnd}| + i$ on outside wire $i$, else $W + t$ when it is the partner of either port of
link $t$, else 0; its inside label is $W + |\text{links}| + i$ on inside wire $i$, $W + t$ for the
second port of link $t$, $2W + t$ for its first port, else 0. \py!nets()! gives \lean!chTl!
(nodes off the box and off the heads, summed labels the outside wires), \lean!chCp! (the box's
nodes at inside labels, summed the inside wires) and \lean!chHdN! (the heads, their boundary
labels below $|\textit{bnd}|$). \py!check()! mirrors \lean!ChainFamOK! and raises
\code!the boundary meets the chain box!, \code!a link does not run from a head node to a tail node!,
\code!a wire's ends read different labels! or \code!a leg off every wire, link and slot!.

\begin{algorithm}[\py!export_chain_family!]\label{alg:lean-export-net-chainfam}
\begin{enumerate}
\item Both sides by \py!_chain_side!; equal multiplicities and link counts
  (\code!the two chain boxes differ in count or links!), equal boundary sizes and dimensions
  (\code!the two sides differ in boundary!).
\item $W = 1 + \max$ of both sides' label counts (\py!sizes()!); check both sides; $m$ the
  boundary size, $u = [W, \dots, W + |\text{links}| - 1]$.
\item \py!_per_copy(s, k)!: solve the multiplicity $= k$ for its single symbol (else
  \code!the box multiplicity has no single symbol!; several solutions:
  \code!the box multiplicity is not linear!), read the scalar at $k$, and return its value at
  $k = 0$ and the ratio $s(k+1)/s(k)$, which must not depend on $k$
  (\code!the scalar does not grow by a fixed ratio per copy!).
\item \py!_fam_pieces! builds Lean's \lean!blkNet W u cp hd 1! (one copy with its shifted head),
  \lean!headNet u hd! and \lean!famNet W u tl cp hd (range m) 0! (the side at 0 copies) on the
  Python side. The three normalizer pieces are \code!PA! (left block at its ratio against the left
  head at 1), \code!PB! (the same on the right) and \code!P0! (both sides at 0 copies at their
  scalars at 0). Each piece's networks become engine diagrams (\py!_net_diagram!: summed labels
  as wires, boundary labels as boundary legs, outputs first; a label off two (one) legs raises
  \code!a summed label not on two legs! (\code!a boundary label not on one leg!)) and go to
  \py!export_symbolic!; a \py!NetUnsupported! becomes
  \code!the normaliser does not prove !\textit{tag}\code!: !\textit{reason}.
\item From each piece's source: the theorem \code!main! statement (else
  \code!!\textit{tag}\code!: a normaliser theorem this reader does not parse!); the dimension
  table must start \code!| 0 => d! (\code!!\textit{tag}\code!: the dimension table does not start at d!);
  the label maps from its networks \code!N0! and \code!N1! (or \code!N0!) onto the Python nets,
  node by node in order (\py!_piece_map!: equal kinds and arities, dimension and phase index 0, one
  to one), written as relabel lists \code!π!\textit{tag}\code!0!, \code!π!\textit{tag}\code!1!.
  All three must take the same hypotheses (\code!the normaliser pieces take different hypotheses!).
\item Emit the header comment \code!# A closed chain family!, the namespace with options
  \code!maxRecDepth 100000! and \code!maxHeartbeats 2000000!, each piece's source as namespace
  \textit{tag} (\py!_piece_source!: import, header and diagram restatement dropped, macros kept for
  the first piece only), the literals \code!DA!, \code!cA!, \code!DB!, \code!cB!,
  \code!def W : ℕ!, the environment \code!Ec d! (dimension $d$, phases zero, no labels),
  \code!hDc!, \code!okA!/\code!okB : ChainFamOK! (15 \code!decide +kernel! proofs),
  \code!fA!/\code!fB : famOK ... = true!, the moved pieces \code!sA!, \code!sB!, \code!s0!
  (each \code!(tag.main args).env0 (E' := Ec d)!, then \lean!Equal.mapL2! with the relabelings
  and \lean!ren_injective!, then \lean!reorderL!/\lean!reorder! onto the Lean pieces), and
\begin{lstlisting}[style=grammar]
theorem family HYPS (k : ℕ) :
    Equal (Ec d) (A0 * (RA) ^ k) ((DA.expandChain cA k).toNet M)
      (B0 * (RB) ^ k) ((DB.expandChain cB k).toNet M) :=
  chain_pair okA okB (by decide) (by decide) (Ec d) (hDc d) M (by decide) fA fB
    (sA ARGS) (sB ARGS) (s0 ARGS) k
theorem NAME HYPS (n : ℕ) : ...at the multiplicity... := family ARGS MULT
\end{lstlisting}
  with the multiplicity as a Lean natural in its symbol (\py!_mult_nat!).
\end{enumerate}
\end{algorithm}

\subsection{Pair boxes over their node boxes}\label{sec:lean-export-net-pairs}

\begin{signature}
\py!has_pairs(left: Diagram, right: Diagram) -> bool! --- either side has a pair box.\\
\py!export_pair_family(left: Diagram, right: Diagram, name: str = "main") -> str!
\end{signature}
\noindent proves two sides equal at every count of their node boxes through \lean!pair_family!
(\file{ArchytasZX/PairBox.lean}, \file{PairFam.lean}). Refusals raise \py!PathSumUnsupported!.
\py!_pair_side(d)! requires exactly two boxes, one a pair box over the other
(\code!not one node box with one pair box!, \code!the pair box is not over the node box!); a
plain node box (no links, groups, cyclic, port scope, parent, tree, feed, alignment) and a pair
box without port scope, links or tree
(\code!the boxes are not a plain node box and its pair box!); the pair box one node between one
left and one right port (\code!the pair box is not one node between two ports!), that node an
\code!F! (\code!the pair node is not a Fourier box!) whose input is wired to the left port and
output to the right port; no boundary inputs; one leg dimension; phaseless Z, X, F nodes; and the
node box holding every other node (\code!the node box does not hold every node!). It keeps the
node box's and the pair box's copy scalars (1 when absent).

\py!_Template! mirrors \lean!D.label m!, \lean!tmpl! (each node without its pair legs, wire $i$ at
$m+i$, the $m$ boundary slots) and \lean!PairOK! (\py!check!: the pair ports are the last legs of
their sides, off every wire and slot; every other leg on exactly one wire or slot).
\py!pin(q)! finds by breadth-first search over Z spiders sharing labels a chain from the node of
pair port $q$ to a boundary slot, returning the certificate $[(i, x), \dots]$ and the slot
(\code!a pair port off a Z spider!, \code!a pair port not tied to a slot by Z spiders!).

\begin{algorithm}[\py!export_pair_family!]\label{alg:lean-export-net-pairs}
\begin{enumerate}
\item Equal multiplicities and boundary sizes, equal dimensions; both templates checked; the
  pins of both pair ports must reach the same slots on both sides
  (\code!the pair legs are pinned to different slots!).
\item \py!_per_copy! as for chains, the ratio multiplied by the node box's copy scalar; both
  scalars at 0 copies must be 1 (\code!a scalar other than 1 at 0 copies!); equal pair copy
  scalars (\code!the pair copy scalars differ!).
\item The two templates at their per-copy ratios go to \py!export_symbolic! (as
  \code!PT!) and are moved onto \lean!tmpl! as above; the dimension must be one symbol $d$. The
  pair copy scalar $q = d^{r}$ ($r$ an integer or half integer, else
  \code!a pair copy scalar other than a power of the dimension!) is written
  \code!(((d : ℕ) : ℂ) ^ p)!, \code!((rt d) ^ p)!, or with \leanx!⁻¹! for $p<0$ (\py!_pair_text!).
\item Emit \code!PA!, \code!PB : PDiagram! with \code!pA!, \code!pB : DPair!
  \code!⟨1, 0, left, right, ⟨fid, .F, 0, 0, [0], [0]⟩⟩!, \code!bA!, \code!bB : DBox!, \code!Ec!,
  \code!okA!, \code!okB : PairOK!, \code!sT! (the templates on \code!Ec!), and
  \code!family : ... Equal (Ec d) (1 * (A) ^ k * Q ^ (k * (k - 1) / 2)) ((PA.expandP k).toNet (k * m)) ...!
  by \lean!Diagram.pair_family! with the four pin certificates checked by
  \code!pinOK ... = true! (\code!decide +kernel!), then \textit{name} at the multiplicity.
\end{enumerate}
\end{algorithm}

\subsection{Copy-index phases}\label{sec:lean-export-net-copyphase}

\begin{signature}
\py!gadget(d: Diagram) -> Gadget!\\
\py!export_copy_phase(left: Diagram, right: Diagram, describe=None) -> str!
\end{signature}
\noindent \py!gadget! requires one box, indexed, with no parent, links, cyclic flag, pair, copy
scalar, port scope, tree or alignment (\code!copy-index phases need one bang box!,
\code!copy-index phases need a plain indexed node box!); no leg dimension reading the index
(\code!a leg dimension reads the box index!); exactly one node $x$ whose phase reads it
(\code!copy-index phases need exactly one node reading the index!), a Z or X spider in the box,
whose phase is a rule without entries equal to $s\,j\,k/d$ ($j$ the box index, $k$ the phase
index variable, $s = 1$ for Z and $-1$ for X) over legs of one dimension symbol $d$
(\code!the index phase is not a rule!, \code!the index phase's legs are not of one dimension symbol!,
\code!the index phase is not s j k / d!). The gadget diagram replaces $x$ by a phaseless copy
with one more input: for Z it is wired from a new Fourier box in the box whose input becomes the
new last boundary input $K$; for X the new input is $K$ itself. The box loses its index.
\py!Gadget! holds (diagram, new node, kind, sign, Fourier box or \py!None!). All refusals raise
\py!NetUnsupported!.

\begin{algorithm}[\py!export_copy_phase!]\label{alg:lean-export-net-copyphase}
\begin{enumerate}
\item \begin{sloppypar}Both gadgets; \py!export_symbolic(ga, gb, "main", describe)! (with $K$ open). Its
  \code!main_diagrams! theorem (the \lean!DEqual_expandK_of_checks! form of
  \cref{alg:lean-export-net-family-diagrams}) and the binders of \leanx!Φ! are read back; else
  \texttt{the gadget pair's export has no family theorem on its diagrams} or
  \texttt{the gadget pair's family theorem does not read as written}. The dimension symbol must be a
  parameter (\code!the index phase's dimension is not a family parameter!).\end{sloppypar}
\item The boxed literals \code!D0!, \code!D1! are parsed and undone (\py!_undo!): the Fourier box
  and its wire removed, $x$ without its last input and with phase index $P = 1 + $ the largest
  phase index (left) or $P+1$ (right), the literal wrapped as an \lean!LDiagram! with
  \code!cph := [(x, box, 0, 0)]!. The phaseless indices of the two sides must agree
  (\code!the gadget pair's phaseless indices differ!); a misplaced $K$ raises
  \code!the gadget's Fourier box is not read as written! or
  \code!the gadget's boundary input is not read as written!.
\item Appended to the normalizer file before its \code!end!: the table
\begin{lstlisting}[style=leannet]
noncomputable def Ψ PARAMS (p : ℕ) : ℕ → ℝ :=
  if p < P then Φ ARGS p else fun k =>
    (if (Nat.unpair p).1 = P then (s0 : ℝ) else (s1 : ℝ)) *
      (((Nat.unpair p).2 : ℝ) * k / d)
\end{lstlisting}
  the theorem \code!dimx!, the linked diagrams \code!L0!, \code!L1!, the nodes \code!X0!,
  \code!X1!, \code!G0!, \code!G1 : LDiagram.Gad ...!, \code!W0!, \code!W1 : WF ...!,
  \code!CPh0/1! (rows below $P$ read as the family's), \code!CPx0/1! (row
  \lean!Nat.pair X.ph u! at $k$ is $s\,u\,k/d$), \code!CPz!, \code!CPb!, \code!CPs!, and
  \code!main_diagrams_original!: for every boundary assignment $\sigma$ in range on the copy
  indices, $s_0\cdot$\lean!L0.evalL! $= s_1\cdot$\lean!L1.evalL! in the environment with phase
  \leanx!Ψ!, by \lean!LDiagram.evalL_eq_of_gad!. Copy $t$ of $x$ reads row \lean!Nat.pair P t!;
  the file holds no copy count.
\end{enumerate}
\end{algorithm}

\subsection{Finite case splits}\label{sec:lean-export-net-finite}

\begin{signature}
\py!class FiniteUnsupported(ValueError)!\\
\py!nest(source, tag) -> tuple[str, list[str], str, str, str]!\\
\py!combine(instances, integral, text, var, shared) -> str!
\end{signature}
\noindent \py!kernel.finite_symbolic! (\cref{ch:kernel}) passes \py!combine! one export per value
combination of a restriction's finite clauses. \py!nest(source, tag)! drops the imports, renames
every top-level namespace $N$ other than \code!ArchytasZX! (and not nested in another) to
$N$\code!.!\textit{tag} in all its qualified references, turns \code!namespace ArchytasZX! /
\code!end ArchytasZX! into the tagged form, makes every \code!macro "! line \code!local!, and
returns the body, the imports, the namespace and name of the last top-level theorem, and its
proposition \code!∀ BINDERS, TYPE! (\py!_statement!: the text up to the first \code!:=! at bracket
depth 0, split at the first depth-0 colon). No theorem raises
\code!an instance export states no theorem!.

\begin{algorithm}[\py!combine!]\label{alg:lean-export-net-finite}
\begin{enumerate}
\item No instance raises \code!no instance!. Nest instance $k$ (from 1) with tag \code!I!$k$; a
  final theorem outside a namespace raises
  \code!an instance's final theorem is outside a namespace!. Imports are merged in order.
\item The case split applies when \py!shared! (every instance a family corollary), there is one
  finite symbol, it is integral (natural values of a multiplicity or dimension clause), the
  values are distinct, and the propositions agree once each instance's \code!(v : ℕ)! is written
  as the symbol's Lean name \textit{v} (\py!_uniform!). Then
\begin{lstlisting}[style=grammar]
theorem main_finite : ∀ v : ℕ, lo ≤ v → v ≤ hi → PROP := by
  intro v h_lo h_hi
  have h : v = a ∨ v = b ∨ ... := by omega
  rcases h with rfl | rfl | ...
  · exact NS1.NAME1
  ...
\end{lstlisting}
  for consecutive values (else the hypothesis is \code!v = a ∨ ... → ! and \code!intro v h!).
\item Otherwise each instance's namespace gains \code!def finite_claim : Prop! (its proposition)
  and \code!theorem finite_holds!, and \code!main_finite : NS1.finite_claim ∧ ...! is the
  anonymous constructor of the \code!finite_holds!.
\item \code!main_finite! is placed in the first instance's namespace with its \code!open! lines;
  the file starts with the comment
  \code!ArchytasZX proof export: restricted to !\textit{text}\code!, decided instance by instance; ...!
  naming \code!the case split over its values! or \code!the conjunction of the instance claims!.
\end{enumerate}
\end{algorithm}

\subsection{Callers in \texttt{certify\_decision}}\label{sec:lean-export-net-certify}

\module{archytaszx.semantics.certify_decision} (\cref{ch:proofs}) reaches these exporters as
follows.
\begin{itemize}
\item \py!decision_proof! sends a restricted decision with \py!instances! to
  \py!finite_proof!, which calls \py!restricted_proof! per instance (a \py!RecursionError!
  counts as a missing proof) and returns \py!kernel.FiniteProof(proofs, restriction, left, right)!,
  or \py!None! with \code!instances without a proof: !$a$\code! = !$v$\code! (!\textit{why}\code!); ...!.
\item \py!restricted_proof!, for a family decision without other evidence, tries in order:
  \py!export_chain_family! when the decision has a doubled pair (failure noted
  \code!chain family: ...; !), \py!export_pair_family! when \py!has_pairs! (\code!pair family: ...; !),
  \py!export_copy_phase! without a doubled pair (\code!copy-index phases: ...; !), each wrapped as
  \py!kernel.Symbolic!; then \py!pair_proof! with \py!concrete_kernel! true unless the pair is
  doubled and not a family.
\item \py!pair_proof(..., concrete_kernel=True)! tries the kernel's \py!from_decision! first for
  any non-contraction verdict and, with \py!concrete_kernel!, for a symbolic-contraction verdict on
  a concrete pair (\py!_concrete_pair!: no parameter and no free symbol in either diagram).
\item \py!path_sum_proof! tries, in order, \py!export_net!, \py!export_qudit_net!,
  \texttt{export\_\allowbreak{}qudit\_\allowbreak{}digits\_\allowbreak{}net} and \py!export_mixed_net! (\cref{ch:lean-export}), and joins the four
  reasons when all fail.
\end{itemize}

\section{Rule translators}\label{sec:lean-export-net-rules}

\py!_RULES! maps an engine rule name (\cref{ch:rules}) to a \term{translator}: a function from a
\py!_Ctx! to a \py!_Spec!. Every translator reads the matched nodes through the match fields of
the rule's pattern (\code!a_id!, \code!wire!, \code!state_id!, …, or the template node keys
\code!nodes!), checks the shape it needs with \py!_need(cond, message)! (raising
\py!NetUnsupported(message)! on failure), and writes the lemma of \file{ArchytasZX/Rules.lean} or
its companions with every label, dimension index and phase index as a literal and every side
condition as a proof text. \Cref{tab:lean-export-net-rules} lists every entry. The lemma names are
Lean declarations of the library (\cref{ch:lean-core}); $d$ is the value of the step's dimension
index, $\sqrt d$ the model \py!Rt!, and "copy" the copy scalar of \cref{def:lean-export-net-copy}.

\subsection{Wrappers}\label{sec:lean-export-net-wrappers}

\begin{itemize}
\item \py!_template(builder, **keys)! adapts a template rule: the match is replaced by a
  \py!_Shim! whose field $f$ is the template node \py!nodes[keys[f]]! (a Boolean key passes as
  is) and whose field \code!wire! is the match's first wire.
\item \py!_shape(builder)! converts a \py!ValueError!, \py!StopIteration!, \py!IndexError! or
  \py!KeyError! raised by an unpacking in the translator into
  \code!!\textit{rule}\code! step shape!.
\item \py!_s2(builder)! handles \code!zxw_s2! and \code!zxw_s2_pink!: without new nodes it is the
  identity removal on the template node; with a template node \code!z! (whose phase must be 0) it
  is \py!_z_perm!; otherwise the X spider \code!p! (one input, one output) becomes the Z spider of
  the result by \lean!identityXZ!.
\end{itemize}

\subsection{Proof-argument helpers}\label{sec:lean-export-net-helpers}

\begin{xltabular}{\linewidth}{@{}p{3.5cm}X@{}}
\toprule
Helper & Text \\
\midrule
\texttt{\_\allowbreak{}chks(n)} & $n$ copies of \texttt{(by chk)} separated by spaces. \\
\texttt{\_\allowbreak{}kdec(n)} & $n$ copies of \texttt{(by decide +kernel)}. \\
\texttt{\_\allowbreak{}phaseless(c, n)} & \texttt{(by phz)}; raises \texttt{a phase where the rule needs none} when the
  node's phase index is not 0. \\
\texttt{\_\allowbreak{}phase\_\allowbreak{}sum()} & \texttt{(by phs)}. \\
\texttt{\_\allowbreak{}positive(c, dz)} & \texttt{(by pos)}. \\
\texttt{\_\allowbreak{}dims\_\allowbreak{}mul()} & \texttt{(by simp only [E, $\Delta$] \textless{};\textgreater{} first \textbar{} rfl \textbar{} norm\_\allowbreak{}num \textbar{} ring)}. \\
\texttt{\_\allowbreak{}chks\_\allowbreak{}all(n)} & \texttt{(by intro w hw; simp only [List.mem\_\allowbreak{}cons, List.not\_\allowbreak{}mem\_\allowbreak{}nil, or\_\allowbreak{}false] at hw; rcases hw with rfl \textbar{} ... \textbar{} rfl \textless{};\textgreater{} chk)}
  ($n$ alternatives): every wire of a list has the step's dimension. \\
\texttt{\_\allowbreak{}each(proofs)} & \texttt{(List.forall\_\allowbreak{}mem\_\allowbreak{}cons.2 $\langle$p$_1$, ... (List.forall\_\allowbreak{}mem\_\allowbreak{}nil \_\allowbreak{})$\rangle$)}: a proof
  of \texttt{$\forall$ c $\in$ cs, P c} for a literal list. \\
\texttt{\_\allowbreak{}zc(n)} & \texttt{(}\textit{n.lean()}\texttt{ : Node)}. \\
\texttt{\_\allowbreak{}as\_\allowbreak{}r(lemma, made, want)} & The lemma unchanged when \texttt{made == want}; else, after checking
  \texttt{\_\allowbreak{}equiv\_\allowbreak{}b} pairwise (\texttt{rewritten nodes differ}),
  \texttt{LocalEq.equivR (lemma) rfl (by decide)}. \\
\texttt{\_\allowbreak{}as\_\allowbreak{}l(...)} & The same on the left: \texttt{LocalEq.equivL (lemma) rfl (by decide)}
  (\texttt{matched nodes differ}). \\
\bottomrule
\end{xltabular}

\subsection{Coefficient proofs}\label{sec:lean-export-net-coef}

Many lemmas take a hypothesis on spider coefficients, \lean!zc E n k! ($1\cdot e(\varphi_k)$ for
\code!Z!, \lean!coef ph k! for \code!ZW!) or \lean!xc!.

\py!_zc_terms(c, n, k)! returns the model of node $n$'s coefficient at $k$ and the rewrite that
turns its Lean phase into that model: for a weighted node, the model from \py!coef_models! (1 at
$k=0$) and no rewrite; for a plain node, its phase vector must be written entry by entry
(\code!phases written entry by entry!), the entry $t$ (turns) at $k$ must be rational
(\code!a rational phase!), and with $t = z + p/m$ ($z = \lfloor t\rfloor$, $m$ the denominator)
the model is $\omega_m^{p}$ and the rewrite
\leanx!rw [show ex (ENV.phase PH K) = om M ^ P from ex_of M P (Z) (by norm_num) (by norm_num [U, Φ])]!.

\py!_coef_case(c, lhs, rhs, rewrites)! proves one coefficient identity once \lean!zc! is
unfolded: the rewrites; then, when roots of unity of several orders occur (with $i$ counted as
order 4), every $\omega_m$ rewritten as $\omega_M^{M/m}$ for the least common multiple $M$ by
\code!show om m = om M ^ (M/m) from om_sub M m (M/m) (by norm_num) (by norm_num)! and
\code!show Complex.I = om M ^ (M/4) from I_sub M (M/4) (by norm_num) (by norm_num)!; finally
\code!exact (!\textit{sc.prove(lhs, rhs)}\code!)!. Each tactic is prefixed \code!all_goals !
and the tactics are joined by \code!; !. \py!_pair_rewrites! deduplicates the rewrites of
(node, index) pairs and appends \leanx!norm_num only [U, Ψ, ite_true, ite_false]! when a weighted
node occurs.

\py!_coef_hyp(c, left, right, fn="zc")! proves \leanx!∀ k < E.dimv dz, fn E left k = ∏ fn E right k!:
at a symbolic dimension, \leanx!(by intro k hk; simp only [fn, U, Φ]; SYM)! with \code!SYM! the
generic closer \py!_coef_sym!
(\leanx!(try simp only [Ψ]); (try split_ifs) <;> first | (exfalso; omega) | (simp [ex_zero]; done) | (norm_num [ex_zero]; done) | ring1 | (subst_vars; norm_num [ex_zero, rt_one]; done) | (subst_vars; push_cast; ring; done) | (exfalso; simp_all; done) | (simp_all; done)!,
the \leanx!Ψ! part only when weights exist); at a concrete dimension $d$,
\leanx!(by intro k hk; simp only [U, Δ] at hk; interval_cases k; · simp only [fn]; CASE₀ ... )!
with one bullet per $k<d$. \py!_coef_at(c, k, lhs, rhs, nodes)! is the single case
\code!simp only [zc]; !\textit{case}.

\begin{definition}[$K_j$ phases]\label{def:lean-export-net-kj}
\py!_kj_proof(c, ph, dz, j)! proves \lean!∀ q < E.dimv dz, ex (E.phase ph q) = om (E.dimv dz) ^ (q * j)!,
the phase $K_j$ with entries $qj/d$ turns. For phase 0 and $j=0$ it is
\leanx!(by intro q _; simp [E, Φ, ex_zero])!. At a symbolic dimension the phase must be a rule
equal to $j\cdot\textit{index}/d$ with no entries (\code!phase is not K_!$j$), proved by
\code!rw [om_eq_ex]!, \leanx!simp only [E, Φ, Δ]!, \code!split_ifs! and \code!ring_nf!. At a
concrete $d$ the phase must be written entry by entry and each entry at $q\ge1$ must differ from
$qj/d$ by an integer $z$ (\code!phase at !$q$\code! is not !$qj$\code!/!$d$\code! turns!); one
bullet \leanx!exact ex_of d (q*j) (z) (by norm_num) (by norm_num [E, Φ])! per $q$, the bullet for
$q = 0$ being \leanx!simp [E, Φ, ex_zero]!. \py!_kj_antipode(c, a, b, d, j)! proves at a symbolic
dimension that phase $b$ at $k$ equals phase $a$ at $(d-k)\bmod d$ for $a = K_j$, $b = K_{-j}$,
by \lean!ex_congr_int! with the integer $-j$.
\end{definition}

\subsection{The catalogue}\label{sec:lean-export-net-catalogue}

\begin{xltabular}{\linewidth}{@{}p{4.3cm}p{3.6cm}p{2.2cm}X@{}}
\caption{Rule translators (\texttt{\_\allowbreak{}RULES}). Column $k$ is the step scalar of the local
equation $L = k\cdot R$; ``$i_L$/$i_R$'' lists summed labels set explicitly.}\label{tab:lean-export-net-rules}\\
\toprule
Engine rules & Lean lemma & $k$ & Shape and notes \\
\midrule
\endfirsthead
\toprule
Engine rules & Lean lemma & $k$ & Shape and notes \\
\midrule
\endhead
\texttt{spider\_\allowbreak{}fusion}, \texttt{z\_\allowbreak{}fusion\_\allowbreak{}prime\_\allowbreak{}d} & \texttt{fusionZ}, \texttt{fusionX},
  \texttt{zcFusion}, \texttt{xcFusion}, \texttt{fusionZG}, \texttt{fusionXG} & 1 & One merged node, no
  deltas, one dimension. Two Z: \texttt{fusionZ} with the merged phase by \texttt{(by phs)}. Two X:
  \texttt{fusionX}, the node whose output carries the wire first. A weighted Z among them:
  \texttt{zcFusion} with a coefficient hypothesis; weighted X: \texttt{xcFusion}. Bundled legs:
  \texttt{fusionZG} / \texttt{fusionXG} through \texttt{\_\allowbreak{}as\_\allowbreak{}r}. \\
\texttt{identity\_\allowbreak{}removal} & \texttt{identityZ}, \texttt{identityX} & 1 & One delta, no new node; $R$ is
  the bare wire at the node's dimension. Phaseless node. \\
\texttt{zx\_\allowbreak{}cap}, \texttt{zx\_\allowbreak{}cap\_\allowbreak{}swapped} & \texttt{cap}, \texttt{capSwapped} & $\sqrt d$ & A Z state
  into an X effect (swapped: X into Z), both phaseless; $R$ empty. \\
\texttt{fourier\_\allowbreak{}state\_\allowbreak{}color\_\allowbreak{}change}(\texttt{\_\allowbreak{}swapped}) & \texttt{colorState}, \texttt{colorStateX},
  \texttt{colorEffect}, \texttt{colorEffectX} & 1 & A spider state into F (or F into an effect). \\
\texttt{state\_\allowbreak{}copy}(\texttt{\_\allowbreak{}swapped}) & \texttt{stateCopy}, \texttt{stateCopySwapped} & copy & An
  X state copied through a Z spider of $n$ outputs; the new states share one phase. \\
\texttt{triangle\_\allowbreak{}inverse\_\allowbreak{}cancellation} & \texttt{triangle}, \texttt{triangle\textquotesingle{}} & 1 & T then $T^{-1}$
  (or $T^{-1}$ then T); $R$ the bare wire. \\
\texttt{fourier\_\allowbreak{}cancellation} & \texttt{fourier4} & 1 & Four F in a row become one node; $i_L$ the
  three inner wires. \\
\texttt{bialgebra}(\texttt{\_\allowbreak{}swapped}), \texttt{zxw\_\allowbreak{}b2\_\allowbreak{}converse} (template) & \texttt{bialgebra},
  \texttt{bialgebraSwapped} & $\sqrt d$ & X with two inputs into Z with two outputs; $R$ two Z and
  two X; $i_R$ the four inner wires. \\
\texttt{hopf}(\texttt{\_\allowbreak{}swapped}) & \texttt{hopf}, \texttt{hopfW}, \texttt{hopfXW}, \texttt{hopfWXW} and
  \texttt{hopfSwapped...} & $1/d$ & A Z/X pair joined through an F path; \texttt{W} marks a weighted Z,
  \texttt{XW} a weighted X; hypothesis \texttt{k * d = 1}; $i_L = (u, v_1, v_2, v_3)$. \\
\texttt{w\_\allowbreak{}fusion}, \texttt{qf\_\allowbreak{}wf} (template) & \texttt{wFusion} & 1 & W into W. \\
\texttt{w\_\allowbreak{}identity} & \texttt{wIdentity} & 1 & A one-output W becomes a bare wire. \\
\texttt{w\_\allowbreak{}zero\_\allowbreak{}copy}, \texttt{triangle\_\allowbreak{}zero\_\allowbreak{}state}, \texttt{zxw\_\allowbreak{}bs0} (template) & \texttt{wZeroCopy},
  \texttt{triZeroState}, \texttt{triInvZeroState} & copy or 1 & An X state into W (copied, copy
  scalar) or into T, $T^{-1}$ (an X state on the output). \\
\texttt{triangle\_\allowbreak{}zero\_\allowbreak{}effect} & \texttt{triZeroEffect} & $\sqrt d$ & T into an effect. \\
\texttt{w\_\allowbreak{}z\_\allowbreak{}effect}, \texttt{w\_\allowbreak{}zero\_\allowbreak{}effect} & \texttt{wZEffect}, \texttt{wZeroEffect} & 1, $\sqrt d$ &
  W into a Z or X effect; $i_R=(v)$ for \texttt{wZEffect}. \\
\texttt{w\_\allowbreak{}z\_\allowbreak{}effect\_\allowbreak{}converse} & \texttt{LocalEq.symm\textquotesingle{} (wZEffect ...)} & 1 & $i_L=(v)$, $i_R=(w)$. \\
\texttt{connective\_\allowbreak{}states}(\texttt{\_\allowbreak{}swapped}) & \texttt{connStates}, \texttt{connEffects},
  \texttt{connStatesX}, \texttt{connEffectsX} & 1 & States into a binder or effects on a splitter;
  dimension product by \texttt{\_\allowbreak{}dims\_\allowbreak{}mul}. \\
\texttt{connective\_\allowbreak{}inverse}(\texttt{\_\allowbreak{}bind\_\allowbreak{}first}) & \texttt{connInverse}, \texttt{connInverseBind} & 1 &
  S then B (one wire) or B then S (two wires). \\
\texttt{wz\_\allowbreak{}bialgebra}, \texttt{zxw\_\allowbreak{}bzw\_\allowbreak{}converse} (template) & \texttt{wzBialgebra} & 1 & Z into W;
  rows, columns and pin pairs written as lists; $i_R$ the row legs. \\
\texttt{wz\_\allowbreak{}bialgebra\_\allowbreak{}converse} & \texttt{LocalEq.symm\textquotesingle{} (wzBialgebra ...)} & 1 & \\
\texttt{port\_\allowbreak{}box\_\allowbreak{}unfusion} & \texttt{unfuseZB}, \texttt{unfuseXB}, \texttt{unfuseXBi} & 1 & Splits one
  bundled leg off a bundled spider; wrapped by \texttt{\_\allowbreak{}as\_\allowbreak{}r} and \texttt{\_\allowbreak{}as\_\allowbreak{}l}; $i_R=(y)$. \\
\texttt{zxw\_\allowbreak{}sym}, \texttt{qf\_\allowbreak{}w\_\allowbreak{}sym} (and converses) & \texttt{wPerm} & 1 & W outputs permuted. \\
\texttt{qf\_\allowbreak{}z\_\allowbreak{}flex}(\texttt{\_\allowbreak{}converse}) & \texttt{zPerm}, \texttt{zcCongr} & 1 & Z legs permuted;
  weighted nodes go to \texttt{zcCongr}. \\
\texttt{phase\_\allowbreak{}to\_\allowbreak{}weight}, \texttt{weight\_\allowbreak{}to\_\allowbreak{}phase} & \texttt{zcCongr} & 1 & Same coefficients. \\
\texttt{scalar\_\allowbreak{}box} & \texttt{zcScalar} & step scalar & A legless Z box; the sum of its coefficients
  equals $k$. \\
\texttt{zxw\_\allowbreak{}ept} & \texttt{zcEpt} & $\sqrt d$ & X state into a spider effect. \\
\texttt{zxw\_\allowbreak{}zer} & \texttt{zcZer} & $1/\sqrt d$ & A one-legged spider of coefficients $[k=0]$ becomes
  an X state (\texttt{(by zcz)}). \\
\texttt{zxw\_\allowbreak{}ad} & \texttt{zcAd} & 1 & Two spider effects on a two-output W become one effect. \\
\texttt{zxw\_\allowbreak{}pcy} & \texttt{zcPcy} & 1 & A spider before a W copies onto both outputs. \\
\texttt{zxw\_\allowbreak{}p1} & \texttt{fourierSq\textquotesingle{}}, \texttt{fourierSq} & 1 & F;F becomes an X cup (or a cup and a
  bending Z). \\
\texttt{zxw\_\allowbreak{}p1\_\allowbreak{}multiplier} & \texttt{fourierMult} & $\sqrt d^{|w|-1}$ & F;F becomes a multiplier. \\
\texttt{zxw\_\allowbreak{}ta} & \texttt{zxwTA} & $1/\sqrt d$ & \\
\texttt{zxw\_\allowbreak{}bsj} & \texttt{zxwBsj} & $\sqrt d$ & Concrete $d$; $j$ from the bindings. \\
\texttt{zxw\_\allowbreak{}k1} & \texttt{kOne} & 1 & Requires step scalar 1 (\texttt{K1 scalar}). \\
\texttt{zxw\_\allowbreak{}k0} & \texttt{kCopy} & copy & $K_j$ state copied. \\
\texttt{zxw\_\allowbreak{}k2} & \texttt{zcK2} & step scalar & Concrete $d$. \\
\texttt{x\_\allowbreak{}spider\_\allowbreak{}definition} & \texttt{xDef}, \texttt{xDefW} & 1 & X spider as Z with F legs;
  $i_R$ the F wires. \\
\texttt{zxw\_\allowbreak{}h1} & \texttt{LocalEq.trans\textquotesingle{} (LocalEq.symm\textquotesingle{} (xDef ...)) (xDefT ...)} & 1 & With
  \texttt{kjNeg} at concrete $d$, \texttt{\_\allowbreak{}kj\_\allowbreak{}antipode} at symbolic $d$. \\
\texttt{qf\_\allowbreak{}mw} & \texttt{qfMW} & $\sqrt d^{|w|-1}$ & Concrete $d$. \\
\texttt{binder\_\allowbreak{}definition} & \texttt{bindDef0}, \texttt{bindDef}, \texttt{bindDefP} & 1 & A binder as a
  splitter, with none, all or some bent legs. \\
\texttt{qf\_\allowbreak{}e}, \texttt{qf\_\allowbreak{}k0}, \texttt{qf\_\allowbreak{}dd}, \texttt{qf\_\allowbreak{}bzw}, \texttt{qf\_\allowbreak{}dzx}, \texttt{qf\_\allowbreak{}ta},
  \texttt{qf\_\allowbreak{}s1} (all through \texttt{\_\allowbreak{}shape}) & \texttt{qfE}, \texttt{qfK0}, \texttt{qfDD}, \texttt{qfBZW},
  \texttt{qfDZX}, \texttt{qfTA}, \texttt{qfS1}, \texttt{qfS1c} & see \cref{sec:lean-export-net-mixed} &
  Mixed-dimension gadgets; the converse directions of \texttt{qf\_\allowbreak{}bzw}, \texttt{qf\_\allowbreak{}dzx}, \texttt{qf\_\allowbreak{}ta}
  raise \texttt{BZW forward}, \texttt{DZX forward}, \texttt{TA forward}. \\
\texttt{zxw\_\allowbreak{}d1} & \texttt{zxwD1} & 1 & $i_L=(w,v)$, $i_R=(v_2,w_2)$. \\
\texttt{qf\_\allowbreak{}d1}(\texttt{\_\allowbreak{}converse}) & \texttt{qfD1} (\texttt{LocalEq.symm\textquotesingle{}}) & 1 & F;F chains on a
  spider's legs. \\
\texttt{zxw\_\allowbreak{}vw} & \texttt{zxwVW} & $\sqrt d^{t-1}$ & V boxes; $t$ the multiplier wire count. \\
\texttt{zxw\_\allowbreak{}zv} & \texttt{zxwZV} & $(1/\sqrt d)^{t-1}$ & \\
\texttt{zxw\_\allowbreak{}va}(\texttt{\_\allowbreak{}converse}) & \texttt{zxwVA} (\texttt{LocalEq.symm\textquotesingle{}}) & 1 & \\
\texttt{qf\_\allowbreak{}va}(\texttt{\_\allowbreak{}converse}) & \texttt{qfVA} (\texttt{LocalEq.flip\textquotesingle{}}) & $(1/\sqrt d)^{t-1}$
  ($\sqrt d^{t-1}$) & \\
\texttt{zxw\_\allowbreak{}kz}(\texttt{\_\allowbreak{}converse}) & \texttt{zxwKZ} (\texttt{LocalEq.flip\textquotesingle{}}) & $(1/\sqrt d)^{t-2}$
  ($\sqrt d^{t-2}$) & $t\ge2$ (\texttt{KZ size}). \\
\texttt{zxw\_\allowbreak{}ww}(\texttt{\_\allowbreak{}converse}) & \texttt{zxwWW}, wrapped by \texttt{LocalEq.idR} / \texttt{idL}
  (\texttt{LocalEq.symm\textquotesingle{}}) & 1 & Identity Z spiders on the produced side inserted one by one. \\
\texttt{qf\_\allowbreak{}ww}(\texttt{\_\allowbreak{}converse}) & \texttt{qfWW} (as above) & 1 & \\
\texttt{qf\_\allowbreak{}hd}(\texttt{\_\allowbreak{}converse}) & \texttt{qfHD} (\texttt{LocalEq.flip\textquotesingle{}}) &
  $(-1)^{d+1}\cdot{}$\allowbreak$\sqrt d^{\,4-d-t}$ & Concrete $d$. \\
\texttt{zxw\_\allowbreak{}hd}(\texttt{\_\allowbreak{}converse}) & \texttt{zxwHD} (\texttt{LocalEq.flip\textquotesingle{}}) & $\sqrt d^{\,4-t}$ & \\
\texttt{triangle\_\allowbreak{}inverse\_\allowbreak{}definition\_\allowbreak{}converse} & \texttt{zcTi} & 1 & \\
\texttt{triangle\_\allowbreak{}definition\_\allowbreak{}converse} & \texttt{zcT} & 1 & \\
\texttt{zxw\_\allowbreak{}s1} & the fusion lemmas or \texttt{zcFusionN} & 1 & One wire: as spider fusion of the
  template nodes \texttt{a}, \texttt{b}; several wires: \texttt{zcFusionN}. \\
\texttt{zxw\_\allowbreak{}s2}, \texttt{zxw\_\allowbreak{}s2\_\allowbreak{}pink} & \texttt{identityZ}, \texttt{identityX}, \texttt{zPerm},
  \texttt{identityXZ} & 1 & See \texttt{\_\allowbreak{}s2}. \\
\texttt{zxw\_\allowbreak{}aso}(\texttt{\_\allowbreak{}converse}) & \texttt{wAso} & 1 & W associativity. \\
\bottomrule
\end{xltabular}

The rule library holds 151 rules and \py!_RULES! 91 of them. The 60 rules without a translator are \texttt{bialgebra\_\allowbreak{}swapped\_\allowbreak{}converse}, \texttt{binder\_\allowbreak{}definition\_\allowbreak{}converse}, \texttt{connective\_\allowbreak{}inverse\_\allowbreak{}converse}, \texttt{connective\_\allowbreak{}states\_\allowbreak{}effects\_\allowbreak{}converse}, \texttt{connective\_\allowbreak{}states\_\allowbreak{}states\_\allowbreak{}converse}, \texttt{connective\_\allowbreak{}states\_\allowbreak{}swapped\_\allowbreak{}effects\_\allowbreak{}converse}, \texttt{connective\_\allowbreak{}states\_\allowbreak{}swapped\_\allowbreak{}states\_\allowbreak{}converse}, \texttt{fourier\_\allowbreak{}cancellation\_\allowbreak{}converse}, \texttt{fourier\_\allowbreak{}effect\_\allowbreak{}color\_\allowbreak{}change\_\allowbreak{}converse}, \texttt{fourier\_\allowbreak{}effect\_\allowbreak{}color\_\allowbreak{}change\_\allowbreak{}swapped\_\allowbreak{}converse}, \texttt{fourier\_\allowbreak{}state\_\allowbreak{}color\_\allowbreak{}change\_\allowbreak{}converse}, \texttt{fourier\_\allowbreak{}state\_\allowbreak{}color\_\allowbreak{}change\_\allowbreak{}swapped\_\allowbreak{}converse}, \texttt{hopf\_\allowbreak{}converse}, \texttt{hopf\_\allowbreak{}swapped\_\allowbreak{}converse}, \texttt{qf\_\allowbreak{}bzw\_\allowbreak{}converse}, \texttt{qf\_\allowbreak{}dd\_\allowbreak{}converse}, \texttt{qf\_\allowbreak{}dzx\_\allowbreak{}converse}, \texttt{qf\_\allowbreak{}e\_\allowbreak{}converse}, \texttt{qf\_\allowbreak{}k0\_\allowbreak{}converse}, \texttt{qf\_\allowbreak{}mw\_\allowbreak{}converse}, \texttt{qf\_\allowbreak{}ta\_\allowbreak{}converse}, \texttt{qf\_\allowbreak{}wf\_\allowbreak{}converse}, \texttt{scalar\_\allowbreak{}box\_\allowbreak{}converse}, \texttt{state\_\allowbreak{}copy\_\allowbreak{}converse}, \texttt{state\_\allowbreak{}copy\_\allowbreak{}swapped\_\allowbreak{}converse}, \texttt{triangle\_\allowbreak{}definition}, \texttt{triangle\_\allowbreak{}inverse\_\allowbreak{}cancellation\_\allowbreak{}converse}, \texttt{triangle\_\allowbreak{}inverse\_\allowbreak{}cancellation\_\allowbreak{}ti\_\allowbreak{}t\_\allowbreak{}converse}, \texttt{triangle\_\allowbreak{}inverse\_\allowbreak{}definition}, \texttt{triangle\_\allowbreak{}zero\_\allowbreak{}effect\_\allowbreak{}converse}, \texttt{triangle\_\allowbreak{}zero\_\allowbreak{}state\_\allowbreak{}t\_\allowbreak{}converse}, \texttt{triangle\_\allowbreak{}zero\_\allowbreak{}state\_\allowbreak{}ti\_\allowbreak{}converse}, \texttt{w\_\allowbreak{}identity\_\allowbreak{}converse}, \texttt{w\_\allowbreak{}zero\_\allowbreak{}copy\_\allowbreak{}converse}, \texttt{w\_\allowbreak{}zero\_\allowbreak{}effect\_\allowbreak{}converse}, \texttt{x\_\allowbreak{}fusion\_\allowbreak{}converse}, \texttt{x\_\allowbreak{}spider\_\allowbreak{}definition\_\allowbreak{}converse}, \texttt{zx\_\allowbreak{}cap\_\allowbreak{}converse}, \texttt{zx\_\allowbreak{}cap\_\allowbreak{}swapped\_\allowbreak{}converse}, \texttt{zxw\_\allowbreak{}ad\_\allowbreak{}converse}, \texttt{zxw\_\allowbreak{}b2}, \texttt{zxw\_\allowbreak{}bs0\_\allowbreak{}converse}, \texttt{zxw\_\allowbreak{}bsj\_\allowbreak{}converse}, \texttt{zxw\_\allowbreak{}bzw}, \texttt{zxw\_\allowbreak{}d1\_\allowbreak{}converse}, \texttt{zxw\_\allowbreak{}ept\_\allowbreak{}converse}, \texttt{zxw\_\allowbreak{}h1\_\allowbreak{}converse}, \texttt{zxw\_\allowbreak{}k0\_\allowbreak{}converse}, \texttt{zxw\_\allowbreak{}k1\_\allowbreak{}converse}, \texttt{zxw\_\allowbreak{}k2\_\allowbreak{}converse}, \texttt{zxw\_\allowbreak{}p1\_\allowbreak{}converse}, \texttt{zxw\_\allowbreak{}p1\_\allowbreak{}multiplier\_\allowbreak{}converse}, \texttt{zxw\_\allowbreak{}pcy\_\allowbreak{}converse}, \texttt{zxw\_\allowbreak{}s1\_\allowbreak{}converse}, \texttt{zxw\_\allowbreak{}s2\_\allowbreak{}converse}, \texttt{zxw\_\allowbreak{}s2\_\allowbreak{}pink\_\allowbreak{}converse}, \texttt{zxw\_\allowbreak{}ta\_\allowbreak{}converse}, \texttt{zxw\_\allowbreak{}vw\_\allowbreak{}converse}, \texttt{zxw\_\allowbreak{}zer\_\allowbreak{}converse}, \texttt{zxw\_\allowbreak{}zv\_\allowbreak{}converse}. A step of one of them is proved through \py!_through_forward! when that succeeds and is an \code!unproved_! axiom otherwise.

\subsection{Local equation conventions}\label{sec:lean-export-net-conventions}

Every lemma states \lean!LocalEq ENV L iL R iR iface 1 k!. Its explicit arguments, in order, are:
the environment; the phase indices of the matched and new nodes; the dimension index; the labels
(wires, then the boundary labels of the region, then the lists of legs); the interface
\code!IFACE!; the scalar $k$ when the lemma takes it (\code!kt!); then the side conditions:
dimension positivity \code!(by pos)!, structural facts \code!(by chk)! (label dimensions,
disjointness, list memberships, permutations), \code!(by decide)! for literal facts, phase facts
\code!(by phz)!, \code!(by phs)!, coefficient hypotheses, $K_j$ proofs and the scalar identity.
The rewritten networks' node order is: context nodes in network order, then the region's nodes in
the order of \py!l_ids! (left) or of \py!r! (right). \code!LocalEq.symm'! turns a lemma around,
\code!LocalEq.flip' (lemma) (nonzero) _ (scalar identity)! turns it around and inverts its scalar,
\code!LocalEq.trans'! chains two lemmas, and \code!LocalEq.idL!, \code!LocalEq.idR! insert one
identity Z spider on the left or right side
(\code!LocalEq.idR (lemma) dz out in x a ph (Or.inl ⟨rfl, rfl⟩) (by chk)×7 (by phz) SIDE rfl!).

\subsection{Argument lists of the core lemmas}\label{sec:lean-export-net-arglists}

The following templates give the exact argument order the translators write for the lemmas of
the ZX, ZW and connective rules. \code!ENV! is the environment text; a name like \code!a.ph! is
the phase index of matched node \code!a!; \code![x.outs]! is a label list; \code!(n)! is a node
literal \code!(⟨...⟩ : Node)!; \code!CHK×k! stands for $k$ copies of \code!(by chk)!, \code!PHZ!
for \code!(by phz)!, \code!POS! for \code!(by pos)!, \code!PHS! for \code!(by phs)!;
\code!KT! and \code!HK! are the scalar text and identity proof of \cref{def:lean-export-net-copy}
(for \code!hopf...!: $k = 1/d$ and \code!(show k * d = 1 from (...))!); \code!HC! is a coefficient
hypothesis (\cref{sec:lean-export-net-coef}); \code!KJ(p)! is the $K_j$ proof of the phase of
node \code!p! (\cref{def:lean-export-net-kj}); \code!DIMS_MUL! is \py!_dims_mul()!.

\begin{lstlisting}[style=leannet]
fusionZ ENV a.ph b.ph r.ph dz w [a.outs] [a.ins] [b.outs] [b.ins] [r.outs] [r.ins] IFACE CHK×4 PHS
fusionX ENV f.ph s.ph r.ph dz w [f.outs] [f.ins] [s.outs] [s.ins] [r.outs] [r.ins] IFACE POS CHK×9 PHS
zcFusion ENV (a) (b) (r) dz w IFACE (by decide)×3 rfl rfl rfl (by decide)×3 CHK HC
xcFusion ENV (f) (s) (r) dz w IFACE (by decide)×3 rfl rfl rfl POS CHK×8 CHK HC(xc)
fusionZG ENV [c₁] [c₂] a.ph b.ph r.ph dz w [a.outs] [i₁] [t₁] [b.outs] [i₂] [t₂] [r.outs] [i₃] IFACE CHK×9 PHS
fusionXG ENV [f.co] [f.ci] [g.co] [g.ci] f.ph g.ph r.ph dz w [f.o] [f.ot] [f.i] [f.it] [g.o] [g.ot] [g.i] [g.it] [r.o] [r.i] IFACE CHK×4 POS CHK×10 PHS
zcFusionN ENV (a) (b) (r) dz [ws] IFACE (by decide)×3 rfl rfl rfl (by decide) w₀ (by decide)×4 CHKS_ALL HC
identityZ ENV n.ph dz p q w x IFACE CHK PHZ
identityX ENV n.ph dz p q w x IFACE POS CHK×4 PHZ
identityXZ ENV n.ph r.ph dz n.out n.in r.in r.out IFACE POS (by decide)×2 CHK×2 PHZ PHZ
cap ENV st.ph ef.ph dz w IFACE POS CHK PHZ PHZ                  -- capSwapped alike
colorState ENV s.ph r.ph dz w x IFACE POS CHK×2 PHZ PHZ         -- colorEffect alike
colorStateX ENV s.ph r.ph dz w x IFACE POS CHK×2 CHK×2 PHZ PHZ  -- colorEffectX alike
stateCopy ENV st.ph sp.ph cph dz w [sp.outs] IFACE KT POS CHK×3 PHZ PHZ PHZ HK
triangle ENV dz w o i x y IFACE CHK×6                           -- triangle' alike
fourier4 ENV a b r.ph dz w₁ w₂ w₃ IFACE POS CHK×14 PHZ
bialgebra ENV x.ph z.ph z₁.ph z₂.ph x₁.ph x₂.ph dz w a b c e n₀₀ n₀₁ n₁₀ n₁₁ IFACE POS CHK×5 PHZ×6
hopf ENV up.ph lo.ph dz u v₁ v₂ v₃ [up.outs] [up.ins] [lo.outs] [lo.ins] [nu.outs] [nu.ins] [nl.outs] [nl.ins] IFACE KT POS CHK×6 HK
wFusion ENV dz w i [first.outs] [second.outs] [r.outs] IFACE CHK×6 POS CHK
wIdentity ENV dz n.out n.in w x IFACE CHK×4
wZeroCopy ENV st.ph cph dz w [tg.outs] IFACE KT POS CHK×3 PHZ PHZ HK
triZeroState ENV st.ph dz w o IFACE POS CHK×3 PHZ               -- triInvZeroState alike
triZeroEffect ENV ef.ph r.ph dz w tr.in IFACE POS CHK×3 PHZ PHZ
wZeroEffect ENV ef.ph dz w i [wn.outs] [r.outs] IFACE POS CHK×6 PHZ
wZEffect ENV ef.ph dz w i v [wn.outs] [rest.outs] IFACE POS CHK×10 PHZ
connStates ENV ds dt dk a.ph b.ph r.ph a b k IFACE POS POS DIMS_MUL CHK×6 PHZ×3
connInverse ENV ds dt dk a b k₁ k₂ x y IFACE POS DIMS_MUL CHK×6
connInverseBind ENV ds dt dk k a₁ b₁ a₂ b₂ x₁ y₁ x₂ y₂ IFACE DIMS_MUL CHK×8
wzBialgebra ENV z.ph q dz w [zi] [wo] IFACE [[row]...] [[col]...] [(e, o)...] POS CHK×14 PHZ PHZ
wPerm ENV n.ph r.ph dz n.in [n.outs] [r.outs] IFACE (by decide)
zPerm ENV n.ph dz [n.outs] [n.ins] [r.outs] [r.ins] IFACE (by decide)
zcCongr ENV (n) (r) IFACE (by decide) (by decide) rfl (by decide) HC
zcScalar ENV (n) IFACE _ (by decide) rfl rfl HS
zcEpt ENV st.ph dz w (ef) IFACE (by decide) rfl (by decide) CHK POS PHZ HZ
zcZer ENV r.ph dz o (n) IFACE (by decide) rfl (by decide) (by decide) CHK POS PHZ (by zcz)
kOne ENV p.ph g.ph z.ph cph dz w i j [os] [us] IFACE POS CHK×10 KJ(p) KJ(c) PHZ PHZ
kCopy ENV st.ph sp.ph cph dz w j [sp.outs] IFACE KT POS CHK×3 KJ(st) PHZ KJ(c) HK
xDef ENV x.ph z.ph dz [(z, o)...] [(i, (a, b), z)...] IFACE POS (by decide)×2 CHK×2 HB
wAso ENV dz w w₂ i [first.outs] [second.outs] [p₁.outs] [p₂.outs] [o₃] IFACE CHK×6 POS CHK×7
\end{lstlisting}

The stems with a suffix share the template of their base: \code!hopfW!, \code!hopfXW!,
\code!hopfWXW! and the \code!hopfSwapped...! family, \code!bialgebraSwapped!,
\code!stateCopySwapped!, \code!connEffects!, \code!connStatesX!, \code!connEffectsX!,
\code!xDefW!. In \code!xDef!, \code!HB! is \code!(fun _ _ => rfl)! when both spiders have one
phase index, else \leanx!(by intro k _; simp only [U, Φ])! (\leanx!Ψ! for \code!xDefW!). In
\code!zcScalar!, \code!HS! sums the coefficients: at a concrete $d$,
\code!(by dsimp only; rw [show ENV.dimv dz = d from rfl]; simp only [Finset.sum_range_succ, Finset.sum_range_zero, zero_add, zc]; ...)!
with the step scalar against $\sum_{j<d} c_j$; at a symbolic dimension (phase index 0) the sum of
ones is rewritten by \code!Finset.sum_congr!, \code!Finset.sum_const!, \code!Finset.card_range!,
and the casts are closed by \code!push_cast! and \code!ring! (with \code!Fam.dimv! unfolded in a
family); at a symbolic dimension with a phase, \code!Finset.sum_congr rfl fun k _ => ?_! and
\code!by_cases hk : k = 0!. A weighted scalar box at a symbolic dimension raises
\code!a weighted scalar box at a symbolic dimension!.

\begin{example}\label{ex:lean-export-net-identity}
On the diagram of \cref{ex:lean-export-net-fusion}, one \code!identity_removal! step removes the
one-input, one-output Z spider (node 1). The region is that node; the step leaves one delta, the
bare wire from label 1 (the old boundary output) to label 3 (the wire). The step theorem is
\begin{lstlisting}[style=leannet]
theorem step1 :
    Equal E 1 N0 1 N1 := by
  refine Equal.trans (isoB E 1 _ ⟨([⟨.Z, [0, 3], [2], 0, 0⟩, ⟨.Z, [1], [3], 0, 0⟩] : List Node), [3], [0, 1, 2]⟩ [] [0, 1] (by chk)) ?_
  have hR := rewrite E ([⟨.Z, [0, 3], [2], 0, 0⟩] : List Node) ([⟨.Z, [1], [3], 0, 0⟩] : List Node) ([⟨.I, [1], [3], 0, 0⟩] : List Node) [3] [] [] [0, 1, 2] [1, 3] 1 1 1 (identityZ E 0 0 1 3 3 1 [1, 3] (by chk) (by phz)) (by chk) (by chk) (by chk) (one_mul _).symm
  refine Equal.trans hR ?_
  refine Equal.trans (isoB E 1 _ ⟨([⟨.Z, [0, 3], [2], 0, 0⟩, ⟨.I, [1], [3], 0, 0⟩] : List Node), [3], [0, 1, 2]⟩ [] [0, 1] (by chk)) ?_
  have hR := wire_elim E 1 ([⟨.Z, [0, 3], [2], 0, 0⟩] : List Node) ([] : List Node) [3] [0, 1, 2] 1 3 0 (by chk) (by chk) (by chk) (by chk) (by chk)
  refine Equal.trans hR ?_
  refine isoB E 1 _ N1 [] [0] ?_
  chk
\end{lstlisting}
The wire label 3 stays summed in the context ($iC = [3]$, $iL = iR = []$); \lean!wire_elim!
renames 3 to 1 and erases it, leaving the network \code!N1!.
\end{example}

\subsection{Mixed-dimension gadgets}\label{sec:lean-export-net-mixed}

The qufinite rules on mixed dimensions (\cite{wang2023qufinitezxw}) match gadgets of binders
\code!B!, splitters \code!S!, X states and effects (\term{pink dots}), multipliers (a Z spider
whose outputs feed one X spider) and Z boxes. The translators parse them into literals of Lean
records:

\begin{itemize}
\item \code!MixL!: \code!⟨a, k, v, d, ph, d', side⟩! for one leg of a mixed Z box: the pink
  label $a$, the box label $k$, the leg label $v$, the connective's first dimension index, the
  pink dot's phase, the connective's second dimension index and \code!false! for an input
  (binder) or \code!true! for an output (splitter). \py!_mix_legs(nodes, box)! finds for each box
  input the binder whose output it is and the X state on the binder's first input, and for each
  output the splitter and the X effect, raising \code!binder!, \code!splitter!,
  \code!pink state!, \code!pink effect! or \code!pink dimension!.
\item \code!qf_dzx! chains \code!⟨MixL_in, MixL_out, box, [ws], u, pz, px⟩! (the leg pair, the
  box node, the multiplier wires, the far label and the multiplier phases), each paired with its
  side flag.
\item \code!qf_ta! legs \code!TaC!: \code!⟨p, [ws], px, pz, MixL_a, MixL_b, MixL_o, box⟩!
  (\py!_ta_side!).
\item V boxes (\py!_parse_v!): the box, its W node, the pink dot, the multipliers
  \code!(w, [wires], v)! and the node ids in the order of the Lean \code!vNodes!; VA chains
  (\py!_parse_va!) \code!(w, box, (t, [wires], v))!.
\item \code!zxw_ww!, \code!qf_ww! green dots \code!⟨in₀, in₁, out, ph⟩! and columns
  \code!⟨..., ph, [col]⟩!; \code!qf_hd! and \code!zxw_hd! chains of ten fields.
\end{itemize}

Their scalars, with $\rho_x = \sqrt{x}$:
\code!qfE!: $\rho_{d_n}^{3-e}$ with $e$ the multiplier width;
\code!qfDD!: $\rho_{d}^{n-3}$ ($n\ge3$) or $(\rho_d^{-1})^{3-n}$;
\code!qfBZW!: $\rho_{r_0}\rho_{r_1}\rho_{r_2}$ over the three connective dimensions;
\code!qfDZX!: $\rho_{d}^{\,c(1-n)+2}\cdot m^{c-1}$ with $c$ chains, $n$ multiplier width and $m$
the small dimension, proved from
$k\cdot(\rho_m^{-1})^{c}m = (\rho_d\rho_m(\rho_d^{-1})^{n+1}d)^{c}(\rho_d^{-1})^{c}d$;
\code!qfTA!: $(\rho_n^{-1})^{c+1}n$;
\code!qfS1!: $\prod_{L_a}\rho\prod_{L_b}\rho\prod_{L_f}\rho^{-1}$ over the legs' connective
dimensions; \code!qfS1c! (every leg joined): $\prod_{L_a}\rho\prod_{L_b}\rho\sum_{j<\ell}
c_a(j)c_b(j)$ with $\ell$ the least of the legs' second dimensions and the boxes' dimensions;
\code!qfK0!: 0 when the box's coefficient $c$ at $j \bmod d_m$ is 0, else
$\rho_{d_m}\rho_{d_g}\rho_{d_{h_1}}\rho_{d_{h_2}}\rho_{d_{n_1}}^{-1}\rho_{d_{n_2}}^{-1}$
(the coefficient must be 0 or 1, \code!K0 coefficient 0 or 1!, and below the smallest dimension
unless 0). All of them require concrete dimensions (messages of the form
\code!... at concrete dimensions!).

\section{Size, memory and checkpoints}\label{sec:lean-export-net-size}

The network export writes one Lean theorem per proof step (\code!step!$c$), one per boxed
rewrite (\code!bun!$c$) and one per nested induction (\code!ind!$c$); the final theorem chains
them by \lean!Equal.mul_trans!. It emits no other checkpoints and splits no step across theorems;
the checkpointed and chunked path-sum exports are specified in \cref{ch:lean-pathsum}. The
controls on the size and cost of the Lean check are:

\begin{itemize}
\item the options \code!maxRecDepth 100000! and \code!maxHeartbeats 2000000!;
\item one network per diagram structure (\py!nets! cache) and one name per network, so a diagram
  that recurs is written once;
\item structural side conditions by \code!decide +kernel! first (the \code!chk! macro), so the
  kernel evaluates them without elaboration;
\item the packed bridges \lean!DBridge.bridgeD! and \lean!BBridge.bridgeB!
  (\cref{alg:lean-export-net-bridge}), which check a diagram literal against its network in time
  linear in the diagram, and their Python mirrors, which emit a bridge only when it succeeds and
  fall back on a plain \code!decide +kernel! otherwise;
\item the Python mirrors of \lean!nodeEquivB!, \lean!Net.split!, \lean!Net.zero!,
  \lean!Net.rename!, \lean!Kind.outCs! and the family checks, so that every emitted
  \lean!isoB! permutation and every split network is known to pass;
\item the label offsets $H\cdot2^t$ and $10^6$, which keep peeled and identified indices apart
  from ordinary ones.
\item phase vectors of more than \py!PHASE_TABLE_MIN! rational entries written as one packed
  numerator table over a common denominator, chunked into named definitions of
  \py!TABLE_CHUNK! entries past that size (\cref{def:lean-export-net-packed});
\item node, label and \code!labs! lists of more than \py!LIST_SPLIT! entries written as chunk
  definitions of \py!LIST_CHUNK! entries joined by \code!++!, and doc comments cut at
  \py!DOC_MAX! characters (\cref{sec:lean-export-net-labelled});
\item each region rewrite stated by \lean!have! before it is chained (\py!_stated!).
\end{itemize}

\section{Failures}\label{sec:lean-export-net-failures}

The export fails in two ways. A step that the exporter cannot prove becomes an \code!unproved_!
axiom and the export continues; a theorem outside the semantics raises \py!NetUnsupported!.
Inside a rule translator, iso or algebraic step, every \py!NetUnsupported! becomes an axiom; the
messages there are notes of the axiom. The following raise from \py!export! itself.

\begin{xltabular}{\linewidth}{@{}p{6.4cm}X@{}}
\toprule
Message & Condition \\
\midrule
\texttt{export takes a Theorem, got T} & The argument is not a kernel \texttt{Theorem}. \\
\texttt{a }$k$\texttt{ inference has no network proof} & An inference kind outside
  \cref{sec:lean-export-net-term} (or an induction outside a family). \\
\texttt{a chain, indexed, pair or copy-scalar bang box} & An extended \bbox{}. \\
\texttt{node boxes that do not nest}, \texttt{port boxes that do not nest},
  \texttt{a port box outside its node\textquotesingle{}s innermost box}, \texttt{a port box on a wire},
  \texttt{a wire between two node boxes} & \Cref{alg:lean-export-net-translate}. \\
\texttt{the boundary does not match}, \texttt{a boundary slot changes dimension},
  \texttt{the diagrams have different boundaries} & Boundary sizes or dimensions disagree. \\
\texttt{a wire joins dimensions }$d_a$\texttt{ and }$d_b$ & A wire between unequal dimensions not
  assumed equal. \\
\texttt{generator }$g$\texttt{ has no Lean kind} & A generator outside \texttt{\_\allowbreak{}KINDS} (the discard). \\
\texttt{port }$r$\texttt{ is neither wired nor on the boundary} & A dangling port. \\
\texttt{node }$n$\texttt{ has legs of different dimensions} & Mixed legs without \texttt{mixed\_\allowbreak{}ok}. \\
\texttt{a boxed }$k$\texttt{ node}, \texttt{a weighted spider with bundled legs},
  \texttt{a bundled leg on a }$k$\texttt{ node} & Bundling outside its forms. \\
\texttt{a rational ... where a natural number is required}, \texttt{the imaginary unit outside $\mathbb{C}$},
  \texttt{$\pi$ where a natural number is required}, \texttt{symbol ... of type ... used as ...},
  \texttt{... is not a natural number}, \texttt{the power ...}, \texttt{an exponential outside $\mathbb{C}$},
  \texttt{cannot write ... in Lean} & \Cref{def:lean-export-net-render}. \\
messages of \texttt{ScalarUnsupported} & A scalar without a Lean form. \\
\texttt{the label renaming is not injective} & \texttt{\_\allowbreak{}perm\_\allowbreak{}to\_\allowbreak{}bijection} outside a step. \\
\texttt{a diagram does not match its family form}, \texttt{two networks of the proof do not match},
  \texttt{the hypothesis region does not match}, \texttt{the networks do not match node for node} &
  Family and induction bridges. \\
\texttt{an induction over a fixed multiplicity},
  \texttt{a box whose multiplicity depends on the induction index} (\texttt{... shifted index},
  \texttt{... lifted multiplicity}), \texttt{an instance at a fixed multiplicity},
  \texttt{a lifted box inside a box at n \textgreater{}= c} & Inductions, shifts, lifts. \\
\texttt{a wire joins dimensions assumed equal},
  \texttt{no dimension of the start or goal is assumed equal to another},
  \texttt{the diagrams\textquotesingle{} scalars are not compared} & \Cref{sec:lean-export-net-identified}. \\
the messages of \cref{sec:lean-export-net-original} & The original-diagram restatement. \\
\texttt{Lean (lake) is not installed} & \texttt{check} without \texttt{lake}. \\
\texttt{a Zp node without two legs and a pair phase}, \texttt{a pair phase ... that is not c x y / M} &
  \Cref{sec:lean-export-net-zp}. \\
the messages of \cref{sec:lean-export-net-linked} & A linked induction (raised in the rerun). \\
\bottomrule
\end{xltabular}

A scalar without a model or an unprovable scalar identity of the final theorem does not fail:
the export records \py!scalar_gap!, states the theorem with the product of the step scalars, and
notes the gap in the header. A failed bridge plan or literal does not fail either: the export
falls back on \code!decide +kernel! or omits the diagram restatement.
